%%%%%%%%%%%%%%%%%%%%%%%%%%%%%%%%%%%%%%%%%%%%%%%%%%%%%%%%%%%%%%%%%%%%%%
% Lớp 11 | Hàm số mũ và hàm số lôgarít | Hàm số mũ. Hàm số lôgarít
% Dạng 1D6?3-5: Toán thực tế áp dụng hàm số mũ, lôgarít
% Số câu: 149  (sinh tự động bởi tach_cau.py; từ nay thầy sửa trực tiếp file này)
%%%%%%%%%%%%%%%%%%%%%%%%%%%%%%%%%%%%%%%%%%%%%%%%%%%%%%%%%%%%%%%%%%%%%%

%%%%%%%%%%%%%%%% TRẮC NGHIỆM 4 PHƯƠNG ÁN %%%%%%%%%%%%%%%%

% nguồn: Ok_TracNghiemLop11_Lan1
\begin{ex}%[1D6H3-5]%[ma:eadcecc9]
	Cường độ ánh sáng tại độ  sâu $h$ (m) dưới mặt hồ được tính theo công thức $I_h=I_0 \cdot\left(\dfrac{1}{2}\right)^{\tfrac{h}{4}}$ trong đó $I_0$ là cường độ ánh sáng tại mặt hồ. Cường độ ánh sáng tại độ sâu $5$ (m) gấp bao nhiêu lần cường độ ánh sáng tại độ sâu $13$ (m).
	\choice 
		{$4{,}3$} 
		{$5$} 
		 {\True $4$}
		 {$3$}
	\loigiai{Tỉ số cường độ ánh sáng tại độ sâu $5$ (m) so với  cường độ ánh sáng tại độ sâu $13$ (m) là
		\begin{eqnarray*}
			\dfrac{I_{5}}{I_{13}}=\left(\dfrac{1}{2}\right)^{\tfrac{5}{4}-\tfrac{13}{4}}=\left(\dfrac{1}{2}  \right)^{-\tfrac{8}{4}}=\left(\dfrac{1}{2}\right)^{-2}=4.
	\end{eqnarray*}}
\end{ex}

% nguồn: Ok_Toan11_TracNghiem_Lan2
\begin{ex}%[1D6V3-5]%[ma:d51a4393]
	Gia đình nhà bác Long Thắm gửi số tiền $ 100 $ triệu đồng vào ngân hàng với lãi suất $ 7\% $/năm. Biết rằng nếu không rút tiền ra khỏi ngân hàng thì cứ sau mỗi năm, số tiền lãi sẽ được nhập vào vốn ban đầu. Sau $ 10 $ năm, nếu không rút lãi lần nào thì số tiền mà nhà  bác Long Thắm nhận được gồm cả gốc lẫn lãi tính theo công thức nào dưới đây?
	 \choice
		 {$10^8.\left(1+0{,}07\right)^9$}
		{$10^8\left(1+0{,}7\right)^{10}$} 
		{\True $10^8\cdot\left(1+0{,}07\right)^{10}$} 
		{$10^8.0{,}07^{10}$} 
	 \loigiai{
	Áp dụng công thức lãi kép  thì số tiền mà nhà bác Long Thắm nhận được gồm cả gốc lẫn lãi là $$10^8\cdot\left(1+7\%\right)^{10}=10^8\cdot\left(1+0{,}07\right)^{10}.$$}
\end{ex}

% nguồn: Ok_Toan11_TracNghiem_Lan2
\begin{ex}%[1D6V3-5]%[ma:bc9d9a75]
	Ông A gửi $ 100$ triệu vào ngân hàng theo  hình thức lãi kép trong một thời gian khá lâu với lãi suất ổn định trong suốt thời gian tiết kiệm là $10\% $/$ 1 $ năm. Tết năm nay do dịch bệnh nên ông rút hết tiền trong ngân hàng ra để gia đình chi tiêu. Sau khi rút cả vốn lẫn lãi, ông trích ra $ 10 $ triệu để sắm sửa đồ Tết thì ông còn $ 240 $ triệu. Hỏi ông đã gửi tiết kiệm trong bao nhiêu năm?
	 \choice
		{$9$  năm}
		{$20$ năm} 
		{$12$  năm}
		{\True $10$ năm} 
	\loigiai{ 
		Giả sử ông A đã gửi  tiết kiệm trong $n$ năm.\\
		Số tiền ông  đã nhận được là $ 250 $ triệu.\\
		Theo công thức lãi suất kép, ta có 
		$$250\cdot 10^6=100\cdot 10^6\cdot \left(1+0{,}1\right)^n\Leftrightarrow  n=\log_{1{,}1}2{,}5\Leftrightarrow n\approx 9{,}61.$$
	Vậy ông A đã gửi  tiết kiệm trong $ 10 $ năm.}
\end{ex}

% nguồn: Ok_Toan11_TracNghiem_Lan2
\begin{ex}%[1D6V3-5]%[ma:ecfe2922]
	Vào ngày $ 15$ hàng tháng, ông An đều đến gửi tiết kiệm tại ngân hàng với số tiền $ 5$ triệu đồng theo hình thức lãi kép với lãi suất không đổi trong suốt quá trình gửi là $ 0{,}6\%  /$ tháng. Hỏi sau đúng ba năm, ông An thu được số tiền cả gốc lẫn lãi là bao nhiêu?
	\choice 
		{$  195\,251\,000$}
		{$  195\,252\,000$}
		{$  201\,450\,000$}
		{\True $ 201\,453\,000$} 
	 \loigiai{
		Đặt $ A=5$ triệu,  $ r=0{,}6\% $.\\
		Sau một tháng ông An có số tiền cả vốn và lãi là $ A\left(1+r\right)$, tiếp tục gửi vào ngân hàng $ A$ đồng nên số tiền trong ngân hàng lúc này  là $T_1=A\left(1+r\right)+A=A\left[\left(1+r\right)+1\right]$.\\
		Sau hai tháng ông An có số tiền cả vốn và lãi là  $T_2=T_1\left(1+r\right)+A=A\left[\left(1+r\right)^2+\left(1+r\right)+1\right]$.\\
		$\ldots$\\ 
	Sau $ 36$ tháng ông An có số tiền cả vốn và lãi là  $$T_{36}=A\left[\left(1+r\right)^{36}+\left(1+r\right)^{35}+\cdots+\left(1+r\right)+1\right]-A=A\cdot\dfrac{\left(1+r\right)^{37}-1}{r}-A\approx 201\,453\,000 \text{ đồng.}$$}
\end{ex}

% nguồn: Ok_Toan11_TracNghiem_Lan2
\begin{ex}%[1D6V3-5]%[ma:5aecc70b]
	Ông A có số tiền $120$ triệu đồng gửi tiết kiệm theo thể thức lãi suất kép, có hai loại để lựa chọn: loại kì hạn $12$ tháng với lãi suất $12{,}5\%$ trên một năm và loại kì hạn $1$ tháng với  lãi suất $1\%$ trên một tháng. Ông A muốn gửi $12$ năm. Theo anh chị ông A gửi loại nào sau $12$ năm sẽ nhận được tổng số tiền nhiều hơn và nhiều hơn bao nhiêu \textit{(làm tròn đến hàng nghìn)}?
	 \choice
		{Gửi theo kì hạn năm  lãi hơn kì hạn tháng $9\,879\,000$ đồng}
		{\True Gửi theo kì hạn tháng lãi hơn kì hạn năm  $9\,687\,000$ đồng}
		{Gửi theo hai  loại bằng nhau}
		{Gửi theo kì hạn năm lãi hơn  kì hạn tháng $9\,678\,000$ đồng}
	\loigiai{ 
		Loại kì  hạn $12$ tháng\\
		Số tiền có được sau $1$ năm là $120\cdot 10^6  + 120\cdot 10^6\cdot 0{,}125 = 120\cdot 10^6(1+0{,}125)$.\\
		Số tiền có  được sau $2$ năm là $120\cdot 10^6(1+0{,}125) + 120\cdot 10^6(1+0,125)\cdot 0,125 = 120\cdot 10^6(1+0{,}125)^2$.\\
		Tương tự vậy số tiền có được sau  $12$ năm là $120\cdot 10^6(1+0{,}125)^{12} \approx 493\,186\,000$ đồng.\\
		Loại kì hạn $1$ tháng ($12$ năm  là $144$ tháng)\\
		Số tiền có được sau $12$ năm là $120\cdot 10^6(1+0{,}01)^{144} \approx  502\,873\,000$ đồng.\\
		Vậy số tiền gửi theo kì hạn $1$ tháng nhiều hơn kì hạn năm là \\ 
		$502\,873\,000 -  493\,186\,000 = 9\,687\,000$ đồng.
	 }
\end{ex}

% nguồn: Ok_Toan11_TracNghiem_Lan2
\begin{ex}%[1D6V3-5]%[ma:bbba36c3]
	Bác Minh gửi $60$ triệu vào ngân hàng kì hạn $ 1 $  năm với lãi suất $5{,}6\% $ /năm. Biết rằng nếu không rút tiền ra khỏi ngân hàng thì cứ sau mỗi năm số tiền lãi sẽ nhập vào gốc để tính lãi cho năm tiếp theo. Hỏi sau ít nhất bao nhiêu năm bác Minh nhận được số tiền nhiều hơn $120$ triệu đồng bao gồm cả gốc và lãi?
	\choice 
		{$11$ năm} 
		{$12$  năm}
		{\True $13$  năm}
		{$14$ năm} 
	 \loigiai{
		Sau $n$ năm số tiền bác Minh  nhận được cả gốc và lãi là $60\left(1+5{,}6\%\right)^n$.\\
		 Vậy bác Minh nhận được số tiền nhiều hơn $120$ triệu đồng bao gồm cả gốc và lãi khi
		$$60\left(1+5{,}6\%\right)^n> 120 \Leftrightarrow n >\log_{1{,}056}2\approx  12{,}7$$
	Vậy bác Minh cần gửi ít  nhất $ 13 $ năm.}
\end{ex}

% nguồn: Ok_Toan11_TracNghiem_Lan2
\begin{ex}%[1D6V3-5]%[ma:5fdfad4c]
	Để đặc trưng cho độ to nhỏ của âm, người ta đưa ra khái niệm mức cường độ độ âm. Một đơn vị thường dùng để đo mức cường độ độ âm là dexinben (viết tắt là dB). Khi đó mức cường độ $L$ của âm được tính theo công thức $L(\text{dB}) = 10\log\dfrac{I}{I_0}$, trong đó $I$ là  cường độ của âm tại thời điểm đang xét, $I_0$ là cường độ âm ở ngưỡng nghe ($I_0 = 10^{-12}$ w/m$^2$). Một trò chuyện bình thường trong lớp học có mức cường độ âm trung bình là $68$ dB. Hãy tính cường độ độ âm tương ứng ra đơn vị w/m$^2$.
	 \choice
		{\True $6{,}3\cdot  10^{-6}$ w/m$^2$}
		{$6{,}3\cdot  10^{-72}$ w/m$^2$}
		 {$6{,}3\cdot 10^{6}$ w/m$^2$}
		{$6{,}3$  w/m$^2$}
	 \loigiai{
		Theo giả thiết ta có $L(\text{dB}) = 68$ dB, $I_0 = 10^{-12}$ w/m$^2$. Tính $I$.\\ 
		Áp dụng công thức ta có 
		\begin{eqnarray*}
			L(\text{dB}) &=& 10\log\dfrac{I}{I_0}\\ 
			68  &=& 10\log\dfrac{I}{I_0}\\
			\log\dfrac{I}{I_0} &=&  6{,}8\\
			 \dfrac{I}{I_0} &=& 10^{6{,}8}\\
			I  &=& 10^{6{,}8}\cdot I_0\\
			I &=&  10^{6{,}8}\cdot 10^{-12}\\
			 &\approx& 6{,}3\cdot 10^{-6} \, \text{w}/\text{m}^2.
		\end{eqnarray*}
		Vậy cường  độ âm là $I \approx 6{,}3\cdot 10^{-6}$ w/m$^2$.
	 }
\end{ex}

% nguồn: Ok_Toan11_TracNghiem_Lan2
\begin{ex}%[1D6V3-5]%[ma:08a96e50]
	Ông B vay ngân hàng $600$ triệu đồng và trả góp hàng tháng. Cuối mỗi tháng bắt đầu từ tháng thứ nhất anh trả $10$ triệu đồng và chịu lãi suất $0{,}8\%$ trên tháng cho số tiền chưa trả. Với hình thức hoàn  nợ như vậy thì sau bao lâu ông B sẽ trả hết số nợ ngân hàng?
	\choice 
		 {\True $50$}
		{$55$} 
		 {$60$}
		 {$65$}
	\loigiai{ 
		Số tiền còn lại sau $1$ tháng  là $600\cdot 10^6 + 600\cdot 10^6 \cdot 0{,}8\% - 10\cdot 10^6$.\\
		Khi đó ta có thể gọi số tiền vay là $A$, lãi  suất là $r$, số tiền trả mỗi cuối tháng là $m$ và $n$ là số tháng để trả hết tiền.\\
		Vậy số tiền còn nợ cuối tháng $1$ là $A(1+r) - m$.\\ 
		Số  tiền còn nợ cuối tháng $2$ là\\
		$(A(1+r) - m)(1+r) - m = A(1+r)^2 - m\left[(1+r) - m\right] = A(1+r)^2 - m[(1+r) +  1]$.\\
		Số tiền còn nợ cuối tháng  $n$ là $A(1+r)^n - \dfrac{m}{r}\left[(1+r)^n - 1\right]$.\\
		 Có nghĩa là khi trả hết tiền thì
		\begin{eqnarray*}
			A(1+r)^n  - \dfrac{m}{r}\left[(1+r)^n - 1\right] &=& 0\\
			600\cdot 10^6(1+0{,}8\%)^n -  \dfrac{10\cdot 10^6}{0{,}8\%}\left[(1+0{,}8\%)^n - 1\right] &=& 0\\
			 (1+0{,}8\%)^n &=& 1{,}48384\\
			 n &=& \log_{1+0{,}8\%} 1{,}48384 \approx 49{,}5
		\end{eqnarray*}
		Vậy khoảng  $50$ tháng thì ông B trả hết nợ.
	 }
\end{ex}

% nguồn: Ok_Toan11_TracNghiem_Lan2
\begin{ex}%[1D6V3-5]%[ma:133aba27]
	Một người gửi $ 100 $ triệu đồng vào một ngân hàng với lãi suất $ 6\% $/ năm. Biết rằng nếu không rút tiền ra khỏi ngân hàng thì cứ sau mỗi năm số tiền lãi  sẽ được nhập vào gốc để tính lãi cho năm tiếp theo. Hỏi ít nhất bao nhiêu năm người đó nhận được số tiền nhiều hơn $ 300 $ triệu bao gồm cả gốc lẫn lãi?.
	 \choice
		{$  20$ năm}
		{$ 18$  năm}
		{$  21$ năm}
		{\True  $ 19$ năm}
	\loigiai{ 
		Áp dụng công thức lãi kép thì sau $ n$ năm, số tiền  người gửi nhận được là $ A=10^8\cdot 1{,}06^n$.\\
		Để nhận được số tiền hơn $ 300 $ triệu  thì
		$$ A >3\cdot 10^8\Leftrightarrow{10^8}\cdot 1{,}06^n>3\cdot 10^8\Leftrightarrow 1{,}06^n> 3\Leftrightarrow n >\log_{1{,}06}3\approx  18{,}85.$$
	Vậy ít nhất sau $ 19 $  năm thì người đó nhận được số tiền nhiều hơn $ 300 $ triệu.}
\end{ex}

% nguồn: Ok_Toan11_TracNghiem_Lan2
\begin{ex}%[1D6V3-5]%[ma:88a8155b]
	Một người gửi số tiền $ 100$ triệu đồng vào ngân hàng với lãi suất $ 7 \%$/ năm. Biết rằng nếu không rút tiền ra khỏi ngân hàng thì cứ sau mỗi năm số tiền lãi sẽ được nhập vào vốn để  tính lãi cho năm tiếp theo. Hỏi sau ít nhất bao nhiêu năm người đó lĩnh được số tiền nhiều hơn $ 200$ triệu đồng, nếu trong khoảng thời gian này người đó không rút tiền ra và lãi suất không đổi?
	 \choice
		{\True  $ 11$ năm}
		{$ 12$ năm} 
		{$  10$ năm}
		{$ 9$ năm} 
	 \loigiai{
		Áp dụng công thức lãi  kép gửi $ 1 $ lần $ N=A{\left(1+r\right)^n}$, với $ A=100\cdot 10^6$ và $ r=7\%$.\\
	Theo đề bài ta tìm  $ n$ bé nhất sao cho $$10^8\left(1+7\%\right)^n>200\cdot10^6\Leftrightarrow{\left(1+7\%\right)^n}> 2\Leftrightarrow n >\log_{\tfrac{107}{100}}2\approx 10{,}245.$$}
\end{ex}

% nguồn: Ok_Toan11_TracNghiem_Lan2
\begin{ex}%[1D6V3-5]%[ma:70102266]
	Để đặc trưng cho độ to nhỏ của âm, người ta đưa ra khái niệm mức cường độ độ âm. Một đơn vị thường dùng để đo mức cường độ độ âm là dexinben (viết tắt là dB). Khi đó mức cường độ $L$ của âm được tính theo công thức là $L(\text{dB}) = 10\log\dfrac{I}{I_0}$, trong đó $I$ là cường độ của âm tại thời điểm đang xét, $I_0$ là cường độ âm ở ngưỡng nghe ($I_0 = 10^{-12}$ w/m$^2$). Tiếng ồn phát ra từ tiếng gõ phím liên tục ở một bàn phím của máy vi tính, có cường độ âm đo được là $10^{-5}$ w/m$^2$. Giả sử trong phòng làm việc của một công ty có hai nhân viên văn phòng cùng thực hiện thao tác gõ phím trên hai bàn  phím máy vi tính giống nhau thì mức cường độ âm tổng cộng do cả hai bàn phím phát ra cùng lúc là bao nhiêu?
	\choice 
		{\True $73$ dB} 
		{$140$ dB} 
		 {$100$ dB}
		 {$72$ dB}
	 \loigiai{
		Nếu  chỉ có một bàn phím có $L(\text{dB}) = 10\log\dfrac{I}{I_0} = 10\log\dfrac{10^{-5}}{10^{-12}} = 70$ dB.\\
		Cả hai bàn phím cùng gõ  \\
		$L_2 = 10\log\dfrac{2I}{I_0} = 10\log 2 + 10\log\dfrac{I}{I_0} = 10\log 2 +  70 \approx 3 + 70 = 73$ dB.\\
		Vậy có thêm một bàn phím  gõ thì mức cường độ âm tăng $3$ dB.
	} 
\end{ex}

% nguồn: Ok_Toan11_TracNghiem_Lan2
\begin{ex}%[1D6V3-5]%[ma:09f2729d]
	Cường độ một trận động đất $M$ (Richte) được cho bởi công thức $M = \log A - \log A_0$, với $A$ là biên  độ rung chấn tối đa và $A_0$ là một biên độ chuẩn (hằng số). Đầu thế kỷ $20$, một trận động đất ở San Francisco có cường độ $8$ độ Richte. Trong cùng năm đó, trận động đất khác Nam Mỹ có biên độ mạnh hơn gấp $4$ lần. Hỏi cường độ của trận động đất ở Nam Mỹ là bao nhiêu?
	\choice 
		 {$8{,}4$}
		{$32$} 
		{$12$} 
		{\True  $8{,}6$}
	 \loigiai{
		Trận động đất ở San Francisco có cường độ $8$ độ Richter khi đó áp dụng công thức\\ 
		$M_1 = \log A - \log A_0 \Rightarrow  8 = \log A - \log A_0$.\\
		Trận động đất ở Nam Mỹ  có biên độ là $4A$, khi đó cường độ của trận động đất ở Nam Mỹ là\\
		$M_2 = \log (4A) -  \log A_0 = \log 4 + \log A - \log A_0 = \log 4 + 8 \approx 0{,}6 + 8 = 8{,}6$ độ Richte.
	 }
\end{ex}

% nguồn: Ok_Toan11_TracNghiem_Lan2
\begin{ex}%[1D6V3-5]%[ma:e3a78d3d]
	Một người gửi $20$ triệu đồng vào ngân hàng theo hình thức lãi kép có kì hạn là $12$ tháng với lãi suất $6{,}3\% /$ năm. Giả sử qua các năm thì lãi suất không thay đổi và người đó không gửi thêm tiền vào mỗi năm. Để biết sau $y$ (năm) thì tổng số tiền cả vốn và lãi có được là $x$ (đồng), người đó sử dụng công thức $y = \log_{1{,}063} \left(\dfrac{x}{20}\right)$.  Hỏi sau bao nhiêu năm thì người đó có được tổng số tiền cả vốn và lãi là $30$ triệu đồng? (Làm tròn kết quả đến hàng đơn vị).
	\choice 
		{\True  $7$ năm}
		 {$6{,}6$ năm}
		 {$6$ năm}
		{$5$  năm}
	\loigiai{Người đó có được tổng số tiền cả vốn và lãi là $30$ triệu  đồng sau $y=\log_{1{,}063}\left( \dfrac{30}{20}\right) =6{,}6\approx 7$ năm.}
\end{ex}

% nguồn: Ok_Toan11_TracNghiem_Lan2
\begin{ex}%[1D6V3-5]%[ma:3cd14853]
	Đầu mỗi tháng, anh Hiếu gửi tiết kiệm ngân hàng số tiền $ 10 $ triệu đồng với hình thức lãi kép, lãi suất là $ 0{,}5 $\%/ tháng. Hỏi sau đúng $ 5 $ năm thì anh Hiếu nhận được số tiền cả gốc và lãi gần  nhất với số tiền nào dưới đây, giả sử rằng trong suốt quá trình gửi, anh Hiếu không rút tiền ra và lãi suất của ngân hàng không thay đổi.
	 \choice
		{$600+10\cdot 1{,}005^{60}$} 
		{\True $10\cdot 1{,}005\cdot  \dfrac{1{,}005^{60}-1}{0{,}005}$}
		{$10\cdot  1{,}005^{60}$}
		 {$10\cdot \dfrac{1{,}005^{60}-1}{0{,}005}$}
	\loigiai{ 
		Sau đúng  $ 5 $ năm thì anh Hiếu nhận được số tiền cả gốc và lãi là
		\begin{eqnarray*}
			 &T_{60}&=10^7\cdot \left(1+0{,}5\%\right)\left(\dfrac{\left(1+0{,}5\%\right)^{60}-1}{0,5\%}\right)\\
			& &=10\cdot  1{,}005\cdot \dfrac{1{,}005^{60}-1}{0{,}005}{10^6}\\
			& &=10\cdot 1{,}005\cdot  \dfrac{1{,}005^{60}-1}{0{,}005}.
		\end{eqnarray*}
	 }
\end{ex}

% nguồn: Ok_Toan11_TracNghiem_Lan2
\begin{ex}%[1D6V3-5]%[ma:e8baff5a]
	Ông Nguyễn Văn B là thương binh hạng $ 4/4$, được hưởng trợ cấp hàng tháng là $ 2\,082\,000$ đồng. Do tình hình dịch bệnh Covid-19 diễn biến phức tạp nên từ tháng $ 4 $ năm $ 2\,021 $ ông không  đi lĩnh tiền mà nhờ thủ quỹ lập một sổ tiết kiệm ở ngân hàng để gởi số tiền hàng tháng vào đó với lãi suất $ 0{,}5\% /$tháng. Hỏi đến đầu tháng $ 4 $ năm $ 2\,022 $ ông đến ngân hàng nhận được số tiền là bao nhiêu?
	\choice 
		{\True $  25\,811\,054$ đồng}
		{$ 2\,210\,413$ đồng} 
		{$ 25\,682\,641$ đồng} 
		{$  27\,893\,054$ đồng}
	\loigiai{ 
		 Ta có $$ T=A(1+r)\cdot \dfrac{\left(1+r\right)^n-1}{r}=2\,082\,000\left(1+0{,}5\%\right)\cdot \dfrac{\left(1+0{,}5\%\right)^{12}-1}{0{,}5\%}=25\,811\,054{,}06$$
	là số tiền ông B sẽ nhận  được.}
\end{ex}

% nguồn: Ok_Toan11_TracNghiem_Lan2
\begin{ex}%[1D6V3-5]%[ma:e3dc3e58]
	Áp suất không khí $P$ (đo bằng milimet thủy ngân, kí hiệu là mmHg) suy giảm mũ so với độ cao $x$ (đo bằng mét), tức $P$  giảm theo công thức $P = P_0\cdot\mathrm{e}^{ix}$ trong đó $P_0 = 760$ mmHg là áp suất ở mực nước biển ($x=0$), $i$ là hệ số suy giảm. Biết rằng ở độ cao $1\,000$m thì áp suất của không khí là $672{,}71$ mmHg. Hỏi áp suất không khí ở độ cao $3\,000$m gần với số nào sau đây nhất?
	 \choice
		 {\True $530{,}23$ mmHg}
		{$540{,}23$ mmHg} 
		{$520{,}23$ mmHg} 
		{$510{,}23$  mmHg}
	 \loigiai{
		 Áp dụng công thức $P = P_0\cdot\mathrm{e}^{ix}$.\\
		Ở độ cao $1\,000$m ta có $P_0 = 760$  mmHg, $x = 1\,000$ m, $P = 672{,}71$ mmHg, từ giả thiết này ta tìm được hệ số suy giảm.\\
		Ta có 
		\begin{eqnarray*}
			672{,}71 &=&  760\mathrm{e}^{1000i}\\
			\mathrm{e}^{1000i} &=& \dfrac{672{,}71}{760}\\ 
			 1000i &=& \ln\dfrac{672{,}71}{760}\\
			i &=& \dfrac{1}{1000}\ln\dfrac{672{,}71}{760}  \approx -0{,}00012.
		\end{eqnarray*}
		 Khi đó ở độ cao $3000$ m, áp suất của không khí là\\
		 $P = 760\mathrm{e}^{-0{,}00012 \cdot 3\,000} \approx 530{,}2340078$ mmHg.
	} 
\end{ex}

% nguồn: Ok_Toan11_TracNghiem_Lan2
\begin{ex}%[1D6V3-5]%[ma:b65fca4d]
	Năm $2020$, dân số thế giới là $7{,}795$ tỉ người và tốc độ tăng dân số $1{,}05\%/$năm. Nếu tốc độ tăng này tiếp tục duy trì ở những năm tiếp theo  thì đến năm bao nhiêu dân số đạt $10$ tỉ người.
	 \choice
		{$2040$} 
		{\True $2044$} 
		{$2048$} 
		 {$2052$}
	\loigiai{ 
		Dân số thế giới sau $1$ năm tính từ năm $2020$ là\\ $7{,}795\cdot (1+1{,}05\%)^{1} = 7{,}795\cdot 1{,}0105^1$  tỉ người.\\
		Dân số thế giới sau $2$ năm tính từ năm  $2020$ là\\ $7{,}795\cdot 1{,}0105\cdot (1+1{,}05\%)^{1} = 7{,}795\cdot 1{,}0105^2$ tỉ người.\\
		 Dân số thế giới sau $t$ năm tính từ năm $2020$ là $P(t) = 7{,}795\cdot 1{,}0105^t$ tỉ người.\\
		Giả sử sau ít nhất $t$ năm tính từ năm $2020$  thì dân số thế giới đạt $10$ tỉ người.\\
		Khi đó $7{,}795\cdot 1{,}0105^t = 10 \Leftrightarrow t \approx  23{,}85$.\\
		Do đó sau ít nhất $24$ năm thì  dân số thế giới đạt $10$ tỉ người.\\
		Vậy đến năm $2020  + 24 = 2044$ dân số đạt $10$ tỉ người.
	} 
\end{ex}

% nguồn: Ok_Toan11_TracNghiem_Lan2
\begin{ex}%[1D6V3-5]%[ma:a1b84182]
	Dân số thế giới được ước tính theo công thức $S = A\cdot \mathrm{e}^{ni}$, trong đó $A$ là dân số của  năm lấy mốc, $S$ là dân số sau $n$ năm, $i$ là tỷ lệ gia tăng dân số hàng năm. Biết năm $2023$ dân số thành phố Cần Thơ năm $2023$ ước tính là $1\,282\,000$ người và tỉ lệ gia tăng dân số là $1{,}03\%$. Hỏi đến năm bao nhiêu thì dân số thành phố Cần Thơ đạt hơn $1{,}5$ triệu người?
	\choice 
		{$2038$} 
		 {\True $2039$}
		 {$2040$}
		 {$2041$}
	\loigiai{ 
		Lấy năm $2023$ làm mốc, khi đó $A  = 1\,282\,000$.\\
		Giả sử sau $n$ năm thì dân số thành phố Cần Thơ đạt  hơn $1\,500\,000$, tức là\\
		\begin{eqnarray*}
			1\,500\,000  &=& 1\,282\,000\cdot \mathrm{e}^{0{,}0103n}\\
			 \mathrm{e}^{0{,}0103n} &=& \dfrac{1\,500\,000}{1\,282\,000} = \dfrac{750}{641}\\
			 0{,}0103n&=& \ln\dfrac{750}{641}\\
			\ln\dfrac{750}{641} &=&\dfrac{\ln\dfrac{750}{641}}{0{,}0103}  \approx 15{,}25
		\end{eqnarray*}
		Có nghĩa là theo tốc độ tăng dân số này thì sau $16$ năm dân số thành phố sẽ đạt $1\,500\,000$ người vào năm $2039$. 
	} 
\end{ex}

% nguồn: Ok_Toan11_TracNghiem_Lan2
\begin{ex}%[1D6V3-5]%[ma:f348793d]
	Cho biết chu kì bán hủy của chất phóng xạ plutonium Pu$^{239}$ là $24\,360$ năm (tức là lượng Pu $^{239}$ sau $24\,360$ năm  phân hủy thì chỉ còn lại một nửa). Sự phân hủy được tính bởi công thức $S = A\mathrm{e}^{rt}$, trong đó $A$ là lượng chất phóng xạ ban đầu, $r$ là tỉ lệ phân hủy hàng năm ($r < 0$), $t$ là thời gian phân hủy, $S$ là lượng còn lại sau thời gian phân hủy $t$. Hỏi $10$ gam Pu$^{239}$ sau bao nhiêu năm phân hủy sẽ còn $1$ gam?
	\choice 
		 {$48\,720$}
		 {$10\,000$}
		 {$82\,333$}
		{\True  $82\,235$}
	\loigiai{ 
		Trước  tiên, ta tìm tỉ lệ phân hủy hàng năm $r$ của Pu$^{239}$.\\
		Pu$^{239}$ có chu kì bán hủy là  $24\,360$ năm, do đó
		\begin{eqnarray*}
			5 &=&10\mathrm{e}^{r24\,360} \\ 
			r &=& \dfrac{\ln\dfrac{5}{10}}{24\,360}\\ 
			r &=& \dfrac{\ln5-\ln10}{24\,360}\\ 
			 r &\approx& -2{,}84544\cdot10^{-5}\\
			&\approx&  -0{,}000028.
		\end{eqnarray*}
		Vậy sự phân hủy của Pu$^{239}$ được tính bởi công  thức $S = A\mathrm{e}^{-0{,}000028t}$, trong đó $S$, $A$ tính bằng gam, $t$ tính bằng năm.\\
		Theo  đề bài cho ta có $1 = 10\mathrm{e}^{-0{,}000028t}\Leftrightarrow t = \dfrac{\ln 10}{0{,}000028} \approx 82\,235$ năm.\\
		Vậy sau khoảng $82\,235$ năm thì $10$  gam Pu$^{239}$ sẽ phân hủy còn lại $1$ gam.
	 }
\end{ex}

% nguồn: Ok_Toan11_TracNghiem_Lan2
\begin{ex}%[1D6V3-5]%[ma:4dbe28a1]
	Một sinh viên ra trường đi làm ngày $1/1/2023$ với mức lương khởi điểm là $a$ đồng mỗi tháng và cứ sau $2$ năm lại được tăng thêm $10\%$ và chi tiêu hàng tháng của anh ta là $60\%$ lương, phần còn lại tiết kiệm hết để mua nhà. Giá trị hiện tại của căn nhà là $1$ tỷ đồng và cứ sau $2$ năm thì giá trị tăng thêm $5\%$. Với mức lương khởi điểm $a$ là bao nhiêu thì sau $12$ năm anh ta mua được nhà (kết quả quy tròn đến hàng nghìn đồng). 
	 \choice
		 {$19\,028\,000$ đồng}
		{$16\,092\,000$ đồng} 
		 {$20\,092\,000$ đồng}
		 {\True $18\,092\,000$ đồng}
	 \loigiai{
		Đầu tiên ta tính giá trị của ngôi nhà sau 12 năm: \\ 
		Giá trị  ngôi nhà sau 2 năm: $10^9 + 10^9 \cdot 0{,}05 = 10^9 \cdot (1+0{,}05)$. \\
		Giá trị ngôi nhà sau 4 năm: $10^9 + 10^9 \cdot 0{,}05 + (10^9 + 10^9 \cdot 0{,}05) \cdot 0{,}05 = 10^9 \cdot  (1+0{,}05)^2$. \\
		Lần lượt ta có giá trị ngôi nhà sau 12 năm: $10^9 + 10^9 \cdot 0{,}05 + (10^9 + 10^9 \cdot 0{,}05) \cdot 0{,}05 \dots = 10^9 \cdot (1+0{,}05)^6$. \\ 
		Sau khi chi tiêu hàng tháng thì số tiền tiết  kiệm là $40\%$ lương. \\
		Có nghĩa là trong hai năm $2023-2024$, số tiền tiết kiệm là $24 \cdot 0{,}4a$. \\ 
		Trong hai năm tiếp theo $2025-2026$, số tiền tiết kiệm là $24 \cdot 0{,}4a(1+0,01)$.  \\
		Tương tự vậy số tiền tiết kiệm được trong  12 năm là \\
		$24 \cdot  0{,}4a\left[1 + (1+0{,}01) + (1+0{,}01)^2 + (1+0{,}01)^3 + (1+0{,}01)^4 + (1+0{,}01)^5\right] = 74{,}069856a$. \\
		Để mua được nhà thì số tiền trên phải bằng số tiền sau 12  năm: \\
		$74{,}069856a = 10^9 \cdot 1{,}05^6 \Rightarrow a \approx 18\,092\,321$.  \\
		Vậy số  $a$ gần bằng $18\,092\,000$.
	} 
\end{ex}

% nguồn: Ok_Toan11_TracNghiem_Lan2
\begin{ex}%[1D6V3-5]%[ma:c13ccfb9]
	Cường độ một trận động đất $M$ (Richte) được cho bởi công thức $M  = \log A - \log A_0$, với $A$ là biên độ rung chấn tối đa và $A_0$ là một biên độ chuẩn (hằng số). Đầu thế kỷ $20$, một trận động đất ở San Francisco có cường độ $8$ độ Richte. Trong cùng năm đó, trận động đất khác ở Nhật Bản có cường độ đo được $6$ độ Richte. Hỏi trận động đất ở San Francisco có biên độ gấp bao nhiêu lần biên độ trận động đất ở Nhật Bản.
	 \choice
		{$\dfrac{1}{100}$} 
		 {$\dfrac{4}{3}$}
		 {$2$}
		{\True $100$} 
	\loigiai{ 
		Trận động  đất ở San Francisco có cường độ $8$ độ Richter khi đó áp dụng công thức\\
		$M_1  = \log A_1 - \log A_0 \Rightarrow 8 = \log A_1 - \log A_0 $ (1).\\
		Trận  động đất ở Nhật Bản có cường độ $6$ độ Richter khi đó áp dụng công thức là\\
		$M_2 = \log A_2 - \log  A_0 \Rightarrow 6 = \log A_2 - \log A_0$ (2).\\
		Lấy (1) $-$ (2) vế với vế ta  được
		\begin{eqnarray*}
			2 &=& \log A_1 - \log A_2  \\
			 2 &=& \log \dfrac{A_1}{A_2}\\
			\dfrac{A_1}{A_2}  &=& 100\\
			 A_1 &=& 100\cdot A_2.
		\end{eqnarray*}
		Vậy trận động đất ở San Francisco  có biên độ gấp $100$ lần biên độ trận động đất ở Nhật Bản.
	 }
\end{ex}

% nguồn: Ok_Toan11_TracNghiem_Lan2
\begin{ex}%[1D6V3-5]%[ma:a85087bc]
	Để đặc trưng cho độ to nhỏ của âm, người ta đưa ra khái niệm mức cường độ độ âm. Một đơn vị thường dùng để đo mức cường độ độ âm là dexinben (viết tắt là dB). Khi đó mức cường độ $L$ của âm được tính theo công thức: $L(\text{dB}) = 10\log\dfrac{I}{I_0}$, trong đó $I$ là cường độ của âm tại thời điểm đang xét, $I_0$ là cường độ âm ở ngưỡng nghe ($I_0 = 10^{-12}$ w/m$^2$). Hai cây đàn ghita giống nhau, cùng hoà tấu một  bản nhạc. Mỗi chiếc đàn phát ra âm có mức cường độ âm trung bình là $60$ dB. Hỏi mức cường độ âm tổng cộng do hai chiếc đàn cùng phát ra là bao nhiêu?
	 \choice
		{$120$} 
		 {$80$}
		 {$180$}
		 {\True $63$}
	\loigiai{ 
		Mức cường độ âm do một chiếc đàn ghita phát ra $L_1  = 10\log\dfrac{I_1}{I_0} = 60$ dB.\\
		Mức cường độ âm  do hai chiếc đàn ghita cùng phát ra là\\
		$L_2 = 10\log\dfrac{I_2}{I_0} = 10\log\dfrac{2I_1}{I_0} = 10\log 2 + 10\log\dfrac{I_1}{I_0} = 10\log 2 +  60 \approx 63$ dB.\\
		Vậy mức cường độ âm  tổng cộng do hai chiếc đàn cùng phát ra là $63$ dB.
	 }
\end{ex}

% nguồn: Ok_Toan11_TracNghiem_Lan2
\begin{ex}%[1D6V3-5]%[ma:6762b353]
	Trong vật lí, sự phân rã của các chất phóng xạ được biểu diễn bằng công thức: $m(t) = m_0\left(\dfrac{1}{2}\right)^{\frac{t}{T}}$, trong đó $m_0$ là khối lượng chất phóng xạ ban đầu (tại thời điểm $t = 0$), $m(t)$ là khối lượng chất phóng xạ tại thời điểm $t$, $T$ là chu  kì bán rã (tức là khoảng thời gian để một nửa số nguyên tử của chất phóng xạ bị biến thành chất khác). Cho biết chu kì bán rã của một chất phóng xạ là $24$ giờ ($1$ ngày đêm). Hỏi $250$ gam chất đó sẽ còn lại bao nhiêu sau $3{,}5$ ngày đêm? \textit{(Kết quả làm tròn đến $3$ chữ số thập phân sau dấu phẩy)}
	 \choice
		 {\True $22{,}097$ gam}
		{$23{,}097$ gam} 
		 {$20{,}097$ gam}
		{$24{,}097$ gam} 
	 \loigiai{
		Áp dụng công  thức $m(t) = m_0\left(\dfrac{1}{2}\right)^{\tfrac{t}{T}}$.\\
		Với $m_0 = 250$, $T = 24$ giờ $= 1$ ngày đêm,  $t = 3{,}5$ ngày đêm.\\
		Ta có $m(3{,}5) = 250\left(\dfrac{1}{2}\right)^{\tfrac{3{,}5}{1}} \approx 22{,}097$ gam. 
	 }
\end{ex}

% nguồn: Ok_Toan11_TracNghiem_Lan2
\begin{ex}%[1D6V3-5]%[ma:d6cdbd6a]
	Một học sinh A khi đủ $ 18 $ tuổi được cha mẹ cho $200\,000\,000$ VNĐ. Số tiền này được bảo quản trong ngân hàng MSB với kì hạn thanh toán $ 1 $ năm và học sinh A chỉ nhận được số tiền này khi học xong $ 4 $ năm đại học. Biết rằng khi đủ $ 22 $ tuổi, số tiền mà học sinh A được nhận sẽ là $243\,101\,250$ VNĐ. Vậy lãi suất kì hạn một năm của ngân hàng MSB là bao  nhiêu?
	 \choice
		 {$ 8\% $}
		{$ 7\%  $}
		{$ 6\%  $}
		 {\True $ 5\% $}
	 \loigiai{
		Gọi lãi suất kỳ hạn một năm của  ngân hàng MSB là $ r $.\\
		Áp dụng công thức lãi suất kép $ P=a{\left(1+r\right)^n}$ trong đó ta có 
		\begin{eqnarray*}
			 & &243\,101\,250=200\,000\,000\left(1+r\right)^4\\
			&\Leftrightarrow &{\left(1+r\right)^4}=\dfrac{243\,101\,250}{200\,000\,000}\\ &\Leftrightarrow &1+r=\sqrt[4]{\dfrac{243\,101\,250}{200\,000\,000}}\Leftrightarrow  r=\sqrt[4]{\dfrac{243\,101\,250}{200\,000\,000}}-1\Leftrightarrow r=0{,}05.
		\end{eqnarray*}
	} 
\end{ex}

% nguồn: Ok_Toan11_TracNghiem_Lan2
\begin{ex}%[1D6V3-5]%[ma:6e2760ad]
	Một người gửi số tiền $3$ triệu đồng vào một ngân hàng với lãi suất $0{,}55\% /$ tháng. Biết rằng nếu người đó không rút tiền ra  khỏi ngân hàng thì cứ sau mỗi tháng, số tiền lãi sẽ được nhập vào vốn ban đầu. Số tiền người đó lãnh được sau hai năm, nếu trong khoảng thời gian này không rút tiền ra và lãi suất không đổi là
	 \choice
		 {$3\cdot\left(1{,}0055\right)^2$ triệu đồng}
		{\True  $3\cdot\left(1{,}0055\right)^{24}$ triệu đồng}
		{$3\cdot\left(1{,}055\right)^{24}$  triệu đồng}
		 {$\left(3{,}0055\right)^{24}$ triệu đồng}
	 \loigiai{
		Áp dụng công thức lãi kép $T_n=T_0\left(1+r\right)^n$ trong đó $T_0$ là số tiền gửi  ban đầu.\\
		$T_n$ là số tiền cả gốc lẫn lãi  sau $ n$ kì.\\
		$ r$ là lãi suất một kì.\\ 
		 Thời gian gửi là $ 2$ năm với chu kì $ 1$ tháng nên số chu kì là $ n=24$.\\
		 Suy ra số tiền người đó lãnh được sau hai năm là $$T_n=3\cdot\left(1+0{,}55\%\right)^{24}=3\cdot\left(1{,}0055\right)^{24}\text{ triệu đồng.}$$
	 }
\end{ex}

% nguồn: Ok_TracNghiemLop11_Lan1
\begin{ex}%[1D6V3-5]%[ma:ee0dd817]
	Sự tăng trưởng dân số được ước tính theo công thức sau $A=P\cdot\mathrm{e}^{r\cdot t}$, trong đó $P$ là dân số năm lấy làm mốc, $A$ là dân số sau $t$ năm, $r$ là tỉ lệ tăng dân  số hàng năm. Biết rằng vào năm $2020$ dân số Việt Nam khoảng $97{,}34$ triệu người và tỉ lệ tăng dân số là $0{,}91\%$. Nếu tỉ lệ tăng dân số này giữ nguyên, hãy ước tính dân số Việt Nam vào năm $2030$.
	\choice 
		{\True $106{,}61$  triệu người}
		{$105{,}61$ triệu người} 
		 {$241{,}82$ triệu người}
		{$100$ triệu người} 
	 \loigiai{
		Ta có  $A=97{,}34\cdot\mathrm{e}^{0{,}91\%\cdot 10}=106{,}61$ triệu người.
	} 
\end{ex}

% nguồn: Ok_TracNghiemLop11_Lan1
\begin{ex}%[1D6V3-5]%[ma:83491c03]
	Dân số thế giới được ước tính theo công thức $S=A\cdot \mathrm{e}^{\pi t}$, trong đó $A$ là dân số của năm lấy làm mốc, $S$ là dân số sau $n$ năm, $i$ là tỉ lệ tăng dân số hằng năm. Dân số Việt Nam năm $2019$ là $95{,}5$ triệu người, tỉ lệ tăng dân số hằng năm từ 2009 đến nay là $1{,}14\%$. Hỏi dân số Việt Nam năm $2009$ gần với  số nào nhất trong các số sau?
	 \choice
		 {$94{,}4$ triệu người}
		 {\True $85{,}2$ triệu người}
		{$86{,}2$ triệu  người}
		 {$83{,}9$ triệu người}
	\loigiai{ 
		Áp dụng công thức $S=A\cdot e^{\pi t}$ trong đó $S=95{,}5$ triệu người, $n=10$ năm, $i=1{,}14\%$.\\ 
	 Ta có số dân Việt Nam năm 2009 là \[A=\dfrac{S}{\mathrm{e}^{\pi t}}=\dfrac{95{,}5}{\mathrm{e}^{10\cdot 1{,}14\%}} \approx 85{,}2 \ \text{triệu người.}\]}
\end{ex}

% nguồn: Ok_TracNghiemLop11_Lan1
\begin{ex}%[1D6V3-5]%[ma:88cb635f]
	Một học sinh $A$ khi $15$ tuổi được hưởng tài sản thừa kế $200\,000\,000$ VNĐ. Số tiền này được bảo quản trong ngân hàng $B$ với kì hạn thanh toán $1$ năm và học sinh $A$ chỉ nhận được số  tiền này khi $18$ tuổi. Biết rằng khi $18$ tuổi, số tiền mà học sinh $A$ được nhận sẽ là $231\,525\,000$ VNĐ. Vậy lãi suất kì hạn một năm của ngân hàng $B$ là bao nhiêu?
	\choice 
		 {$8 \%$/năm}
		 { $7\%$/năm}
		{\True $6 \%$/năm} 
		{$5\%$/năm} 
	 \loigiai{
		Số tiền nhận được của  gốc và lãi là
		\begin{eqnarray*}
			 200\,000\,000(1+r)^3&=&231\,525\,000\\
			 (1+r)^3&=&1{,}157625\\
			 1+r&=&\sqrt[3]{1{,}157625}\\
			r&=&\sqrt[3]{1{,}157625}-1=0{,}05=5\%. 
		\end{eqnarray*}
	} 
\end{ex}

% nguồn: Ok_TracNghiemLop11_Lan1
\begin{ex}%[1D6V3-5]%[ma:14383592]
	Anh An gửi số tiền 58 triệu đồng vào một ngân hàng theo hình thức lãi kép và ổn định trong $9$ tháng thì lĩnh về được $61\,758\,000$ đồng. Hỏi lãi suất ngân hàng  hàng tháng là bao nhiêu? Biết rằng lãi suất không thay đổi trong thời gian gửi.
	\choice 
		 {$0{,}8 \%$}
		{  $0{,}6 \%$}
		 {\True $0{,}7 \%$}
		{$0{,}5\%$} 
	 \loigiai{
		Áp dụng công  thức $A_n=A_0\left(1+r\right)^n$ với $n$ là số kỳ hạn, $A_0$ là số tiền ban đầu, $A_n$ là số tiền sau $n$ kỳ hạn, $r$ là lãi suất.\\
		Suy ra  $A_9=A_0\left(1+r\right)^9\Rightarrow r=\sqrt[9]{\dfrac{A_9}{A_0}}-1$ = $0{,}7\%$.
	} 
\end{ex}

% nguồn: Ok_TracNghiemLop11_Lan1
\begin{ex}%[1D6V3-5]%[ma:4f47ed5f]
	Ông An gửi vào ngân hàng $60$ triệu đồng theo hình thức lãi kép. Lãi suất ngân hàng là $8\%$/năm. Sau $5$ năm ông An tiếp tục  gửi thêm $60$ triệu đồng nữa. Hỏi sau $10$ năm kể từ lần gửi đầu tiên ông An đến rút toàn bộ tiền gốc và tiền lãi được là bao nhiêu triệu đồng?
	\choice 
		 {$231{,}815$}
		 { $197{,}201$}
		 {\True $217{,}695$}
		{$190{,}271$} 
	\loigiai{ 
		Số tiền  ông An nhận được sau $5$ năm đầu là $60(1+8\%)^5=88{,}160$ (triệu đồng). \\
		 Số tiền ông An nhận được sau $10$ năm là
		\[(88{,}16+60)(1+8\%)^5=217{,}695  \ \text{(triệu đồng)}.\]
	} 
\end{ex}

% nguồn: Ok_TracNghiemLop11_Lan1
\begin{ex}%[1D6V3-5]%[ma:93b95417]
	Ông Tuấn gửi $100$ triệu đồng vào ngân hàng với hình thức lãi kép, kỳ hạn $1$ năm với lãi suất $8\%$. Sau $5$ năm ông rút toàn bộ tiền và dùng $1$ năm nữa để sửa nhà, số tiền còn lại ông tiếp tục gửi ngân hàng với lãi suất như lần trước. Số tiền lãi ông tuấn nhận được sau  $10$ năm gửi gần nhất với giá trị nào dưới đây?
	\choice 
		{$46{,}933$  triệu đồng}
		{$24{,}480$  triệu đồng}
		{\True  $81{,}413$ triệu đồng}
		 {$107{,}946$ triệu đồng}
	\loigiai{ 
		\begin{itemize}
			\item $5$ năm đầu ông Tuấn có số tiền cả gốc và lãi là \[T_1=100\cdot (1+0{,}08)^5=146{,}933 \ \text{(triệu  đồng).}\]
			\item Sau khi sửa  nhà số tiền còn lại gửi vào ngân hàng trong 5 năm thì số tiền cả gốc và lãi là
			 \[T_2=\dfrac{146{,}932}{2}(1+0{,}08)^5=107{,}946 \ \text{(triệu đồng).}\]
			 \item Số tiền lãi trong 10 năm là \[L=(149{,}933-100)+(107{,}946-73{,}466)=81{,}413\ \text{(triệu đồng)}.\]
		\end{itemize}
	} 
\end{ex}

% nguồn: Ok_TracNghiemLop11_Lan1
\begin{ex}%[1D6V3-5]%[ma:09aef5ef]
	Tại một xí nghiệp, công thức $P(t)=500\cdot \left(\dfrac{1}{2}\right)^{\frac{t}{3}}$ được dùng để tính giá trị còn lại (tính theo triệu đồng) của một chiếc máy sau thời gian $t$ (tính theo năm) kể từ khi đưa vào sử dụng. Sau một năm đưa vào sử dụng, giá trị còn lại của  máy bằng bao nhiêu phần trăm so với ban đầu?
	\choice 
		{$84,3\%$} 
		{$57,1\%$} 
		 {$39,3\%$}
		{\True  $79,4\%$}
	\loigiai{ 
		\begin{itemize}
			\item Giá trị  của máy khi chưa sử dụng là $P(0)=500\cdot \left(\dfrac{1}{2}\right)^0=500$ (triệu đồng).
			\item Giá trị còn lại của máy sau 1 năm là $P(1)=500\cdot \left(\dfrac{1}{2}\right)^{\frac{1}{3}}=397$ (triệu đồng). 
		\end{itemize}
		 Sau một năm đưa vào sử dụng, giá trị còn lại của máy bằng $\dfrac{397}{500}\cdot 100\%=79{,}4\%$ so với ban đầu.
	} 
\end{ex}

% nguồn: Ok_TracNghiemLop11_Lan1
\begin{ex}%[1D6V3-5]%[ma:056a6262]
	Anh Nam gửi $100$ triệu đồng vào ngân hàng theo thể thức lãi kép kì hạn là một quý với lãi suất $3\%$/quý. Sau đúng $6$ tháng anh Nam gửi thêm $100$ triệu đồng với kì hạn và lãi suất như  trước đó. Hỏi sau $1$ năm số tiền anh Nam nhận được là bao nhiêu?
	\choice 
		 {\True $218{,}64$ triệu đồng}
		{ $208{,}25$ triệu  đồng}
		{$210{,}45$ triệu  đồng}
		{$209{,}25$ triệu  đồng}
	 \loigiai{
		\begin{itemize}
			\item Số tiền anh Nam nhận  được sau 6 tháng là
			\[T_1=100(1+3\%)^2=106{,}09 \  \text{triệu đồng}.\]
			\item Số tiền anh Nam nhận  được sau một năm là
			 \[T_2=(106{,}09+100)(1+3\%)^2\approx 218{,}64 \ \text{triệu đồng}.\]
		\end{itemize}
	 }
\end{ex}

% nguồn: Ok_TracNghiemLop11_Lan1
\begin{ex}%[1D6V3-5]%[ma:e2756745]
	Giả sử cường độ ánh sáng $l$ dưới mặt biển giảm dần theo độ sâu theo công thức $l=l_0\cdot a^d$, trong đó $l_0$ là cường độ ánh sáng tại mặt nước biển, $a$ là một hằng số dương, $d$ là độ sâu tính từ mặt nước biển (tính bằng mét). Biết rằng ở một vùng biển $X$, cường độ ánh sáng tại độ sâu $1$ m bằng $95\%$ cường độ ánh sáng tại mặt nước biển thì tại độ sâu $15$ m ở vùng biển đó,  cường độ ánh sáng bằng bao nhiêu phần trăm so với cường độ ánh sáng tại mặt nước biển?
	 \choice
		 {$45{,}9\%$}
		{$38{,}1\%$} 
		{$59{,}7\%$} 
		{\True  $46{,}3\%$}
	\loigiai{ 
		Ta có cường độ ánh sáng tại độ sâu $1$ m bằng $95\%$ cường độ ánh sáng tại mặt nước  biển. \\ Suy ra với $d=1$ thì $l=\dfrac{95}{100}l_0$.\\
		Vậy $\dfrac{95}{100}l_0=l_0\cdot a^1\Rightarrow  a = \dfrac{19}{20}$. \\
		Tại độ sâu $15$ m, suy ra  $d=15$ và $a=\dfrac{19}{20}$.\\
		Ta có $l=l_0\cdot \left(\dfrac{19}{20}\right)^{15}\Rightarrow \dfrac{l}{l_0}=\left(\dfrac{19}{20}\right)^{15}\approx 0{,}4633 \approx  46{,}33\%$.
	 }
\end{ex}

% nguồn: Ok_TracNghiemLop11_Lan1
\begin{ex}%[1D6V3-5]%[ma:691ca246]
	Để dự báo dân số của một quốc gia, người ta sử dụng công thức $S=A\mathrm{e}^{\pi t}$ trong đó $A$ là dân số của năm lấy làm mốc tính,  $S$ là dân số sau $n$ năm, $r$ là tỉ lệ tăng dân số hàng năm. Năm $2017$, dân số Việt nam là $93\,671\,600$ người. Giả sử tỉ lệ tăng dân số hàng năm không đổi là $0{,}81\%$ dự báo dân số Việt nam năm $2035$ là bao nhiêu người?
	\choice 
		 {$109\,256\,100$}
		 {\True $108\,374\,700$}
		 {$107\,500\,500$}
		 {$108\,311\,100$}
	 \loigiai{
		Lấy năm 2017 làm mốc, ta có $A=93\,671\,600$, $n=2035-2017=18$.\\ 
		Suy ra dân số Việt Nam vào năm $2035$ là \[S=93\,671\,600\cdot \mathrm{e}^{18\cdot\frac{0{,}81}{100}}\approx 108\,374\,700 \  \text{người}.\]
	} 
\end{ex}

% nguồn: Ok_Toan11_TracNghiem_Lan2
\begin{ex}%[1D6C3-5]%[ma:63850dcd]
	Bạn Bình được gia đình gửi vào sổ tiết kiệm $ 200 $ triệu đồng với lãi suất $0{,}45\%$ một tháng theo hình thức lãi kép. Nếu mỗi tháng Bình rút ra một số tiền như nhau vào ngày ngân hàng trả lãi thì hàng tháng Bình rút ra số tiền gần nhất với số nào sau đây để đúng 4  năm vừa hết số tiền trong sổ tiết kiệm?
	 \choice
		 {$4\,620\,000$}
		{$4\,529\,000$} 
		{$4\,756\,000$} 
		{\True  $4\,642\,000$}
	\loigiai{ 
		Gọi số tiền bạn Bình rút ra hàng tháng là $x$ $\left(x >  0\right)$, số tiền ban đầu là $P$ $\left(P > 0\right)$, lãi suất tiền gửi hàng tháng là $r$ $\left(r > 0\right)$.\\
		Lãi suất nhận được sau tháng thứ nhất là $P\cdot r$.\\ 
		Số tiền cuối tháng thứ nhất sau  khi rút còn lại $P_1=P\left(1+r\right)-x$.\\
		Lãi  suất nhận được sau tháng thứ nhất là $P_1\cdot r$.\\
		Số tiền cuối  tháng thứ nhất sau khi rút còn lại $P_2=P_1\left(1+r\right)-x=P{\left(1+r\right)^2}-x\left(1+r\right)-x$.\\
		Cứ như thế, số tiền còn  lại sau $n$ tháng là:
		\begin{eqnarray*}
			&P_n&=P{\left(1+r\right)^n}-x{\left(1+r\right)^{n-1}}-x{\left(1+r\right)^{n-2}}-\cdots-x\left(1+r\right)-x\\ 
			 & &=P{\left(1+r\right)^n}-x\cdot\dfrac{\left(1+r\right)^n-1}{r}.
		\end{eqnarray*}
		Sau $ 48 $ tháng, số vừa  hết khi và chỉ khi
	$$P_n=0 \Leftrightarrow P{\left(1+r\right)^{48}}-x\cdot\dfrac{\left(1+r\right)^{48}-1}{r}=0 \Leftrightarrow 200\left(1{,}0045\right)^{48}-x\cdot\dfrac{\left(1{,}0045\right)^{48}-1}{0{,}0045}=0 \Leftrightarrow x\approx 4{,}642.$$} 
\end{ex}

% nguồn: Ok_Toan11_TracNghiem_Lan2
\begin{ex}%[1D6C3-5]%[ma:869fe45a]
	Vợ chồng nhà chị Thơm vay ngân hàng 400 triệu đồng để mua nhà với hình thức trả góp, chị chọn gói lãi suất ưu đãi cố định $ 0{,}5\% $/tháng trong $ 12 $ tháng đầu và sang tháng thứ $ 13 $ trở  đi thì ngân hàng tính lãi suất thả nổi theo quy định. Gia đình chị hoàn nợ cho ngân hàng theo cách sau đúng một tháng kể từ ngàỳ vay thì bắt đầu hoàn nợ; hai lần hoàn nợ liên tiếp cách nhau đúng một tháng, số tiền hoàn nợ ở mỗi tháng là $ 15 $ triệu đồng. Sau khi hết $ 12 $ tháng ưu đãi thì chị Thơm phải trả lãi suất thả nổi là $ 1\% $/tháng. Biết rằng mỗi tháng ngân hàng chỉ tính lãi trên số dư nợ thực tế của tháng đó và lãi suất thả nổi của ngân hàng không thay đổi trong thời gian chị Thơm hoàn nợ. Hỏi chị Thơm cần bao nhiêu tháng để trả hết nợ ngân hàng kể từ khi vay?
	\choice 
		 {$17$ tháng}
		{$29$ tháng} 
		 {\True $30$ tháng}
		{$18$ tháng} 
	 \loigiai{
		Số tiền  chị Thơm nợ sau một tháng là $ 400+400\cdot0{,}5\%=400(1+0{,}5\%)$.\\
		Sau 1 tháng  thì số tiền chị Thơm phải trả là $ 400(1+0{,}5\%)-15$.\\
		Sau 2 tháng thì số tiền chị Thơm phải  trả là\\
		$  400(1+0{,}5\%)-15+[400(1+0{,}5\%)-15]\cdot 0{,}5\%-15$\\
		 $=400(1+0{,}5\%)^2-15[(1+0{,}5\%)+1]$.\\
		Sau $ 12 $ tháng thì số tiền chị  Thơm phải trà là
		\begin{eqnarray*}
			&A&=400(1+0{,}5\%)^{12}-15\left[(1+0{,}5\%)^{11}+(1+0{,}5\%)^{10}+\cdots+(1+0{,}5\%)+1\right]\\ 
			&  &=400(1+0{,}5\%)^{12}-15\dfrac{(1+0{,}5\%)^{12}-1}{(1+0{,}5\%)-1}=400(1+0{,}5\%)^{12}-\dfrac{15}{0{,}5\%}\left[(1+0{,}5\%)^{12}-1\right]\\
			 & &=239{,}637.
		\end{eqnarray*}
		Gọi $  n$ là số tháng tiếp theo mà chị Thơm cần đề trả hết nợ, tương tự như trên ta được
		$$ A{(1+1\%)^n}-\dfrac{15\left[(1+1\%)^n-1\right]}{1\%}=0\Leftrightarrow n=\log_{1+1\%}\dfrac{-15}{A\cdot1\%-15}\approx 17{,}49.$$ 
		Tức là  chị Thơm cần thêm $ 18 $ tháng để trả hết nợ.\\
	Vậy chị  Thơm cần $ 12+18=30$ tháng để trả hết nợ ngân hàng kể từ khi vay.}
\end{ex}

% nguồn: Ok_TracNghiemLop11_Lan1
\begin{ex}%[1D6C3-5]%[ma:82916a80]
	Ông Chính gửi $200$ triệu đồng vào  một ngân hàng với lãi suất $7\%$ năm. Biết rằng nếu không rút tiền ra khỏi ngân hàng thì cứ sau mỗi năm số tiền lãi sẽ được nhập vào gốc để tính lãi cho năm tiếp theo và từ năm thứ $2$ trở đi, mỗi năm ông gửi thêm vào tài khoản với số tiền $20$ triệu đồng. Hỏi sau $18$ năm số tiền ông Chính nhận được cả gốc lẫn lãi là bao nhiêu? Giả định trong suốt thời gian gửi lãi suất không thay đổi và ông Chính không rút tiền ra.
	 \choice
		 {$1\,686\,898\,000$ VNĐ}
		{ $743\,585\,000$ VNĐ} 
		{$739\,163\,000$  VNĐ}
		{\True $1\,335\,967\,000$ VNĐ} 
	 \loigiai{
		Gọi $a=200$  triệu đồng, $b=20$ triệu đồng; $\alpha = 7\%$.
		\begin{itemize}
			\item Số tiền sau  1 năm là $a(1+\alpha)$.
			\item Số tiền sau 2 năm là $a(1+\alpha)^2+b(1+\alpha)$. 
			\item Số tiền sau 3 năm là $a(1+\alpha)^3+b(1+\alpha)^2+b(1+\alpha)$. 
			\item[$\ldots$] 
			\item Số tiền  sau 18 năm là
			\allowdisplaybreaks\begin{eqnarray*}
			A&=&a(1+\alpha)^{18}+b\left[(1+\alpha)^{17}+(1+\alpha)^{16}+\ldots (1+\alpha)\right]\\ 
			 &=&a(1+\alpha)^{18}+b\left[(1+\alpha)\cdot \dfrac{(1+\alpha)^{17} -1}{\alpha}\right]\\
			&=&200(1+7\%)^{18}+20\left[(1+7\%)\cdot \dfrac{(1+7\%)^{17}  -1}{7\%}\right]\\
			&\approx&1\,335,967 \ \text{(triệu đồng)} \approx  1\,335\,967\,000 \ \text{đồng}.
		\end{eqnarray*}
	\end{itemize}
	Vậy số tiền ông Chính nhận sau 18 năm là  $1\,335\,967\,000$ VNĐ.
 }
\end{ex}


%%%%%%%%%%%%%%%% ĐÚNG/SAI %%%%%%%%%%%%%%%%

% nguồn: Ok_dungsai_lop11_Lan2
\begin{ex}%[1D6H3-5]%[ma:a4288ef4]
	Lạm phát là sự tăng mức giá chung một cách liên tục của hàng hoá và dịch vụ theo thời gian, tức là sự mất giá trị của một loại tiền tệ nào đó. Chẳng hạn, nếu lạm phát là $5\%$ một năm thì sức mua của $1$ triệu đồng sau một năm chỉ còn là $950$ nghìn đồng (vì đã giảm mất $5\%$ của $1$ triệu đồng, tức là $50\,000$ đồng). Nói chung, nếu tỉ lệ lạm phát trung bình là $r\%$  một năm thì tổng số tiền $P$ ban đầu, sau $n$ năm số tiền đó chỉ còn giá trị là $A=P \cdot\left(1-\dfrac{r}{100}\right)^n$.
	 \choiceTF
		{\True Nếu  tỉ lệ lạm phát là $6\%$ một năm thì sức mua của $100$ triệu đồng sau một năm sẽ còn lại là $94$ triệu đồng}
		{\True Nếu tỉ lệ  lạm phát là $8\%$ một năm thì sức mua của $100$ triệu đồng sau hai năm sẽ còn lại là $84{,}64$ triệu đồng}
		{Nếu sức mua của $100$ triệu đồng sau hai năm chỉ còn là $90$ triệu đồng thì  tỉ lệ lạm phát trung bình của hai năm đó là $5{,}03\%$}
		{Nếu tỉ lệ lạm phát là $5\%$ một năm thì sau $14$  năm sức mua của số tiền ban đầu chỉ còn lại một nửa}
	\loigiai{ 
		\begin{itemchoice}
			\itemch Ta có $A=100 \cdot  \left(1-\dfrac{6}{100}\right)^1=100 \cdot 0{,}94=94$.\\
			Vậy tỉ lệ lạm phát là $6\%$ một năm thì sức  mua của $100$ triệu đồng sau một năm sẽ còn lại là $94$ triệu đồng.
			\itemch  Ta có $A=100 \cdot \left(1-\dfrac{8}{100}\right)^2=100 \cdot (0{,}92)^2=100 \cdot 0{,}8464=84{,}64$.\\
			Vậy tỉ lệ lạm phát là $8\%$ một năm thì sức mua của $100$ triệu đồng sau hai  năm sẽ còn lại là $84{,}64$ triệu đồng.
			\itemch Ta có $90=100 \cdot \left(1-\dfrac{r}{100}\right)^2 \Leftrightarrow \left(1-\dfrac{r}{100}\right)^2=0{,}9  \Leftrightarrow r\approx 5{,}13\%$.
			\itemch Ta có $A=100  \cdot \left(1-\dfrac{5}{100}\right)^{14}=100 \cdot (0{,}95)^{14}\approx 48{,}77$.
		\end{itemchoice}
	} 
\end{ex}

% nguồn: Ok_dungsai_lop11_Lan2
\begin{ex}%[1D6H3-5]%[ma:24c98561]
	Công thức $\log x=11{,}8+1{,}5M$ cho  biết mối liên hệ giữa năng lượng $x$ tạo ra (tính theo erg, $1$ erg tương đương $10^{-7}$ jun) với độ lớn $M$ theo thang Richter của một trận động đất.
	 \choiceTF
		{Trận động đất có độ lớn $2$ độ Richter tạo ra năng lượng khoảng $6{,}3\cdot  10^{34}$ erg}
		{\True Trận động đất có độ lớn $3$ độ Richter tạo  ra năng lượng khoảng $2\cdot 10^{9}$ (J)}
		{Trận động đất có độ lớn $5$ độ Richter tạo  ra năng lượng gấp $100$ lần so với trận động đất có độ lớn $3$ độ Richter}
		{\True Người ta ước lượng rằng một trận động đất có độ lớn khoảng từ $4$ đến $6$ độ Richter. Năng lượng do trận động đất đó tạo ra nằm trong khoảng $10^{17{,}8}\leqslant x\leqslant 10^{20{,}8}$ erg} 
	\loigiai{ 
	\begin{itemchoice}
		\itemch Tại $M=2 \Rightarrow \log x=11{,}8+1{,}5\cdot 2=14{,}8 \Rightarrow x= 10^{14{,}8}\approx  6{,}3\cdot 10^{14}$ erg.
		 \itemch Tại $M=3\Rightarrow \log x=11{,}8+1{,}5\cdot 3 =16{,}3 \Rightarrow x= 10^{16{,}3} erg \approx 10^{16{,}3}\cdot 10^{-7} \approx 2\cdot 10^9$.
		\itemch Gọi $x_1$, $x_2$ lần  lượt là năng lượng tạo ra của hai trận động đất có độ lớn lần lượt là \\
		 $M_1=5$ và $M_2=3$ (độ Richter). \\
		Ta có $\log  x_1=11{,}8+1{,}5\cdot M_1$, $\log x_2=11{,}8+1{,}5M_2$.\\
		Suy ra $\log x_1-\log x_2=1{,}5(M_1-M_2) \Rightarrow \log\dfrac{x_1}{x_2}=3  \Rightarrow \dfrac{x_1}{x_2}=10^3=1\,000$. \\
		Vậy trận động đất có độ lớn $5$ độ Richter tạo ra năng lượng gấp $1\,000$ lần so với  trận động đất có độ lớn $3$ độ Richter.
		\itemch Ta có $11{,}8+1{,}5\cdot  4 \leqslant \log x \leqslant 11{,}8+1{,}5\cdot 6 \Rightarrow 17{,}8\leqslant \log x\leqslant 20{,}8 \Rightarrow 10^{17{,}8}\leqslant x\leqslant 10^{20{,}8}$ erg.\\
		Vậy trận động đất đó tạo ra nằm trong khoảng $10^{17{,}8}\leqslant  x\leqslant 10^{20{,}8}$ erg.
	\end{itemchoice}
	 }
\end{ex}

% nguồn: Ok_dungsai_lop11_Lan2
\begin{ex}%[1D6V3-5]%[ma:d52a7a94]
	Tại sông Sài Gòn, cường  độ ánh sáng mặt trời đi qua môi trường nước được tính theo công thức
	$I = I_0 \cdot 10^{-\tfrac{x}{3}}$, trong đó $x$ là độ sâu (tính bằng mét) so với mặt nước sông, $I_0$  là cường độ ánh sáng tại mặt nước sông.
	\choiceTF 
		{\True Tại độ sâu $3$ mét cường độ ánh sáng  là $\dfrac{I_0}{10}$}
		{Tại độ sâu $6$ mét, cường độ ánh sáng  giảm đi $1\,000$ lần so với cường độ ánh sáng tại mặt nước sông}
		{\True Cường độ ánh sáng tại độ sâu $12$ mét bằng $\dfrac{1}{1\,000}$ lần cường  độ ánh sáng tại độ sâu $3$ mét}
		{Để cường độ ánh sáng giảm đi $\dfrac{1}{1\,000}$ so với cường độ ánh sáng tại mặt nước sông thì phải  ở độ sâu $5$ mét trở đi}
	\loigiai{ 
		\begin{itemchoice}
			 \itemch
			Tại độ sâu $3$ mét, ta có $x = 3 \Rightarrow I  = I_0 \cdot 10^{-3/3} = I_0 \cdot 10^{-1} = \dfrac{I_0}{10}$.
			\itemch 
			Tại độ sâu $6$ mét, ta có $x = 6 \Rightarrow I = I_0 \cdot 10^{-6/3} = I_0 \cdot  10^{-2} = 0{,}01 I_0$.
			 \itemch
			Tại độ sâu $12$ mét, ta có $x = 12 \Rightarrow I_{12} = I_0 \cdot 10^{-12/3} =  I_0 \cdot 10^{-4} = 0{,}0001 I_0$.\\
			Tại độ sâu $3$ mét, ta có $x = 3 \Rightarrow I_3 =  I_0 \cdot 10^{-1} = 0{,}1 I_0$.\\
			Suy ra $\dfrac{I_{12}}{I_3} = \dfrac{0{,}0001 I_0}{0{,}1 I_0} =  \dfrac{1}{1\,000}$.
			 \itemch
			Ta có $I = I_0 \cdot 10^{-x/3} = \dfrac{1}{1\,000} I_0  \Rightarrow 10^{-x/3} = \dfrac{1}{1\,000}$.\\
			Suy ra $-\dfrac{x}{3} = \log_{10} \dfrac{1}{1\,000} =  -3 \Rightarrow x = 9$.
		\end{itemchoice}
	} 
\end{ex}

% nguồn: Ok_dungsai_lop11_Lan2
\begin{ex}%[1D6V3-5]%[ma:061b8972]
	Một người gửi tiết kiệm vào ngân hàng theo hình thức lãi  kép với lãi suất $0{,}5\%$ mỗi tháng.
	 \choiceTF
		{\True Nếu số tiền ban đầu là $T$ thì số tiền sau $n$ tháng là $T\left(1{,}005\right)^n$} 
		{Nếu gửi $15$ triệu đồng thì sau $1$ năm người đó nhận được hơn $1{,}2$ triệu đồng  tiền lãi}
		{Nếu muốn số tiền tăng từ $50$ triệu lên $100$ triệu thì cần ít  nhất $140$ tháng}
		{\True Nếu gửi đều mỗi tháng $700$ nghìn đồng, sau $12$ tháng người đó nhận được hơn $9$ triệu đồng cả gốc lẫn lãi} 
	 \loigiai{
		\begin{itemchoice}
			\itemch Tiền sau $n$ tháng theo công thức lãi kép là 
			 \[
			A = T\left(1 + 0{,}5\%\right)^n = T\left(1{,}005\right)^n. 
			\] 
			 \itemch Tiền lãi sau $12$ tháng:
			 \[
			 A = 15\left(1{,}005\right)^{12} \approx 15 \cdot 1{,}0617 \approx 15{,}93 \Rightarrow Lãi \approx 0{,}93\ \text{triệu} < 1{,}2\ \text{triệu}.
			\] 
			\itemch Cần  tìm $n$ để:
			\[ 
			50\left(1{,}005\right)^n \geq 100  \Rightarrow \left(1{,}005\right)^n \geq 2
			\Rightarrow n  \geq \dfrac{\ln 2}{\ln\left(1{,}005\right)} \approx 139.
			\] 
			Cần ít nhất $139$ tháng,  không phải $140$ tháng.
			\itemch Gửi đều $T = 0{,}7$ triệu/tháng trong  $12$ tháng:
			\begin{eqnarray*}
				&&T\left(1{,}005\right)^{11} + T\left(1{,}005\right)^{10} + \cdots +  T\left(1{,}005\right)^0\\
				&=&T \cdot  \dfrac{\left(1{,}005\right)^{12} - 1}{0{,}005}\\
				&\approx&0{,}7  \cdot \dfrac{1{,}0617 - 1}{0{,}005}\\
				&\approx&0{,}7  \cdot 12{,}34 \approx 8{,}64.
			\end{eqnarray*}
			Tổng số tiến bằng $8{,}64 + 0{,}7  \cdot 12 = 8{,}64 + 8{,}4 = 17{,}04\ \text{triệu đồng}$ (hơn $9$ triệu đồng cả gốc lẫn lãi).
		\end{itemchoice}
	} 
\end{ex}

% nguồn: Ok_dungsai_lop11_Lan2
\begin{ex}%[1D6V3-5]%[ma:bc294dcb]
	Lạm phát là sự tăng mức giá chung một cách liên tục của hàng hoá và dịch vụ theo thời gian, tức là sự mất giá trị của một loại tiền tệ nào đó. Chẳng hạn, nếu lạm phát là $5\%$ một năm thì sức mua của $1$ triệu đồng sau một năm chỉ  còn là $950$ nghìn đồng (vì đã giảm mất $5\%$ của $1$ triệu đồng, tức là $50\,000$ đồng).
	Nói chung, nếu tỉ lệ lạm  phát trung bình là $r\%$ một năm thì tổng số tiền $P$ ban đầu sau $n$ năm số tiền đó chỉ còn giá trị là
	 \[
	A = P\left(1 - \dfrac{r}{100}\right)^n. 
	 \]
	 \choiceTF
		{\True Nếu tỉ lệ lạm phát là $7\%$ một năm thì sức mua của $100$ triệu đồng sau hai năm sẽ  còn lại $86\,490\,000$ đồng}
		{Nếu tỉ lệ lạm  phát là $7\%$ một năm thì sức mua của $100$ triệu đồng sau hai năm sẽ còn lại $96\,490\,000$ đồng}
		{Nếu sức mua của $100$ triệu đồng sau ba năm chỉ còn lại $80$ triệu đồng thì tỉ lệ lạm phát trung bình mỗi năm  là $9{,}17\%$ (làm tròn kết quả đến phần trăm)}
		 {Nếu tỉ lệ lạm phát trung bình là $6\%$ một năm thì sau $15$ năm sức mua của số tiền ban đầu chỉ còn lại một nửa}
	\loigiai{ 
		\begin{itemchoice}
			\itemch 
			Giả thiết $P = 100$ triệu đồng, $r = 7\%$, $n = 2$ năm.\\ 
			Ta  có
			 \[
			A = 100 \cdot 10^6 \left(1 - \dfrac{7}{100} \right)^2 = 86\,490\,000 \text{  đồng}.
			\] 
			\itemch 
			Vì sau hai năm sức mua còn lại của $100\,000\,000$ là $86\,490\,000$ đồng, không phải $96\,490\,000$ đồng. 
			\itemch 
			Giả thiết $P = 100$ triệu đồng, $A = 80$  triệu đồng, $n = 3$ năm.\\
			 Ta có
			 \[
			80 = 100 \left(1 - \dfrac{r}{100} \right)^3 \Rightarrow 1 - \dfrac{r}{100} = \sqrt[3]{\dfrac{4}{5}}  \Rightarrow r \approx 7{,}17.
			 \]
			Vậy tỉ lệ lạm  phát trung bình mỗi năm là $r \approx 7{,}17\%$.
			 \itemch
			Giả thiết cho $P = X$  triệu đồng, $A = \dfrac{X}{2}$ triệu đồng, $r = 6\%$.\\
			Ta  có
			\[ 
			\dfrac{X}{2} = X\left(1 - \dfrac{6}{100}  \right)^n \Rightarrow (0{,}94)^n = \dfrac{1}{2} \Rightarrow n \approx \log_{0{,}94}\left(\dfrac{1}{2} \right) \approx 11{,}2 \text{ (năm)}.
			 \]
			Vậy sau khoảng $12$ năm thì sức  mua của số tiền còn lại là một nửa.
		\end{itemchoice}
	} 
\end{ex}

% nguồn: Ok_dungsai_lop11_Lan2
\begin{ex}%[1D6V3-5]%[ma:d471aa63]
	Một người gửi $80$ triệu đồng vào ngân hàng theo hình thức  lãi kép với lãi suất $0{,}6\%$ mỗi tháng. Biết rằng người đó không rút tiền ra trong suốt quá trình gửi. Hỏi sau ít nhất bao nhiêu tháng thì người đó nhận được tổng số tiền (gốc và lãi) ít nhất là $100$ triệu đồng?
	 \choiceTF
		 {\True Số tiền sau $n$ tháng được tính theo công thức $A = 80(1{,}006)^n$}
		{Số tiền lãi sau tháng đầu tiên là $600$ nghìn  đồng}
		 {Sau $24$ tháng, số tiền người đó có được vượt quá $110$ triệu đồng}
		{\True Cần  ít nhất $37$ tháng để tổng số tiền đạt $100$ triệu đồng}
	\loigiai{ 
		\begin{itemchoice}
			\itemch Công thức tiền  gửi lãi kép hàng tháng là
			\[ 
			A =  P(1 + r)^n = 80(1{,}006)^n.
			\] 
			\itemch  Lãi tháng đầu tiên là
			\[ 
			80 \cdot 0{,}6\% = 80 \cdot  0{,}006 = 0{,}48\ \text{triệu} = 480\ \text{nghìn đồng}.
			 \]
			\itemch  Sau $24$ tháng:
			 \[
			 A = 80(1{,}006)^{24} \approx 80 \cdot 1{,}1275 \approx 90{,}2\ \text{triệu} < 110\ \text{triệu}.
			\] 
			\itemch  Giải bất phương trình:
			 \[
			80(1{,}006)^n \geq  100 \Rightarrow (1{,}006)^n \geq 1{,}25
			\Rightarrow n \geq \dfrac{\ln 1{,}25}{\ln 1{,}006}  \approx \dfrac{0{,}2231}{0{,}00598} \approx 37{,}3.
			 \]
			Do đó cần ít nhất $38$  tháng.
		\end{itemchoice}
	} 
\end{ex}

% nguồn: Ok_dungsai_lop11_Lan2
\begin{ex}%[1D6V3-5]%[ma:79a1fd2d]
	Số người trong cộng đồng sinh viên đã nghe một tin đồn nào đó là \( N = P\left(1 - \mathrm{e}^{-0{,}15d}\right) \), trong đó $P$ là tổng số sinh viên của cộng đồng và $d$ là số ngày trôi qua kể từ khi tin đồn bắt  đầu.
	 \choiceTF
		{Trong  một cộng đồng $1\,000$ sinh viên, sau ngày thứ nhất có $140$ sinh viên nghe được tin đồn}
		{\True  Trong một cộng đồng $1\,450$ sinh viên, để $525$ sinh viên nghe được tin đồn phải mất $3$ ngày}
		{Sau $3$ ngày $500$  sinh viên nghe được tin đồn thì cộng đồng có $1\,400$ sinh viên}
		{\True Số  ngày để một nửa cộng đồng sinh viên nghe được tiếng đồn là $5$ ngày}
	\loigiai{ 
		\begin{itemchoice}
			\itemch  Với $P = 1000$, $d = 1$:
			\[ 
			N_1 =  1000\left(1 - \mathrm{e}^{-0{,}15 \cdot 1}\right) \approx 139{,}292.
			 \]
			Suy ra sau ngày thứ  nhất có khoảng $139$ người.
			\itemch Với $P = 1\,450$, $d =  3$:
			\[ 
			N_3 =  1450\left(1 - \mathrm{e}^{-0{,}15 \cdot 3}\right) \approx 525{,}43.
			 \]
			Suy ra có khoảng $525$  sinh viên nghe tin.
			 \itemch Với $N = 500$, $d = 3$:
			 \[
			500  = P\left(1 - \mathrm{e}^{-0{,}15 \cdot 3}\right)
			\Rightarrow  P = \dfrac{500}{1 - \mathrm{e}^{-0{,}45}} \approx 1\,379{,}798.
			 \]
			Sau $3$  ngày $500$ sinh viên nghe được tiếng đồn thì cộng đồng có $1\,380$ sinh viên.
			\itemch  Với $N = \dfrac{P}{2}$:
			 \[
			 \dfrac{P}{2} = P\left(1 - \mathrm{e}^{-0{,}15d}\right)
			 \Rightarrow 1 - \mathrm{e}^{-0{,}15d} = \dfrac{1}{2}
			\Rightarrow  \mathrm{e}^{-0{,}15d} = \dfrac{1}{2}
			\Rightarrow d = \dfrac{\ln(2)}{0{,}15}  \approx 4{,}62.
			\] 
			Số ngày để một nửa cộng đồng  sinh viên nghe được tiếng đồn là $5$ ngày.
		\end{itemchoice}
	 }
\end{ex}

% nguồn: Ok_dungsai_lop11_Lan2
\begin{ex}%[1D6V3-5]%[ma:dafe03d2]
	Lạm phát là sự tăng mức giá chung một cách liên tục của hàng hoá và dịch vụ theo thời gian, tức là sự mất giá trị của một loại tiền tệ nào đó. Chẳng hạn, nếu lạm phát là $5\%$ mỗi năm thì sức mua của $1$ triệu đồng sau một năm chỉ còn là $950$ nghìn đồng (vì đã giảm  mất $5\%$ của $1$ triệu đồng, tức là $50$ nghìn đồng). Nói chung, nếu tỉ lệ lạm phát trung bình là $r\%$ một năm thì tổng số tiền $P$ ban đầu, sau $n$ năm sẽ còn lại có giá trị là
	\[ 
	A = P\left(1 - \dfrac{r}{100}\right)^n. 
	\] 
	 \choiceTF
		{\True Nếu tỉ lệ lạm phát trung bình là $7\%$ một năm thì tổng số tiền $P$ ban đầu sau $n$ năm sẽ còn lại có giá trị  là $A = P\left(1 - \dfrac{7}{100}\right)^n$}
		{Nếu tỉ lệ lạm phát là $8\%$ một năm thì sức mua của $100$ triệu  đồng sau hai năm sẽ còn lại khoảng $83{,}6$ triệu đồng}
		{\True Nếu sức mua của $100$ triệu đồng sau hai năm còn lại $90$ triệu đồng thì tỉ lệ lạm phát trung bình  của hai năm đó là $5{,}13\%$}
		{Nếu tỉ lệ lạm phát là $5\%$ mỗi năm thì sau ít nhất $14$ năm sức mua của số  tiền ban đầu chỉ còn lại một nửa}
	\loigiai{ 
		Sau $1$ năm số  tiền đó chỉ còn giá trị là $A_1=P-P\cdot \dfrac{r}{100}=P\left(1-\dfrac{r}{100}\right)$.\\
		Sau $2$ năm  số tiền đó chỉ còn giá trị là $A_2=A_1-A_1\cdot \dfrac{r}{100}=P\left(1-\dfrac{r}{100}\right)^2$.\\
		 $\ldots\ldots$\\
		 Sau $n$ năm số tiền đó chỉ còn giá trị là $A_n=P\left(1-\dfrac{r}{100}\right)^n$.
		\begin{itemchoice}
			\itemch Áp dụng công thức: 
			 \[
			A = P\left(1 - \dfrac{r}{100}\right)^n \Rightarrow A = P\left(1  - \dfrac{7}{100}\right)^n.
			 \]
			\itemch Với $P = 100$, $r = 8$, $n =  2$:
			 \[
			 A = 100\left(1 - \dfrac{8}{100}\right)^2 = 100(0{,}92)^2 \approx 84{,}64.
			 \]
			 Vậy nếu tỉ lệ lạm phát là $8\%$ một năm thì sức mua của $100$ triệu đồng sau hai năm sẽ còn lại khoảng $84{,}6$ triệu đồng.
			\itemch Với $P = 100$, $A = 90$,  $n = 2$ ta có
			\begin{eqnarray*}
				 &&90 = 100\left(1 - \dfrac{r}{100}\right)^2\\
				&\Rightarrow&\left(1 - \dfrac{r}{100}\right)^2  = 0{,}9\\
				&\Rightarrow&1 - \dfrac{r}{100} = \sqrt{0{,}9} \approx 0{,}9487\\ 
				 &\Rightarrow&r \approx 5{,}13\%.
			\end{eqnarray*}
			Vậy nếu sức mua của $100$ triệu đồng sau hai năm còn  lại $90$ triệu đồng thì tỉ lệ lạm phát trung bình của hai năm đó là $5{,}13\%$.
			\itemch  Với $A = \dfrac{P}{2}$, $r = 5$:
			\[ 
			\dfrac{1}{2} = \left(1 - \dfrac{5}{100}\right)^n \Rightarrow \dfrac{1}{2} =  0{,}95^n
			\Rightarrow n =  \dfrac{\ln(1/2)}{\ln(0{,}95)} \approx 14{,}21.
			 \]
			Cần khoảng $14$ năm và $3$ tháng, không phải ít nhất $14$ năm. 
		\end{itemchoice}
	 }
\end{ex}

% nguồn: Ok_dungsai_lop11_Lan2
\begin{ex}%[1D6V3-5]%[ma:64a38694]
	Ông $X$ gửi vào ngân hàng số tiền $300$ triệu  đồng theo hình thức lãi kép với lãi suất $6\%$/năm.
	 \choiceTF
		 {Số tiền lãi ông $X$ nhận được ở năm đầu tiên là $6$ triệu đồng}
		{\True Công thức tính số tiền ông $X$ nhận được cả gốc và lãi  sau $n$ năm gửi tiền là\break $T_n = 300\,000\,000\cdot (1+6\%)^n\text{ (đồng)}$}
		 {Số tiền ông $X$ nhận được sau $5$ năm là nhiều hơn $410$ triệu đồng}
		{\True Nếu ông $X$ muốn nhận được số tiền cả gốc lẫn lãi nhiều  hơn $500$ triệu đồng thì cần gửi ít nhất $9$ năm}
	 \loigiai{
		\begin{itemchoice}
			\itemch Vì số tiền lãi năm đầu tiên bằng số tiền gửi nhân với lãi suất: $300 \times 6\% = 18$ triệu đồng. 
			\itemch Áp dụng công thức $T_n = A(1 + r)^n$.\\ 
			Với $A =  300\,000\,000$, $r = 6\%$ nên suy ra $T_n = 300\,000\,000\cdot (1 + 6\%)^n$ đồng.
			\itemch Vì số tiền ông nhận được sau $5$ năm là $$T_5 = 300\,000\,000\cdot (1 + 6\%)^5 \approx 401\,467\,673\text{ đồng}.$$ Số  tiền này nhỏ hơn $410$ triệu đồng.
			 \itemch Ta có
			\begin{eqnarray*}
				&&T_n  > 500\,000\,000\\
				&\Leftrightarrow&300\,000\,000\cdot (1+6\%)^n > 500\,000\,000\\ 
				&\Leftrightarrow&n > \log_{1+6\%}  \dfrac{5}{3} \approx 8{,}77.
			\end{eqnarray*}
			Vậy ông $X$ cần gửi ít nhất  $9$ năm.
		\end{itemchoice}
	} 
\end{ex}

% nguồn: Ok_dungsai_lop11_Lan2
\begin{ex}%[1D6V3-5]%[ma:5dc7c475]
	Lạm phát là sự tăng mức giá chung một cách liên tục của hàng hoá và dịch vụ theo thời gian, tức là sự mất giá trị của một loại tiền tệ nào đó. Chẳng hạn, nếu lạm phát là $5\%$ một năm thì sức mua của $1$ triệu đồng sau một năm chỉ còn là $950$  nghìn đồng (vì đã giảm mất $5\%$ của $1$ triệu đồng, tức là $50\,000$ đồng). Nếu chúng ta biết tỉ lệ lạm phát trung bình là $r\%$ một năm thì tổng số tiền $P$ ban đầu, sau $n$ năm sẽ còn giá trị là
	\[ 
	A  = P\left(1 - \dfrac{r}{100}\right)^n.
	\] 
	 \choiceTF
		{Nếu tỉ lệ lạm phát là $6\%$ mỗi năm thì sức  mua của $100$ triệu đồng sau một năm sẽ còn là $95$ triệu đồng}
		{\True Nếu tỉ lệ lạm phát là $8\%$ mỗi năm thì sức mua của $100$ triệu đồng sau hai năm sẽ còn  lại là $84{,}64$ triệu}
		{\True Nếu sức mua của $100$ triệu đồng sau hai năm còn lại $90$ triệu đồng thì tỉ  lệ lạm phát trung bình của hai năm đó là khoảng $5{,}13\%$}
		{\True Nếu sau khoảng $14$ năm thì sức mua của số tiền  ban đầu còn lại một nửa nếu tỉ lệ lạm phát là $5\%$ một năm}
	 \loigiai{
		\begin{itemchoice}
			\itemch Với $r = 6$, $P = 100$, $n = 1$: 
			 \[
			A = 100\left(1 - \dfrac{6}{100}\right)^1 = 94 \text{  triệu đồng}.
			\] 
			Vậy sức mua của $100$ triệu đồng sau một năm với tỉ lệ  lạm phát là $6\%$ một năm chỉ còn lại khoảng $94$ triệu đồng.
			\itemch  Với $r = 8$, $P = 100$, $n = 2$:
			 \[
			A = 100\left(1 - \dfrac{8}{100}\right)^2  = 84{,}64 \text{ triệu đồng}.
			 \]
			Vậy sức mua của $100$ triệu đồng  sau hai năm với tỉ lệ lạm phát là $8\%$ một năm chỉ còn lại khoảng $84{,}64$ triệu đồng.
			\itemch Với $P  = 100$, $A = 90$, $n = 2$:
			 \[
			90 = 100\left(1 -  \dfrac{r}{100}\right)^2 \Rightarrow \left(1 - \dfrac{r}{100}\right)^2 = 0{,}9 \Rightarrow 1 - \dfrac{r}{100} = \sqrt{0{,}9}
			\Rightarrow  r \approx 5{,}13.
			\] 
			Vậy tỉ lệ  lạm phát trung bình của hai năm đó là khoàng $5{,}13\%$.
			\itemch Với $P = 1$,  $A = \dfrac{1}{2}$, $r = 5$:
			\[ 
			\dfrac{1}{2} = \left(1 - \dfrac{5}{100}\right)^n  = (0{,}95)^n
			 \Rightarrow \ln\dfrac{1}{2} = n \ln(0{,}95) \Rightarrow n = \dfrac{\ln\dfrac{1}{2}}{\ln(0{,}95)} \approx 13{,}51.
			\] 
			Vậy sau  khoảng $14$ năm thì sức mua của số tiền ban đầu sẽ chỉ còn lại một nửa nếu tỉ lệ lạm phát là $5\%$ một năm.
		\end{itemchoice}
	} 
\end{ex}

% nguồn: Ok_dungsai_lop11_Lan2
\begin{ex}%[1D6V3-5]%[ma:9ed637bc]
	Anh Nam vay tiền ngân hàng $1$ tỷ đồng theo phương thức trả góp với lãi suất $0{,}5\%$/tháng. Nếu cuối mỗi tháng bắt  đầu từ tháng thứ nhất anh Nam trả $30$ triệu đồng. Biết rằng lãi suất ngân hàng không thay đổi trong suốt thời gian trả nợ.
	\choiceTF 
		{\True Số tiền nợ sau  $8$ tháng là $796\,464\,780{,}4$}
		{\True Số tiền nợ sau $10$ tháng  là $744\,299\,339{,}8$}
		{\True  Sau $37$ tháng thì anh Nam trả hết nợ}
		{\True Nếu cuối mỗi tháng bắt  đầu từ tháng thứ nhất anh Nam trả $45$ triệu đồng thì sau hai năm anh Nam trả hết nợ}
	\loigiai{ 
		Gọi $a$ là số tiền vay,  $r$ là lãi suất, $m$ là số tiền thanh toán mỗi tháng.\\
		Số tiền nợ sau tháng  thứ nhất là $N_1=a(1+r)-m$.\\
		Số tiền nợ sau tháng thứ hai  là
		$$N_2=[a(1+r)-m]+[a(1-r)-m]r-m=a(1+r)^2-m[(1+r)+1].$$ 
		$\ldots  \ldots$\\
		Số tiền nợ  sau $n$ tháng là
		\[ 
		N_n = a(1 + r)^n - m \cdot \dfrac{(1 + r)^n -  1}{r}.\]
		Với $a = 10^9$, $r = 0{,}5\% = 0{,}005$, $m = 30 \cdot  10^6$.
		\begin{itemchoice}
			\itemch Số tiền nợ sau $8$ tháng  là
			 \[
			N_8 = 10^9 (1 + 0{,}005)^8 - 30 \cdot 10^6 \cdot \dfrac{(1  + 0{,}005)^8 - 1}{0{,}005} \approx 796464780{,}4.
			 \]
			\itemch Số tiền nợ sau $10$ tháng là 
			 \[
			N_{10} = 10^9  (1 + 0{,}005)^{10} - 30 \cdot 10^6 \cdot \dfrac{(1 + 0{,}005)^{10} - 1}{0{,}005} \approx 744\,299\,339{,}8.
			\] 
			 \itemch Tìm $n$ để trả hết nợ:
			 \[
			1\,000(1 +  0{,}005)^n - 30 \cdot \dfrac{(1 + 0{,}005)^n - 1}{0{,}005} = 0
			\Rightarrow n  \approx 36{,}55.
			\] 
			Vậy sau $37$ tháng anh  Nam trả hết nợ.
			\itemch Nếu $m = 45$, ta có 
			 \[
			1\,000(1 + 0{,}005)^n - 45 \cdot \dfrac{(1 + 0{,}005)^n - 1}{0{,}005} =  0
			\Rightarrow  n \approx 23{,}615.
			\] 
			Vậy sau hai năm  ($24$ tháng) anh Nam trả hết nợ.
		\end{itemchoice}
	} 
\end{ex}

% nguồn: Ok_dungsai_lop11_Lan2
\begin{ex}%[1D6V3-5]%[ma:5e6a1844]
	 Một người gửi tiết kiệm ngân hàng số tiền $110$ triệu đồng theo hình thức lãi suất kép với kì hạn $3$ tháng và lãi suất $5{,}4\%$/năm.
	 \choiceTF
		 {\True Sau một kì hạn người đó có số tiền lãi là $1\,485\,000$ đồng}
		{Sau $1$ năm không rút tiền thì  số tiền cả gốc lẫn lãi người đó nhận được ít hơn $116$ triệu đồng}
		{Sau $5$ năm tổng tiền cả vốn  lẫn lãi người đó nhận được không quá $143{,}25$ triệu đồng}
		{Tổng tiền cả vốn lẫn lãi người đó nhận được gấp đôi so với  ban đầu ít nhất là $12$ năm}
	\loigiai{ 
		\begin{itemchoice}
			\itemch Sau một kì hạn, người đó có  số tiền lãi là $100\,000\,000 \cdot \dfrac{5{,}4}{100} \cdot \dfrac{1}{4} = 1\,485\,000$ đồng.
			\itemch Sau $1$ năm không rút tiền, số tiền cả gốc lẫn lãi  là $100 \cdot \left(1 + \dfrac{0{,}054}{4}\right)^4 \approx 116{,}06$ triệu đồng, lớn hơn $116$ triệu.
			\itemch Sau $5$ năm, tức $20$  kì hạn, tổng tiền nhận được là $100 \cdot \left(1 + \dfrac{0{,}054}{4}\right)^{20} \approx 143{,}84$ triệu đồng, lớn hơn $143{,}25$ triệu.
			\itemch Gọi $n$ là số quý tối thiểu để tiền nhận được gấp đôi: $$110  \cdot \left(1 + \dfrac{0{,}054}{4}\right)^n = 2 \cdot 110 \Rightarrow n = \log_{1 + 0{,}054/4} 2 \approx 52.$$
			Do đó cần ít nhất $13$ năm ($52$  tháng) để tiền gấp đôi, không phải $12$ năm.
		\end{itemchoice}
	} 
\end{ex}


%%%%%%%%%%%%%%%% TRẢ LỜI NGẮN %%%%%%%%%%%%%%%%

% nguồn: Ok_TraLoiNgan_Lop11_Lan1
\begin{ex}%[1D6H3-5]%[ma:eed849b2]
	Số lượng của loại vi khuẩn A trong một phòng thí nghiệm được tính theo công thức $s(t)=s(0) \cdot 2^t$, trong đó $s(0)$ là số lượng vi khuẩn A lúc ban đầu, $s(t)$ là số lượng vi khuẩn A có sau $t$ phút. Biết sau $3$ phút thì số lượng vi khuẩn A là  $625$ con. Hỏi sau bao nhiêu phút, kể từ lúc ban đầu, số lượng vi khuẩn A là $10$ triệu con?
	 \shortans{7}
	 \loigiai{
		Ta có $s(3)=s(0) \cdot  2^3 \Rightarrow s(0)=\dfrac{s(3)}{8}=78{,}125$ nghìn con.\\
		Do đó $s(t)=10$ triệu con $=10\,000$  nghìn con khi
		$$10\,000=s(0) \cdot 2^t \Rightarrow 2^t=\dfrac{10\,000}{78{,}125}=128=2^7 \Rightarrow t=7\text{ phút}.$$ 
	 }
\end{ex}

% nguồn: Ok_TraLoiNgan_Lop11_Lan1
\begin{ex}%[1D6H3-5]%[ma:6a12b8a8]
	Giả sử số tiền gốc là $ A$, lãi suất là $ r\% /$ kì hạn gửi (có thể là tháng, quý hay năm) thì tồng số tiền nhận được cả gốc và lãi sau $ n$ kì hạn gửi là $ A{(1+r)^n}$. Bà Hạnh gửi 50 triệu vào tài khoản định kỳ tính  lãi kép với lãi suất là $ 8\% /$ năm. Tính số tiền lãi thu được sau 10 năm (làm tròn đến hàng phần mười).
	 \shortans{57,9}
	\loigiai{ 
		Áp dụng công thức tính lãi kép,  sau 10 năm số tiền cả gốc và lãi bà Hạnh thu về là $ A{(1+r)^n}=50(1+0,08)^{10}\approx 107,9$ triệu đồng.\\
	Suy ra số tiền lãi bà Hạnh thu về sau 10 năm là  $ 107{,}9-50=57{,}9$ triệu đồng.}
\end{ex}

% nguồn: Ok_TraLoiNgan_Lop11_Lan1
\begin{ex}%[1D6H3-5]%[ma:966f6d80]
	Số lượng của loại vi khuẩn A trong một phòng thí nghiệm được tính theo công thức $ s(t)=s(0)\cdot{2^t}$, trong đó $ s(0)$ là số lượng vi khuẩn $ A$ lúc ban đầu, $ s(t)$ là số  lượng vi khuẩn $ A$ có sau $ t$ phút. Biết sau 3 phút thì số lượng vi khuẩn $ A$ là 625 con. Hỏi sau bao nhiêu phút, kể từ lúc ban đầu, số lượng vi khuẩn A là 10 triệu con?\\
	\shortans{7} 
	\loigiai{ 
		Ta có $ s(3)=s(0)\cdot{2^3}\Rightarrow s(0)=\dfrac{s(3)}{8}=78,125$ nghìn  con.\\
		 Do đó $ s(t)=10$ triệu con $=10000$ nghìn con khi
		\begin{eqnarray*}
			& &10000=s(0)\cdot{2^t} \\ 
			 &\Leftrightarrow& {2^t}=\dfrac{10000}{78,125}=128=2^7\\
			&\Leftrightarrow & {2^t}=2^7\\ 
			&\Leftrightarrow & t=7. 
		\end{eqnarray*}
	} 
\end{ex}

% nguồn: Ok_TraLoiNgan_Lop11_Lan1
\begin{ex}%[1D6H3-5]%[ma:4b82dfc1]
	Một người gửi $10$ triệu đồng vào ngân hàng theo hình thức lãi kép có kì hạn là 12 tháng vối lãi suất $6\%$/năm. Giả sử qua các năm thì lãi suất không thay đổi và người đó không gửi thêm tiền vào mỗi năm. Để biết sau $y$ (năm) thì tổng số tiền cả vốn và lãi có được là $x$ (đồng), người đó sử dụng công thức $y=\log_{1,06}\left(\dfrac{x}{10}\right)$. Hỏi sau bao  nhiêu năm thì người đó có được tổng số tiền cả vốn và lãi là $20$ triệu đồng? (Làm tròn kết quả đến hàng đơn vị).
	 \shortans{12}
	\loigiai{ 
		Người đó có được tổng số tiền cả vốn lẫn lãi là $20$ triệu đồng sau 
		$$y=\log_{1,06}\left(\dfrac{20}{10}\right)\approx 12\text{ (năm).}$$ 
	} 
\end{ex}

% nguồn: Ok_TraLoiNgan_Lop11_Lan1
\begin{ex}%[1D6H3-5]%[ma:3db0e537]
	Giả sử số tiền gốc là $A$, lãi suất là $r\%$/kì hạn gửi (có thể là tháng, quý hay năm) thì tổng số  tiền nhận được cả gốc và lãi sau $n$ kì hạn gửi là $A(1+r)^n$. Bà Hạnh gửi $100$ triệu vào tài khoản định kỳ tính lãi kép với lãi suất là $8\%$/năm. Số tiền lãi bà Hạnh thu được sau $10$ năm có dạng $1a5{,}8b2$ triệu đồng, với $a$, $b$ là các số tự nhiên. Tính giá trị $T=a b+1$.
	\shortans{10} 
	\loigiai{ 
		Áp dụng công thức tính lãi kép, sau $10$ năm số tiền cả gốc và lãi bà  Hạnh thu về là
		$$A(1+r)^n=100(1+0{,}08)^{10} \approx 215{,}892\text{ triệu đồng}.$$ 
		Suy ra số tiền lãi bà Hạnh thu về sau $10$ năm là $215{,}892-100=115{,}892$ triệu đồng.\\ 
		Do đó $a=1$,  $b=9$. Suy ra $ab+1=10$.
	 }
\end{ex}

% nguồn: Ok_TraLoiNgan_Lop11_Lan2
\begin{ex}%[1D6H3-5]%[ma:59068c94]
	Thang đo Richte được Charles Francis đề xuất và sử dụng lần đầu tiên vào năm $1\,935$ để sắp xếp các số đo độ chấn động của các cơn động đất với đơn vị Richte. Công thức tính độ chấn động như sau: $M_L=\log A-\log A_0$, $M_L$ là độ chấn động, $A$ là biên độ tối đa được đo bằng địa chấn kế và $A_0$ là biên độ chuẩn. Hỏi theo thang độ Richte, cùng với một biên độ chuẩn thì  biên độ tối đa của một trận động đất $7$ độ Richte sẽ lớn gấp mấy lần biên độ tối đa của một trận động đất $5$ độ Richte?
	 \shortans{100}
	\loigiai{ 
		Trận động  đất $7$ độ Richte có $M_7=\log A_7-\log A_0=7 \Leftrightarrow A_7=A_0 \cdot 10^7$.\\
		Trận động đất $7$  độ Richte có
		$M_5=\log A_5-\log A_0=5 \Leftrightarrow A_5=A_0  \cdot 10^5$.\\
		Ta thấy  $\dfrac{A_7}{A_5}=\dfrac{A_0\cdot 10^7}{A_0\cdot 10^5}=10^2=100$.\\
		Vậy trận động đất $7$ độ  Richte có biên độ tối đa gấp $100$ lần trận động đất $5$ độ Richte.
	 }
\end{ex}

% nguồn: Ok_TraLoiNgan_Lop11_Lan2
\begin{ex}%[1D6H3-5]%[ma:56e506bf]
	Các khí thải gây hiệu ứng nhà kính là nguyên nhân chủ yếu làm Trái đất nóng lên. Theo OECD, khi nhiệt độ Trái đất tăng lên thì tổng giá trị kinh tế toàn cầu giảm. Người ta ước tính rằng, khi nhiệt độ Trái đất tăng thêm $2^\circ$C thì tổng giá trị kinh tế toàn cầu  giảm $3\%$; còn khi nhiệt độ Trái đất tăng thêm $5^\circ$C thì tổng giá trị kinh tế toàn cầu giảm $10\%$. Biết rằng, nếu nhiệt độ Trái đất tăng thêm $t^\circ$C, tổng giá trị kinh tế toàn cầu giảm $f(t)\%$ thì $f(t)=k\cdot a^t$, trong đó $k$, $a$ là các hằng số dương Khi nhiệt độ Trái đất tăng thêm bao nhiêu độ thì tổng giá trị kinh tế toàn cầu giảm đến $20\%$ (làm tròn kết quả đến hàng phần mười)?
	 \shortans{6,7}
	 \loigiai{
		Theo bài ra ta  có $\heva{&k\cdot a^2=3\% \\&k\cdot a^5=10\%}\Leftrightarrow \heva{&a=\sqrt[3]{\dfrac{10}{3}} \\&k=\dfrac{3\%}{a^2}.}$\\
		Tổng giá trị kinh tế toàn giảm đến $20\%$, tức là $f(t)=20$.\\ 
		Khi đó $\dfrac{3\%}{a^2}\cdot a^t=20\%\Leftrightarrow a^{t-2}=\dfrac{20}{3} \Leftrightarrow  t=2+\log_a \dfrac{20}{3}\approx 6{,}7$.
	 }
\end{ex}

% nguồn: Ok_TraLoiNgan_Lop11_Lan2
\begin{ex}%[1D6H3-5]%[ma:638060f8]
	Biết rằng vi khuẩn E. coli là vi khuẩn gây tiêu chảy  đường ruột, gây đau bụng dữ dội, ngoài ra cứ sau $20$ phút thì số lượng vi khuẩn tăng gấp đôi, nghĩa là số lượng tính theo công thức $S=S_n\cdot 2^n$, $S_0$ là số lượng ban đầu, $n$ là số lần nhân đôi. Ban đầu chỉ có $40$ con vi khuẩn nói trên trong đường ruột, hỏi sau bao lâu số lượng vi khuẩn là $671\,088\,640$ con?
	 \shortans{8}
	 \loigiai{
		Ta có $S=S_0  \cdot 2^n$.\\
		Suy ra $671\,088\,640=40  \cdot 2^n \Leftrightarrow 2^n=\dfrac{671\,088\,640}{40}=16\,777\,216\Leftrightarrow n=\log_2(16\,777\,216)=24$.\\
		Vậy để số lượng  vi khuẩn là $671\,088\,640$ con thì vi khuẩn đã có $24$ lần nhân đôi.\\
		Thời gian để thực hiện $24$ lần nhân đôi đó là $24\cdot  20=480$ phút$=8$ giờ.
	} 
\end{ex}

% nguồn: Ok_TraLoiNgan_Lop11_Lan2
\begin{ex}%[1D6H3-5]%[ma:40716c68]
	Năm $2\,022$, bạn Yến mua một chiếc điện thoại Iphone $14$ Pro Max với giá $35$ triệu đồng. Công thức $P(t)=35\cdot \left(\dfrac{5}{6}\right)^{\tfrac{t}{15}}$ được dùng để tính giá trị còn lại của một chiếc điện thoại sau thời gian t kể từ khi sử dụng. Hỏi giá trị còn lại của chiếc điện thoại  sau $1$ năm $2$ tháng sử dụng là bao nhiêu (làm tròn kết quả đến hàng đơn vị)?
	 \shortans{30}
	 \loigiai{
		Ta có $1$ năm  $2$ tháng $=14$ tháng.\\
		Giá trị còn lại của chiếc điện thoại sau $1$ năm $2$ tháng sử dụng là 
		\[P(14)=35\cdot \left(\dfrac{5}{6}\right)^{\tfrac{14}{15}}  \approx 30\ (\text{triệu đồng}).\]
	} 
\end{ex}

% nguồn: Ok_TraLoiNgan_Lop11_Lan2
\begin{ex}%[1D6H3-5]%[ma:2b9f602a]
	Biết rằng năm  $2\,023$, dân số Việt Nam là $78\,685\,800$ người và tỉ lệ tăng dân số năm đó là $1{,}7\%$. Cho biết sự tăng dân số được ước tính theo công thức $S=A\mathrm{e}^{Nr}$, cứ tăng dân số với tỉ lệ như vậy thì đến năm nào dân số nước ta ở mức $120$ triệu người.
	\shortans{2048} 
	 \loigiai{
		 Theo đề bài ta có $A=78\,685\,800$, $r=0{,}017$.\\
		Dân số nước  ta ở mức $120$ triệu đồng khi
		$120\,000\,000=78\,685\,800  \mathrm{e}^{0{,}017N}\Leftrightarrow N\approx 24{,}8$ năm.\\
		Vậy đến năm $2\,048$ thì  dân số nước ta ở mức $120$ triệu người.
	} 
\end{ex}

% nguồn: Ok_TraLoiNgan_Lop11_Lan2
\begin{ex}%[1D6H3-5]%[ma:7e51b1fd]
	Cường độ một trận động đất $M$ được cho bởi công thức $M=\log\dfrac{A}{A_0}$, với $A$ là  biên độ rung chấn tối đa và $A_0$ là biên độ chuẩn. Đầu thế kỉ XX, một trận động đất ở San Francisco có cường độ $8{,}3$ độ richter. Trong cùng năm đó, trận động đất khác ở Nam Mỹ có biên độ mạnh hơn gấp $4$ lần. Tính cường độ của trận động đất ở Nam Mỹ (làm tròn kết quả đến hàng phần mười).
	 \shortans{8,9}
	\loigiai{ 
		Vì trận động đất ở San Francisco có  cường độ $8{,}3$ độ richter nên ta có $\log\dfrac{A_1}{A_0}=8{,}3$.\\
		Trận động đất khác ở Nam Mỹ có biên độ mạnh hơn  gấp $4$ lần có cường độ là
		\[\log\dfrac{A_2}{A_0}=\log\dfrac{4\cdot A_1}{A_0}=\log 4+\log\dfrac{A_1}{A_0}=\log 4+8{,}3\approx 8{,}9\ (\text{độ  Richter}).\]
	} 
\end{ex}

% nguồn: Ok_TraLoiNgan_Lop11_Lan2
\begin{ex}%[1D6H3-5]%[ma:61cb6cc0]
	Tờ tiền mệnh giá $500\,000$ VND có kích thước chiều dài $1{,}52\cdot 10^{-1}$ m; chiều rộng $6{,}5\cdot 10^{-2}$ m; bề dày $10^{-4}$ m; nặng $10^{-3}$ kg. Ngày $05/07/2023$ công ty Xổ số điện toán Việt Nam thông báo ông An ở thành phố  Thái Bình trúng thưởng trị giá $39$ tỷ đồng. Công ty Xổ số điện toán Việt Nam đã trả thưởng cho ông An bằng tiền mặt toàn loại tiền mệnh giá $500\,000$ VND. Ông An nhận được số kilogam tiền là
	\shortans{78} 
	\loigiai{ 
		Tổng số tiền ông an nhận được là  $39 \cdot 10^9$ VND.\\
		Số tờ tiền mệnh giá $500\,000$ VND mà ông An nhận được là $\dfrac{39\cdot 10^9}{5\cdot  500\,000}=7{,}8 \cdot 10^4$ (tờ).\\
		Tổng khối lượng tiền mà ông An  nhận được là $7{,}8\cdot 10^4\cdot 10^{-3}=78$ (kg).
	 }
\end{ex}

% nguồn: Ok_TraLoiNgan_Lop11_Lan2
\begin{ex}%[1D6H3-5]%[ma:e9c616f1]
	Thực hiện một mẻ nuôi cấy vi khuẩn  với $1\,500$ vi khuẩn ban đầu, nhà sinh học phát hiện số lượng vi khuẩn tăng thêm $25\%$ sau mỗi hai ngày. Công thức $P(t)=P_0\cdot a^t$ cho phép tính số lượng vi khuẩn của mẻ nuôi cấy sau $t$ ngày kể từ thời điểm ban đầu. Sau $6$ ngày thì số lượng vi khuẩn gần bằng bao nhiêu (làm tròn kết quả đến hàng đơn vị)?
	 \shortans{2930}
	 \loigiai{
		Vì ban  đầu có $1\,500$ vi khuẩn nên $P_0=1\,500$.\\
		Sau mỗi $2$ ngày,  số lượng vi khuẩn tăng thêm $25\%$, tức là \[P(2)=1\,500\cdot 1{,}25=1\,875.\]
		Mà $P(2)=1\,500\cdot a^2$ nên $a^2= 1{,}25$.\\ 
		Số lượng vi khuẩn sau  $6$ ngày là
		\[P(6)=P_0\cdot a^6=1\,500 \cdot  \left(a^2\right)^3=1\,500 \cdot (1{,}25)^3\approx 2\,930.\]
	 }
\end{ex}

% nguồn: Ok_TraLoiNgan_Lop11_Lan2
\begin{ex}%[1D6H3-5]%[ma:4740d1a3]
	Chu kì bán rã của Cacbon $^{14}\mathrm{C}$ là khoảng $5\,730$ năm. Một vật có khối lượng Cacbon $^{14}\mathrm{C}$ ban đầu là $m_0$ thì sau khoảng thời gian $t$ năm, khối lượng Cacbon $^{14}\mathrm{C}$ còn lại của vật đó là $m(t)=m_0 \left(\dfrac{1}{2}\right)^{\tfrac{t}{T}}$. Các nhà khảo cổ đã tìm được trong một mẫu đồ  cổ một lượng Cacbon và xác định được nó đã mất khoảng $25\%$ lượng Cacbon $^{14}\mathrm{C}$ ban đầu của nó. Hỏi mẫu đồ cổ đó có tuổi là bao nhiêu?
	\shortans{2387} 
	 \loigiai{
		Khối lượng của mẫu đồ cổ khi được tìm  thấy là
		\[m(t)=m_0-\dfrac{25}{100}m_0=\dfrac{3}{4}m_0.\] 
		Ta  có $m(t)=m_0\cdot\left(\dfrac{1}{2}\right)^{\tfrac{t}{5\,730}} \Leftrightarrow \dfrac{3}{4}m_0=m_0 \left(\dfrac{1}{2}\right)^{\tfrac{t}{5\,730}} \Leftrightarrow t\approx 2\,387$.\\
		 Vậy mẫu đồ cổ đó có tuổi là $2\,387$ năm.
	 }
\end{ex}

% nguồn: Ok_TraLoiNgan_Lop11_Lan2
\begin{ex}%[1D6H3-5]%[ma:0e3360cb]
	Đầu năm $2\,022$, anh Luân có mua một chiếc xe ô tô Vinfast Lux A 2.0 với giá $948$ triệu đồng để chạy taxi. Biết mỗi tháng xe ô tô hao mòn mất $0{,}6\%$ giá trị xe đồng  thời làm ra được $20$ triệu đồng. Hỏi sau một năm, tổng giá trị chiếc xe và số tiền chiếc xe chạy taxi kiếm được là bao nhiêu triệu đồng (làm tròn kết quả đến hàng đơn vị)?
	\shortans{1122} 
	 \loigiai{
		Số tiền anh Luân kiếm được sau $1$ năm ($12$ tháng) là $20\cdot 12=240$  (triệu đồng).\\
		Sau một năm, giá trị xe còn lại  là $948\cdot (1-0{,}006)^{12}=948\cdot (0{,}994)^{12}$ (triệu đồng).\\
		Sau một năm, tổng giá trị chiếc xe và số tiền chiếc xe chạy taxi kiếm được là  \[240+948\cdot (0{,}994)^{12} \approx 1122 (\text{triệu đồng}).\]
	} 
\end{ex}

% nguồn: Ok_TraLoiNgan_Lop11_Lan2
\begin{ex}%[1D6H3-5]%[ma:ac7bd62b]
	Biết rằng nếu gọi các loại cây xanh trong quá trình quang hợp sẽ nhận được một lượng nhỏ cacbon $14$. Khi một bộ phận của cây xanh đó bị chết thì hiện tượng quang hợp cũng dừng và  nó sẽ không nhận thêm cacbon $14$ nữa. Lượng cacbon $14$ của bộ phận đó sẽ phân hủy một cách chậm chạp và chuyển hóa thành nitơ $14$. Biết rằng nếu gọi $P(t)$ là số phần trăm cacbon $14$ còn lại trong một bộ phận của một cây sinh trưởng từ $t$ năm trước đây thì $P(t)$ được tính theo công thức $P(t)=100\cdot \left(\dfrac{1}{2}\right)^{\tfrac{t}{5\,750}}$ ($\%$). Phân tích mẫu gỗ từ một công trình kiến trúc cổ, người ta thấy lượng cacbon $14$ còn lại trong mẫu gỗ đó là $65\%$. Hãy xác định niên đại của công trình đó.
	 \shortans{3574}
	 \loigiai{
		Theo đề bài ta có  $P(t)=65\Leftrightarrow 65=100\cdot \left(\dfrac{1}{2}\right)^{\tfrac{t}{5\,750}} \Leftrightarrow 0{,}65=\left(\dfrac{1}{2}\right)^{\tfrac{t}{5\,750}}\Leftrightarrow t\approx 3\,574$.\\
		Vậy tuổi của công trình kiến trúc đó là $3\,574$  năm.
	 }
\end{ex}

% nguồn: Ok_TraLoiNgan_Lop11_Lan2
\begin{ex}%[1D6H3-5]%[ma:45b5f41d]
	Vào cuối năm $2\,022$, báo Rossiyskaya Gazeta dẫn lời Bộ trưởng Tài nguyên Nga  cảnh báo nước này sẽ cạn kiệt dầu mỏ sau $28$ năm nữa nếu sản lượng khai thác hằng năm vẫn giữ như năm $2\,022$. Bắt đầu từ năm $2\,023$, nếu nước Nga mỗi năm giảm sản lượng khai thác $2\%$ so với năm trước thì sau bao nhiêu năm nữa nước này cạn kiệt dầu mỏ (làm tròn kết quả đến hàng đơn vị)?
	 \shortans{41}
	\loigiai{ 
		Gọi $S$ là lượng dầu mỏ còn  lại của Nga trên thực thế tính từ cuối năm $2\,022$.\\
		Gọi $x$ là sản lượng khai  thác hằng năm như năm $2\,022$.\\
		Theo bài ra  ta có $S=28x$.\\
		Gọi $n$ là số năm khai thác còn lại với sản Lượng  khai thác thay đổi hằng năm tính từ $2\,023$.\\
		Tổng lượng khai thác sau $n$  năm là $S_n=x \dfrac{1-0{,}98^n}{1-0{,}98}=50x(1-0{,}98^n)$.\\
		Khi khai thác hết ta có $50x(1-0{,}98^n)=28x\Leftrightarrow 0{,}98^n=0{,}44  \Leftrightarrow n=\log_{0{,}98}(0{,}44) \approx 41$.
	} 
\end{ex}

% nguồn: Ok_TraLoiNgan_Lop11_Lan1
\begin{ex}%[1D6V3-5]%[ma:ac8c5040]
	Trong một phòng thí nghiệm, người ta nuôi một loại vi khuẩn. Ban đầu có $300$ vi khuẩn và sau $1$ giờ thì số vi khuẩn là $705$ con. Giả sử số vi khuẩn tại thời điểm $x$ (giờ) được tính theo hàm mũ là $f(x) = C\mathrm{e}^{kx}$ với $C$ và $k$ là các hằng số thực dương. Xác định số vi khuẩn  sau $5$ giờ tính theo đơn vị nghìn con (kết quả lấy $1$ chữ số sau dấu thập phân).
	 \shortans{21,5}
	\loigiai{ 
		Ta có $300 = f(0) = C$ và $705 = f(1) = C\mathrm{e}^k = 300\mathrm{e}^k$ nên $\mathrm{e}^k = \dfrac{705}{300} = 2{,}35$ hay $f(x) =  300\cdot (2{,}35)^x$.\\
		Do đó $f(5) = 300\cdot (2{,}35)^5 \approx 21501 \approx 21{,}5$  nghìn con.
	 }
\end{ex}

% nguồn: Ok_TraLoiNgan_Lop11_Lan1
\begin{ex}%[1D6V3-5]%[ma:9bb9561f]
	Trong một nghiên cứu, một nhóm học sinh được cho xem cùng một danh sách các loài động vật và được kiểm tra lại xem họ còn nhớ bao nhiêu phần trăm danh sách đó sau mỗi tháng. Giả sử sau $t$  tháng, khả năng nhớ trung bình của nhóm học sinh đó được tính theo công thức:
	$$ 
	M(t)=75-20 \ln (t+1), 0 \leq t  \leq 12
	 $$
	 (đơn vị:$\%$)\\
	Đến tháng thứ  mấy thì nhóm học sinh đó nhớ được khoảng một nửa danh sách các loài động vật đã xem?
	 \shortans{3}
	 \loigiai{
		Theo giả thiết, học sinh nhớ được  một nửa danh sách các loài vật, tức là
		 $$
		M(t)=50(\%) 
		$$ 
		Ta có: $50=75-20 \ln (t+1) \Leftrightarrow \ln (t+1)=\dfrac{5}{4} \Leftrightarrow  t+1=e^{\tfrac{5}{4}} \Leftrightarrow t \approx 2,49$ (tháng).\\
		Vậy qua tháng thứ $3$ khả năng nhớ  trung bình của nhóm học sinh là $50 (\%)$.
	 }
\end{ex}

% nguồn: Ok_TraLoiNgan_Lop11_Lan1
\begin{ex}%[1D6V3-5]%[ma:41b4880b]
	Trong vật lí, sự phân rã các chất phóng xạ được cho  bởi công thức:
	 $$
	 m(t)=m_0\left(\dfrac{1}{2}\right)^{\tfrac{t}{T}}
	 $$
	Trong đó, $m_0$ là khối lượng chất phóng  xạ ban đầu (tại thời điểm $t=0$), $m(t)$ là khối lượng chất phóng xạ tại thời điểm $t$ và $T$ là chu kì bán rã. Hạt nhân Poloni $(P o)$ là chất phóng xạ $a$ có chu kì bán rã $138$ ngày. Giả sử lúc đầu có $100$ Poloni. Tính khối lượng Poloni còn lại sau $100$ ngày theo đơn vị gam (làm tròn kết quả đến phần chục).
	 \shortans{60,5}
	 \loigiai{
		Khối lượng  Poloni còn lại sau $100$ ngày là: $m(100)=100 \cdot\left(\dfrac{1}{2}\right)^{\tfrac{100}{138}} \approx 60,5$(g).
	} 
\end{ex}

% nguồn: Ok_TraLoiNgan_Lop11_Lan1
\begin{ex}%[1D6V3-5]%[ma:6f7a5e54]
	Một người lập kế hoạnh gửi tiết kiệm ngân hàng như sau: Đầu tháng 1 năm 2023, người đó gửi 10 triệu đồng; sau mỗi đầu tháng tiếp theo, người đó gửi số tiền nhiều hơn $10\%$ so với số tiền đã gửi ở  tháng liền trước đó. Biết rằng lãi suất ngân hàng không đổi là $0{,}5\%$ mỗi tháng và được tính theo hình thức lãi kép. Với kế hoạnh như vậy, đến hết tháng 12 năm 2024, số tiền của người đó trong tài khoản tiết kiệm là bao nhiêu triệu đồng?
	\shortans{922} 
	 \loigiai{
		Với $A=10$ triệu, $a=0{,}1$,  $r=0{,}005$.\\
		Đầu tháng 2:  $A(1+r)+A(1+a)$.\\
		Đầu tháng 3:  $A(1+r)^2+A(1+a)(1+r)+A(1+a)^2$.\\
		 Đầu tháng 4: $A(1+r)^3+A(1+a)(1+r)^2+A(1+a)^2(1+r)+A(1+a)^3$.\\
		 ….\\
		Đầu tháng $n$: $A\left[(1+r)^{n-1}+(1+r)^{n-2}(1+a)+\cdots  +(1+r)(1+a)^{n-2}+(1+a)^{n-1}\right]$.\\
		Hết  tháng $n$: $A\left[(1+r)^{n-1}+(1+r)^{n-2}(1+a)+\cdots +(1+r)(1+a)^{n-2}+(1+a)^{n-1}\right](1+r)$.\\
		Gọi $B$ là số tiền của  người đó trong tài khoản tiết kiệm đến hết tháng 12 năm 2024.
		Khi đó  $n=24$.\\
	Ta có $B=A\cdot\dfrac{(1+a)^n-(1+r)^n}{a-r}(1+r)\approx  922{,}7$.}
\end{ex}

% nguồn: Ok_TraLoiNgan_Lop11_Lan1
\begin{ex}%[1D6V3-5]%[ma:e4cac3bf]
	Nếu tỉ lệ lạm phát trung bình hằng năm là $4 \%$ thì chi phí $C$ cho việc mua một loại hàng hoá hoặc sử dụng một dịch vụ nào đó sẽ được mô hình hoá bằng công  thức:
	$$C(t)=P(1+0,04)^t,$$ 
	trong đó $t$ là thời gian (tính bằng năm) kể từ  thời điểm hiện tại và $P$ là chi phí hiện tại cho hàng hoá hoặc dịch vụ đó.\\
	Giả sử hiện tại chi phí cho mỗi lần thay dầu ô tô là $ 800 $ nghìn đồng. Hãy ước tính chi phí cho mỗi  lần thay dầu ô tô sau $ 5 $ năm nữa (kết quả tính theo đơn vị nghìn đồng và làm tròn đến hàng đơn vị).
	 \shortans{9730}
	\loigiai{ 
		Chi phí cho mỗi lần thay dầu ô tô sau  $ 5 $ năm nữa là:
		$$C(5)=800(1+0,04)^5 \approx 973\,  \text{(nghìn đồng).}$$
	} 
\end{ex}

% nguồn: Ok_TraLoiNgan_Lop11_Lan2
\begin{ex}%[1D6V3-5]%[ma:46c1abb3]
	Sự tăng trưởng của một loại vi khuẩn tuân theo công thức $S=A\cdot\mathrm{e}^{rt}$, trong đó $A$ là số lượng vi khuẩn ban đầu, $r$ là tỉ lệ tăng trưởng, $t$ là thời gian tăng trưởng. Biết rằng số lượng vi khuẩn ban đầu là $100$  con và sau $5$ giờ có $300$ con. Số vi khuẩn sau $10$ giờ bằng bao nhiêu?
	\shortans{900} 
	\loigiai{ 
		Trước tiên, ta tìm tỉ lệ tăng trưởng mỗi giờ của loại vi  khuẩn này.\\
		Từ giả thiết ta có $300=100\cdot\mathrm{e}^{5r}\Leftrightarrow r=\dfrac{\ln 300-\ln  100}{5}=\dfrac{\ln 3}{5}$.\\
	Sau $10$ giờ, từ $100$ con vi  khuẩn sẽ có $100\cdot\mathrm{e}^{10\cdot\dfrac{\ln 3}{5}}=900$ con.}
\end{ex}

% nguồn: Ok_TraLoiNgan_Lop11_Lan2
\begin{ex}%[1D6V3-5]%[ma:d83300fc]
	Dân số thế giới  được ước tính theo công thức $P_n=P_0\cdot\mathrm{e}^{nr}$, trong đó $P_0$ là dân số của năm lấy làm mốc, $P_n$ là dân số sau $n$ năm, $r$ là tỉ lệ tăng dân số hàng năm. Biết rằng năm $2\,001$ dân số Việt Nam là $76\,685\,800$ người và tỉ lệ tăng dân số năm đó là $1{,}7\%$. Hỏi cứ tăng dân số với tỉ lệ như vậy thì đến năm nào dân số nước ta ở mức $115$ triệu người?
	\shortans{2,03} 
	 \loigiai{
		Theo bài ra ta xét phương trình 
		$$P_n=115\cdot 10^6\Leftrightarrow P_0\cdot\mathrm{e}^{nr}=115\cdot 10^6\Leftrightarrow\ln P_0+nr=\ln({115\cdot 10}^6).$$ 
		Suy ra  $n=\dfrac{\ln({115\cdot 10}^6)-\ln P_0}{r}\approx 23{,}8$.\\
	Như vậy đến năm $2\,025$  dân số nước ta sẽ ở mức $115$ triệu người.}
\end{ex}

% nguồn: Ok_TraLoiNgan_Lop11_Lan2
\begin{ex}%[1D6V3-5]%[ma:6aeb465e]
	Cường độ một trận động đất được cho bởi công thức $M=\log A-\log A_0$ độ Richter, với $A$ là biên độ rung chấn tối đa và $A_0$ là một biên độ chuẩn (hằng số). Đầu thế kỷ $20$, một trận  động đất ở San Francisco có cường độ đo được $8$ độ Richter. Trong cùng năm đó, trận động đất khác ở Nhật Bản có cường độ đo được $6$ độ Richer. Hỏi trận động đất ở San Francisco có biên độ gấp bao nhiêu lần biên độ trận động đất ở Nhật Bản?
	\shortans{100} 
	\loigiai{ 
		Nhận thấy ở San  Francisco trận động đất có cường độ là $M_1=\log A_1-\log A_0=\log\dfrac{A_1}{A_0}=8$.\\
		Ở  Nhật Bản trận động đất có cường độ là $M_2=\log\dfrac{A_2}{A_0}=6$.\\
	Khi đó $8-6=\log\dfrac{\mathrm{A}_1}{A_0}-\log\dfrac{A_2}{A_0}=\log\dfrac{A_1}{A_2}\Leftrightarrow  2=\log\dfrac{A_1}{A_2}\Leftrightarrow\dfrac{A_1}{A_2}=10^2=100$.}
\end{ex}

% nguồn: Ok_TraLoiNgan_Lop11_Lan2
\begin{ex}%[1D6V3-5]%[ma:5b4839c9]
	Sự tăng trưởng của loại vi khuẩn tuân theo công thức $S=A\mathrm{e}^{r\cdot t}$, trong đó $A$ là số lượng vi khuẩn ban đầu, $r$ là tỉ lệ tăng trưởng ($r>0$), $t$ là thời gian tăng trưởng. Biết số vi khuẩn ban đầu là  $100$ con và sau $5$ giờ có $300$ con. Thời gian để số vi khuẩn tăng gấp đôi số vi khuẩn ban đầu bằng bao nhiêu phút? (Kết quảlàm tròn đến hàng đơn vị)
	 \shortans{189}
	\loigiai{ 
		 Vì sau $5$h có $300$ con vi khuẩn, nên suy ra $300=100\cdot\mathrm{e}^{5r}\Leftrightarrow r=\dfrac{\ln 3}{5}$.\\
		Để vi  khuẩn tăng gấp đôi thì ta có phương trình
		$$200=100\cdot\mathrm{e}^{\tfrac{1}{5}\ln 3\cdot t}\Leftrightarrow\mathrm{e}^{\tfrac{1}{5}\ln 3\cdot t}=2\Leftrightarrow 3^{\tfrac{1}{5}\cdot t}=2\Leftrightarrow  t=5\log_32\Leftrightarrow t\approx 3{,}15.$$
	Vậy thời gian để số vi khuẩn tăng gấp đôi số vi khuẩn ban đầu là $3$ giờ  $9$ phút.}
\end{ex}

% nguồn: Ok_TraLoiNgan_Lop11_Lan2
\begin{ex}%[1D6V3-5]%[ma:3b70cb81]
	Giả sử số lượng một bầy ruồi tại thời điểm $t$ được tính theo công thức là $N(t)=N_o\cdot\mathrm{e}^{kt}$, trong đó $N_0$ là số lượng ruồi tại thời điểm $t=0$ và $k$ là hằng số tăng trưởng của bầy ruồi. Biết  số lượng ruồi tăng lên gấp đôi sau $9$ ngày và biết $N_0=100$ con. Hỏi sau bao nhiêu ngày bầy ruồi có $800$ con?
	 \shortans{24}
	\loigiai{ 
		 Ta có $2N_0=N_0\cdot\mathrm{e}^{9k}\Leftrightarrow k=\dfrac{\ln 2}{9}$.\\
	Để được  $800$ con ruồi, ta có $800=100\cdot\mathrm{e}^{t\cdot\dfrac{\ln 2}{9}}\Leftrightarrow t=\dfrac{\ln 8}{\ln 2}\cdot 9=27$ ngày.}
\end{ex}

% nguồn: Ok_TraLoiNgan_Lop11_Lan2
\begin{ex}%[1D6V3-5]%[ma:cfe617e9]
	Người ta thả một lượng bèo vào một hồ nước. Kết quả cho thấy sau $9$ giờ bèo sẽ sinh sôi kín cả mặt hồ. Biết rằng sau mỗi giờ, lượng bèo tăng gấp $10$ lần lượng bèo trước đó và tốc độ tăng không đổi. Hỏi sau mấy giờ thì lượng bèo  phủ kín $\dfrac{1}{3}$ mặt hồ? (kết quả làm tròn đến hàng phần trăm)
	 \shortans{8,52}
	\loigiai{ 
		Gọi $A$ là lượng bèo ban đầu. Sau  mỗi giờ, lượng lá bèo tăng gấp $10$ lần nên sau $9$ giờ ta lượng bèo là $A\cdot10^9$.\\
		 Gọi $t$ là số giờ để lượng bèo trong hồ phủ kín $\dfrac{1}{3}$ mặt hồ. Khi đó ta có
	$$A\cdot10^t=\dfrac{1}{3}\times A\cdot 10^9\Rightarrow t=\log\dfrac{10^9}{3}=9-\log  3\approx 8{,}52.$$}
\end{ex}

% nguồn: Ok_TraLoiNgan_Lop11_Lan2
\begin{ex}%[1D6V3-5]%[ma:97364edf]
	Một nguồn âm đẳng hướng đặt tại điểm $O$ có công suất truyền âm không đổi. Mức cường độ âm tại điểm $M$ cách $O$ một khoảng $R$ được tính bởi công thức $L_M=\log\dfrac{k}{R^2}$ (Ben) với $k$ là hằng số. Biết điểm $O$ thuộc đoạn thẳng $AB$ và mức  cường độ âm tại $A$ và $B$ lần lượt là $L_A=3$ (Ben) và $L_B=5$ (Ben). Tính mức cường độ âm tại trung điểm $AB$ (làm tròn đến $2$ chữ số sau dấu phẩy)
	\shortans{3,69} 
	 \loigiai{
		Ta  có $L_A<L_B\Rightarrow OA>OB$.\\
		Gọi  $I$ là trung điểm $AB$. Ta có $L_A=\log\dfrac{k}{OA^2}\Rightarrow\dfrac{k}{OA^2}=10^{L_A}\Rightarrow OA=\dfrac{\sqrt{k}}{\sqrt{10}^{L_A}}$.\\
		 $L_B=\log\dfrac{k}{OB^2}\Rightarrow\dfrac{k}{OB^2}=10^{L_B}\Rightarrow OB=\dfrac{\sqrt{k}}{\sqrt{10}^{L_B}}$.\\
		$L_I=\log\dfrac{k}{OI^2}\Rightarrow\dfrac{k}{OI^2}=10^{L_I}\Rightarrow OI=\dfrac{\sqrt{k}}{\sqrt{10}^{L_I}}$.\\ 
		Ta có $OI=\dfrac{1}{2}(OA-OB)\Rightarrow\dfrac{\sqrt{k}}{\sqrt{10}^{L_I}}=\dfrac{1}{2}\left(\dfrac{\sqrt{k}}{\sqrt{10}^{L_A}}-\dfrac{\sqrt{k}}{\sqrt{10}^{L_B}}\right)\Rightarrow\dfrac{1}{\sqrt{10}^{L_I}}=\dfrac{1}{2}\left(\dfrac{1}{\sqrt{10}^{L_A}}-\dfrac{1}{\sqrt{10}^{L_B}}\right)$  \\
	 $ \Rightarrow L_I=-2\log\left[\dfrac{1}{2}\left(\dfrac{1}{\sqrt{10}^{L_A}}-\dfrac{1}{\sqrt{10}^{L_B}}\right)\right]\Rightarrow L_I\approx 3{,}69 $.}
\end{ex}

% nguồn: Ok_TraLoiNgan_Lop11_Lan2
\begin{ex}%[1D6V3-5]%[ma:a25ffdc0]
	Sự tăng trưởng của một loại vi khuẩn trong phòng thí nghiệm được tính theo công thức $S(t)=S_0\cdot\mathrm{e}^{r\cdot t}$. Trong đó $S_0$ là số lượng vi  khuẩn ban đầu, $S(t)$ là số lượng vi khuẩn có sau $t$ (phút), $r$ là tỷ lệ tăng trưởng $(r>0)$, $t$ (tính theo phút) là thời gian tăng trưởng. Biết rằng số lượng vi khuẩn ban đầu có $500$ con và sau $5$ giờ có $1\,500$ con. Hỏi sau bao nhiêu giờ kể từ lúc ban đầu có $500$ con để số lượng vi khuẩn đạt $121\,500$ con?
	 \shortans{25}
	 \loigiai{
		Ta có $S_0=500$ (con);  $5$ giờ bằng $ 300$ phút.\\
		Sau $5$ giờ số  vi khuẩn là $S(300)=500\cdot\mathrm{e}^{300r}\Leftrightarrow 1500=500\cdot\mathrm{e}^{300r}\Leftrightarrow r=\dfrac{\ln 3}{300}$.\\
		Vậy khoảng thời gian $t$ kể từ lúc bắt đầu có $500$ con vi khuẩn đến khi số lượng vi  khuẩn đạt $121\,500$ con thỏa mãn $121\,500=500\cdot\mathrm{e}^{r\cdot t}
	\Leftrightarrow t=\dfrac{\ln 243}{r}=\dfrac{300\ln  243}{\ln 3}=1\,500 $ (phút) $=25$ (giờ).}
\end{ex}

% nguồn: Ok_TraLoiNgan_Lop11_Lan2
\begin{ex}%[1D6V3-5]%[ma:15ed6c6c]
	Áp suất không khí $P$ theo công thức $P=P_0\cdot\mathrm{e}^{kx}$ (mmHg), trong đó $x$ là độ cao, $P_0=760$ (mmHg) là áp suất không khí ở mức nước biển $(x=0)$, $k$ là hệ số suy giảm. Biết rằng ở độ cao $1\,000$  m thì áp suất không khí là $672{,}71$ (mmHg). Tính áp suất của không khí ở độ cao $4\,000$ m.
	\shortans{466} 
	 \loigiai{
		Ở độ cao $1\,000$ m áp suất không khí là  $672{,}71$ (mmHg).\\
		Nên ta có $672{,}71=760\mathrm{e}^{1\,000k}\Leftrightarrow\mathrm{e}^{1\,000k}=\dfrac{672{,}71}{760}\Leftrightarrow  k=\dfrac{1}{1000}\ln\dfrac{672{,}71}{760}$.\\
	Áp suất  ở độ cao $4\,000$ m là $P=760\mathrm{e}^{4\,000k}=760\mathrm{e}^{4\,000\cdot\dfrac{1}{1\,000}\ln\dfrac{672{,}71}{760}}\approx 466{,}52$ (mmHg).}
\end{ex}

% nguồn: Ok_TraLoiNgan_Lop11_Lan2
\begin{ex}%[1D6V3-5]%[ma:e9091c9b]
	Một điện thoại đang nạp pin, dung lượng pin nạp được tính theo công thức\\ $Q(t)=Q_0\cdot\left(1-\mathrm{e}^{-t\sqrt{2}}\right)$ với $t$ là khoảng thời gian tính bằng giờ  và $Q_0$ là dung lượng nạp tối đa (pin đầy). Thời gian nạp pin của điện thoại tính từ lúc cạn hết pin cho đến khi điện thoại đạt được $90\%$ dung lượng pintối đa bằng bao nhiêu? (Kết quả được làm tròn đến hàng phần trăm)
	\shortans{1,63} 
	 \loigiai{
		Theo bài ta  có\\
		$$Q_0\cdot\left(1-\mathrm{e}^{-t\sqrt{2}}\right)=0{,}9\cdot Q_0\Leftrightarrow  1-\mathrm{e}^{-t\sqrt{2}}=0{,}9\Leftrightarrow\mathrm{e}^{-t\sqrt{2}}=0{,}1\Leftrightarrow t=-\dfrac{\ln(0{,}1)}{\sqrt{2}}\approx 1{,}63.$$
	Vậy sau khoảng thời gian $t\approx 1{,}63$ giờ thì dung lượng pin của điện  thoại tính từ lúc cạn hết pin sẽ nạp được $90\%$ dung lượng pintối đa.}
\end{ex}

% nguồn: Ok_TraLoiNgan_Lop11_Lan2
\begin{ex}%[1D6V3-5]%[ma:4fbce175]
	Cho áp suất không khí $P$ (đo bằng milimet thuỷ ngân, kí hiệu là mmHg) suy giảm mũ so với độ cao $x$ (đo bằng mét), tức $P$ giảm theo công thức $P=P_0\mathrm{e}^{xi}$ trong đó $P_0=760$ mmHg là áp suất ở mực nước biển $(x=0)$, $i$ là hệ số suy giảm. Biết rằng ở độ cao $1\,000$ m thì áp suất của không khí là $672{,}71$ mmHg. Áp suất không khí ở độ cao $3\,580$ m bằng bao nhiêu? (Kết quả làm tròn đến  hàng đơn vị)
	\shortans{491} 
	 \loigiai{
		Áp dụng công thức  $P=P_0\mathrm{e}^{xi}$.\\
		Ở độ cao $1\,000$ m, ta có $P_0=760$ mmHg, $x=1\,000$ m,  $P=672{,}71$ mmHg, từ giả thiết này ta tìm được hệ số suy giảm $i$. Ta có $672{,}71=760\mathrm{e}^{1\,000\times i}\Leftrightarrow 1\,000i=\ln\dfrac{672{,}71}{760}\Leftrightarrow i=\dfrac{1}{1\,000}\ln\dfrac{672{,}71}{760}$.\\
	Khi đó ở độ cao $3\,580$ m, áp suất của  không khí là $P=760\mathrm{e}^{\tfrac{1}{1\,000}\ln\dfrac{672{,}71}{760}\times 3\,580}\approx 491{,}04$.}
\end{ex}

% nguồn: Ok_TraLoiNgan_Lop11_Lan2
\begin{ex}%[1D6V3-5]%[ma:22dd50c2]
	Gọi $I(t)$ là số ca bị nhiễm bệnh Covid-19 ở quốc gia X sau $t$ ngày khảo sát. Khi đó ta có công thức $I(t)=A\cdot\mathrm{e}^{r_0(t-1)}$ với $A$ là số ca bị nhiễm trong ngày khảo sát đầu tiên, $r_0$ là hệ số lây nhiễm. Biết rằng ngày đầu tiên khảo sát có  $500$ ca bị nhiễm bệnh và ngày thứ $10$ khảo sát có $1\,000$ ca bị nhiễm bệnh. Ngày thứ $20$ số ca nhiễm bệnh bằng bao nhiêu, biết rằng trong suốt quá trình khảo sát hệ số lây nhiễm là không đổi?
	\shortans{2,16} 
	 \loigiai{
		Theo  giả thiết ta có $I(1)=A=500$.\\
		Ngày  thứ $10$ có $1\,000$ ca nên $I(10)=A\cdot\mathrm{e}^{9r_0}\Leftrightarrow 1000=500\cdot\mathrm{e}^{9r_0}\Leftrightarrow r_0=\dfrac{\ln 2}{9}$.\\
	Vậy ngày thứ $20$ số ca nhiễm  bệnh là $I(20)=500\cdot\mathrm{e}^{\tfrac{19\ln 2}{9}}\approx 2\,160$.}
\end{ex}

% nguồn: Ok_TraLoiNgan_Lop11_Lan2
\begin{ex}%[1D6V3-5]%[ma:42ad4951]
	Gọi $N(t)$ là số phần trăm cacbon $14$ còn lại trong một bộ phận của một cây sinh trưởng từ $t$ năm trước đây thì ta có công thức $N(t)=100\cdot\left(\dfrac{1}{2}\right)^{\tfrac{t}{A}}(\%)$ với $A$ là hằng số. Biết rằng một mẫu gỗ có tuổi khoảng $3\,754$ năm thì lượng cacbon $14$  còn lại là $65\%$. Phân tích mẫu gỗ từ một công trình kiến trúc cổ, người ta thấy lượng cacbon $14$ còn lại trong mẫu gỗ là $79\%$. Mẫu gỗ được lấy từ công trình đó bao nhiêu tuổi?
	\shortans{2,06} 
	 \loigiai{
		 Theo bài ta có $65=100\cdot\left(\dfrac{1}{2}\right)^{\tfrac{3\,754}{A}}\Leftrightarrow 0{,}65=\left(\dfrac{1}{2}\right)^{\tfrac{3\,754}{A}}\Leftrightarrow\dfrac{3\,754}{A}=\log_{\tfrac{1}{2}}0{,}65\Leftrightarrow A=\dfrac{3\,754}{\log_{\tfrac{1}{2}}0{,}65}$.\\
		Do mẫu gỗ còn $79\%$ lượng Cacbon $14$ nên ta  có
		\begin{eqnarray*}
			 79=100\cdot\left(\dfrac{1}{2}\right)^{\tfrac{t}{A}}
			&\Leftrightarrow& 0{,}79=\left(\dfrac{1}{2}\right)^{\tfrac{t}{A}}\\ 
			&\Leftrightarrow&\dfrac{t}{A}=\log_{\tfrac{1}{2}}0{,}79\\ 
			 &\Leftrightarrow& t=A\cdot\log_{\tfrac{1}{2}}0{,}79=\dfrac{3\,754}{\log_{\tfrac{1}{2}}0{,}65}\cdot\log_{\tfrac{1}{2}}0{,}79\approx 2\,055.
	\end{eqnarray*}}
\end{ex}

% nguồn: Ok_TraLoiNgan_Lop11_Lan2
\begin{ex}%[1D6V3-5]%[ma:838075af]
	Áp suất không khí $P$ (đo bằng milimet thủy ngân, kí hiệu mmHg) là một đại lượng được tính theo công thức $P=P_0\mathrm{e}^{xi}$ trong đó $x$ là độ cao (đo bằng mét, so với mực nước biển), $P_0=760$ mmHg là áp suất ở mực nước biển, $i$ là hệ số suy giảm. Biết rằng, ở độ cao $1\,000$ m thì áp suất của không khí là $672{,}72$ mmHg. Áp suất của không khí ở độ cao $15$ km bằng bao nhiêu milimet  thủy ngân? (Kết quả được làm tròn đến hàng đơn vị)
	 \shortans{122}
	 \loigiai{
		Do ở độ cao $1\,000$ m, áp suất của  không khí là $672{,}72$ mmHg nên ta có\\
		 $$672{,}72=760\mathrm{e}^{1\,000i}\Leftrightarrow i=\dfrac{1}{1\,000}\ln\dfrac{672{,}72}{760}.$$
		Khi ở độ cao $15$ km tức là $15\,000$  m thì áp suất của không khí
		$$P=760\mathrm{e}^{15\,000\times\dfrac{1}{1\,000}\ln\dfrac{672{,}72}{760}}\approx  121{,}93399.$$
	 Vậy, áp suất của không khí ở độ cao $15$ km gần nhất với số $122$.}
\end{ex}


%%%%%%%%%%%%%%%% TỰ LUẬN %%%%%%%%%%%%%%%%

% nguồn: Ok_TuanLuan_Toan11_Lan1
\begin{bt}%[1D6N3-5]%[ma:31237a57]
	Số lượng của loại vi khuẩn $A$ trong một phòng thí nghiệm được tính theo công thức $s(t)=s(0) \cdot 2^t$ trong đó $s(0)$ là số lượng vi khuẩn $A$ lúc ban đầu $,s(t)$ là số lượng vi khuẩn $A$ có sau $t$ phút. Biết sau $3$ phút thì số lượng vi khuẩn $A$ là $625$ nghìn con. Hỏi sau bao lâu, kể từ lúc ban đầu, số lượng vi khuẩn $A$  là $10$ triệu con?
	 \loigiai{
		Theo đề bài ta có  \\$s(3) = 625 \Rightarrow s(0)\cdot 2^3 = s(3) \Leftrightarrow s(0) = \dfrac{s(3)}{8} = 78{,}125$ (nghìn con).\\
		Do $s(t) = 10$ triệu con $\Rightarrow s(t) = 10\,000$ nghìn con $\Rightarrow 2^t = \dfrac{s(t)}{s(0)}  = 18 \Rightarrow t = 7$ phút.
	 }
\end{bt}

% nguồn: Ok_TuanLuan_Toan11_Lan1
\begin{bt}%[1D6N3-5]%[ma:7be47263]
	Công ty FTK về mua bán xe ô tô đã qua sử dụng, sau khi khảo sát thị trường $6$ tháng đã đưa ra công thức chung về giá trị còn lại của ô tô $4$ chỗ kể từ khi đưa vào sử dụng (các loại xe $4$ chỗ không sử dụng mục đích kinh doanh) được tính $P(t)=A\cdot\left(\frac{3}{4}\right)^\frac{t}{4}$. Trong đó $A$ là giá tiền ban đầu mua xe, $t$ là số năm kể từ khi đưa  vào sử dụng.
	Tính giá trị còn lại  của xe ô tô sau $30$ tháng đưa vào sử dụng. Biết giá trị mua xe ban đầu là $920$ triệu.
	 \loigiai{
		Ta có  $A=920$ triệu; $t=2{,}5$ năm.\\
		Vậy giá trị còn lại của  xe ô tô sau $30$ tháng đưa vào sử dụng là\\ $P(2{,}5)=920\cdot\left(\frac{3}{4}\right)^\frac{2{,}5}{4}=768\,601\,304$(đồng).
	 }
\end{bt}

% nguồn: Ok_TuanLuan_Toan11_Lan2
\begin{ex}%[1D6N3-5]%[ma:3945a6ad]
	Số lượng của loại vi khuẩn $A$ trong một phòng thí nghiệm được tính theo công thức $s(t)=s(0) \cdot 2^t$ trong đó $s(0)$ là số lượng vi khuẩn $A$ lúc ban đầu $,s(t)$ là  số lượng vi khuẩn $A$ có sau $t$ phút. Biết sau $3$ phút thì số lượng vi khuẩn $A$ là $625$ nghìn con. Hỏi sau bao lâu, kể từ lúc ban đầu, số lượng vi khuẩn $A$ là $10$ triệu con?
	 \loigiai{
		Theo đề bài ta có \\$s(3) =  625 \Rightarrow s(0)\cdot 2^3 = s(3) \Leftrightarrow s(0) = \dfrac{s(3)}{8} = 78{,}125$ (nghìn con).\\
		Do $s(t) =  10$ triệu con $\Rightarrow s(t) = 10\,000$ nghìn con $\Rightarrow 2^t = \dfrac{s(t)}{s(0)} = 18 \Rightarrow t = 7$ phút.
	 }
\end{ex}

% nguồn: Ok_TuanLuan_Toan11_Lan2
\begin{ex}%[1D6N3-5]%[ma:707030a3]
	Công ty FTK về mua bán xe ô tô đã qua sử dụng, sau khi khảo sát thị trường $6$ tháng đã đưa ra công thức chung về giá trị còn lại của ô tô $4$ chỗ kể từ khi đưa vào sử dụng (các loại xe $4$ chỗ không sử dụng mục đích kinh doanh) được tính $P(t)=A\cdot\left(\frac{3}{4}\right)^\frac{t}{4}$. Trong đó $A$ là giá tiền ban đầu  mua xe, $t$ là số năm kể từ khi đưa vào sử dụng.
	Tính giá trị còn lại  của xe ô tô sau $30$ tháng đưa vào sử dụng. Biết giá trị mua xe ban đầu là $920$ triệu.
	 \loigiai{
		Ta có $A=920$ triệu; $t=2{,}5$ năm.\\ 
		Vậy giá trị còn lại của xe ô tô sau $30$ tháng đưa  vào sử dụng là\\ $P(2{,}5)=920\cdot\left(\frac{3}{4}\right)^\frac{2{,}5}{4}=768\,601\,304$(đồng).
	} 
\end{ex}

% nguồn: Ok_TuanLuan_Toan11_Lan1
\begin{bt}%[1D6H3-5]%[ma:0ba87729]
	 Trong vật lí, sự phân rã các chất phóng xạ được cho bởi công thức $m(t) = m_0 \left(\dfrac{1}{2}\right)^{\tfrac{t}{T}}$. Trong đó, $m_0$ là khối lượng chất phóng xạ ban đầu (tại thời điểm $t= 0$), $m(t)$ là khối lượng của chất phóng xạ tại thời điểm $t$ và $T$ là chu kì bán rã.\\
	Hạt nhân Poloni ($Po$) là chất phóng xạ $\alpha$ có chu kì bán rã $138$ ngày. Giả sử  lúc đầu có $100$ Poloni. Tính khối lượng Poloni còn lại sau $100$ ngày theo đơn vị gam (làm tròn kết quả đến phần nghìn).
	\loigiai{ 
		Khối lượng của  Poloni sau $100$ ngày là\\
		$m(100)= 100 \cdot  \left(\dfrac{1}{2}\right)^{\tfrac{100}{138}} = 60{,}515$ (gam).
	} 
\end{bt}

% nguồn: Ok_TuanLuan_Toan11_Lan1
\begin{bt}%[1D6H3-5]%[ma:2bb0bb43]
	 Ông Ba gửi vào ngân hàng $200$ triệu đồng theo hình thức lãi kép có kì hạn là $12$ tháng với lãi suất ngân hàng là $6\%$/năm. Giả sử qua các năm thì lãi suất không thay đổi và ông Ba không gửi thêm tiền vào mỗi năm. Hỏi sau $3$ năm ông Ba có số tiền cả gốc và lãi bằng bao nhiêu triệu đồng? (Làm tròn kết quả đến hàng đơn vị)
	\loigiai{ 
		Áp dụng công  thức lãi kép ta có $T=200\cdot\left(1+6\%\right)^3\approx 238$ (triệu đồng).
	} 
\end{bt}

% nguồn: Ok_TuanLuan_Toan11_Lan1
\begin{bt}%[1D6H3-5]%[ma:1e49a910]
	Một người gửi $100$ triệu đồng vào ngân hàng với lãi suất $8\% /$ năm. Biết rằng nếu người đó không rút tiền ra khỏi ngân hàng thì cứ sau mỗi năm số tiền lãi sẽ được nhập vào vốn để tính lãi cho năm tiếp theo. Số tiền người đó nhận sau $n$ năm sẽ được tính theo công thức $T_n=100(1+r)^n$ (triệu đồng), trong đó $r(\%)$ là lãi suất và $n$ là số năm gửi tiền.Hỏi số tiền lãi thu được của người đó  sau $10$ năm là bao nhiêu? (Kết quả trong bài được tính chính xác đến hàng phần trăm).
	\loigiai{ 
		Số tiền người đó  nhận sau $10$ năm là
		\[T_{10}=100\left(1+\dfrac{8}{100} \right)^{10}  \approx 215{,}89 \text{(triệu đồng)}.\]
		Số tiền lãi sau $10$ năm gửi tiền  xấp xỉ là $215{,}89-100=115{,}89$ (triệu đồng).
	} 
\end{bt}

% nguồn: Ok_TuanLuan_Toan11_Lan1
\begin{bt}%[1D6H3-5]%[ma:bf3283fe]
	Giả sử giá trị còn lại (tính theo triệu đồng) của một chiếc ô tô sau $t$ năm  sử dụng được mô hình hóa bằng công thức $V(t)=A\cdot (0{,}905)^t$, trong đó $A$ là giá xe (tính theo triệu đồng) lúc mới mua. Hỏi nếu theo mô hình này, sau bao nhiêu năm sử dụng thì giá trị của chiếc xe đó còn lại không quá $300$ triệu đồng? (Làm tròn kết quả đến hàng đơn vị). Biết $A=780$ (triệu đồng).
	 \loigiai{
	Theo  đề bài ta có $$V(t)\le 300\Leftrightarrow 780\cdot (0{,}905)^t\le 300\Leftrightarrow (0{,}905)^t\le\dfrac{5}{13}\Leftrightarrow t\ge\log_{0{,}095}\left(\dfrac{5}{13} \right)\approx 9{,}6.$$
	Vậy sau khoảng $10$ năm sử dụng thì giá trị của chiếc xe đó còn lại không quá $300$ triệu  đồng.
	 }
\end{bt}

% nguồn: Ok_TuanLuan_Toan11_Lan1
\begin{bt}%[1D6H3-5]%[ma:b32eef2b]
	Một người gửi $100$ triệu đồng vào ngân hàng với lãi suất $8\%/$năm. Biết rằng nếu người đó không rút tiền ra khỏi ngân hàng thì cứ sau mỗi năm số tiền lãi sẽ được nhập vào vốn để tính lãi cho năm tiếp theo.  Số tiền người đó nhạn sau $n$ năm sẽ được tính theo công thức $T_n=100\cdot (1+r)^{n}$ (triệu đồng), trong đó $r\,(\%)$ là lãi suất và $n$ là số năm gửi tiền.\\
	Hỏi số tiền lãi thu được của người đó sau $10$ năm là bao  nhiêu?
	 \loigiai{
	 Số tiền người đó nhận sau $10$ năm là $T_{10}=100\cdot\left(1+\dfrac{8}{100} \right)^{10}\approx 215{,}89$ (triệu đồng).\\
	Số tiền lãi sau $10$ năm xấp xỉ là $215{,}89-199=115{,}89$ (triệu  đồng).
	 }
\end{bt}

% nguồn: Ok_TuanLuan_Toan11_Lan1
\begin{bt}%[1D6H3-5]%[ma:f7b2bda2]
	Nếu $D_0$ là chênh lệch nhiệt độ ban đầu giữa một vật $M$ và các vật xung  quanh, và nếu các vật xung quanh có nhiệt độ $T_S$, thì nhiệt độ của vật $M$ tại thời điểm $t$ được mô hình hóa bởi hàm số: $T(t)=T_S+D_0 \cdot\mathrm{e}^{-kt}\,(1)$ (trong đó $k$ là hằng số dương phụ thuộc vào vật $M$).
	Một con gà tây nướng được lấy từ lò nướng khi nhiệt độ của nó đã  đạt đến $195^\circ F$ và được đặt trên một bàn trong một căn phòng có nhiệt độ là $65^\circ F$.
	Nếu nhiệt độ của gà tây là $150^\circ F$ sau nửa giờ, nhiệt độ của nó sau $60$  phút là bao nhiêu?
	 \loigiai{
		Ta  có $T_S=65^\circ F$ và độ chênh lệch nhiệt độ là $D_0=195^\circ F-65^\circ F=130^\circ F$.\\
		Sau nửa giờ $\left(t=0{,}5\right)$  thì nhiệt độ của gà là $T=150^\circ F$.\\
		Áp dụng  công thức $(1)$ ta được $150=65+130\cdot\mathrm{e}^{-0{,}5k}\Leftrightarrow\mathrm{e}^{-k}=\left(\frac{17}{26}\right)^2$.\\
		 Vậy $T(t)=65+130\cdot\left(\frac{17}{26} \right)^{2t}$.\\
		 Suy ra nhiệt độ của gà sau $60$ phút ($t=1$ giờ) là $65+130\cdot\left(\frac{17}{26}\right)^{2\cdot1} \approx121^\circ F$.
	} 
\end{bt}

% nguồn: Ok_TuanLuan_Toan11_Lan2
\begin{ex}%[1D6H3-5]%[ma:bdf5415e]
	Trong vật lí, sự phân rã các chất phóng xạ được cho bởi công thức $m(t) = m_0 \left(\dfrac{1}{2}\right)^{\tfrac{t}{T}}$. Trong đó, $m_0$ là khối lượng chất phóng xạ ban đầu (tại thời điểm $t= 0$), $m(t)$ là khối lượng của chất phóng xạ tại thời điểm $t$ và $T$ là chu kì bán rã.\\ 
	Hạt nhân Poloni ($Po$) là chất phóng xạ $\alpha$ có chu kì bán rã $138$ ngày. Giả sử lúc đầu có $100$ Poloni. Tính khối lượng Poloni còn lại sau $100$ ngày theo đơn vị gam (làm  tròn kết quả đến phần nghìn).
	 \loigiai{
		Khối lượng của Poloni sau $100$ ngày là\\ 
		$m(100)= 100  \cdot \left(\dfrac{1}{2}\right)^{\tfrac{100}{138}} = 60{,}515$ (gam).
	 }
\end{ex}

% nguồn: Ok_TuanLuan_Toan11_Lan2
\begin{ex}%[1D6H3-5]%[ma:870414d6]
	Ông Ba gửi vào ngân hàng $200$ triệu đồng theo hình thức lãi kép có kì hạn là  $12$ tháng với lãi suất ngân hàng là $6\%$/năm. Giả sử qua các năm thì lãi suất không thay đổi và ông Ba không gửi thêm tiền vào mỗi năm. Hỏi sau $3$ năm ông Ba có số tiền cả gốc và lãi bằng bao nhiêu triệu đồng? (Làm tròn kết quả đến hàng đơn vị)
	\loigiai{ 
		Áp dụng công  thức lãi kép ta có $T=200\cdot\left(1+6\%\right)^3\approx 238$ (triệu đồng).
	} 
\end{ex}

% nguồn: Ok_TuanLuan_Toan11_Lan2
\begin{ex}%[1D6H3-5]%[ma:b73ed348]
	Một người gửi $100$ triệu đồng vào ngân hàng với lãi suất $8\% /$ năm. Biết rằng nếu người đó không rút tiền ra khỏi ngân hàng thì cứ sau mỗi năm số tiền lãi sẽ được nhập vào vốn để tính lãi cho năm tiếp theo. Số tiền người đó nhận sau $n$ năm sẽ được tính theo công thức $T_n=100(1+r)^n$ (triệu đồng), trong đó $r(\%)$ là lãi suất và $n$ là số năm gửi tiền.Hỏi số tiền lãi thu được của người đó sau $10$ năm là bao nhiêu? (Kết quả trong bài được tính chính xác đến hàng phần trăm). 
	 \loigiai{
		 Số tiền người đó nhận sau $10$ năm là
		\[T_{10}=100\left(1+\dfrac{8}{100}  \right)^{10} \approx 215{,}89 \text{(triệu đồng)}.\]
		Số tiền lãi sau $10$ năm gửi tiền  xấp xỉ là $215{,}89-100=115{,}89$ (triệu đồng).
	 }
\end{ex}

% nguồn: Ok_TuanLuan_Toan11_Lan2
\begin{ex}%[1D6H3-5]%[ma:fd55a41a]
	Giả sử giá trị còn lại (tính  theo triệu đồng) của một chiếc ô tô sau $t$ năm sử dụng được mô hình hóa bằng công thức $V(t)=A\cdot (0{,}905)^t$, trong đó $A$ là giá xe (tính theo triệu đồng) lúc mới mua. Hỏi nếu theo mô hình này, sau bao nhiêu năm sử dụng thì giá trị của chiếc xe đó còn lại không quá $300$ triệu đồng? (Làm tròn kết quả đến hàng đơn vị). Biết $A=780$ (triệu đồng).
	 \loigiai{
	Theo  đề bài ta có $$V(t)\le 300\Leftrightarrow 780\cdot (0{,}905)^t\le 300\Leftrightarrow (0{,}905)^t\le\dfrac{5}{13}\Leftrightarrow t\ge\log_{0{,}095}\left(\dfrac{5}{13} \right)\approx 9{,}6.$$
	Vậy sau khoảng $10$ năm sử dụng thì  giá trị của chiếc xe đó còn lại không quá $300$ triệu đồng.
	} 
\end{ex}

% nguồn: Ok_TuanLuan_Toan11_Lan2
\begin{ex}%[1D6H3-5]%[ma:2c6e1b12]
	Một người gửi $100$ triệu đồng vào ngân hàng với lãi suất $8\%/$năm. Biết rằng nếu người đó không rút tiền  ra khỏi ngân hàng thì cứ sau mỗi năm số tiền lãi sẽ được nhập vào vốn để tính lãi cho năm tiếp theo. Số tiền người đó nhạn sau $n$ năm sẽ được tính theo công thức $T_n=100\cdot (1+r)^{n}$ (triệu đồng), trong đó $r\,(\%)$ là lãi suất và $n$ là số năm gửi tiền.\\
	Hỏi  số tiền lãi thu được của người đó sau $10$ năm là bao nhiêu?
	 \loigiai{
	Số tiền người đó nhận sau $10$ năm  là $T_{10}=100\cdot\left(1+\dfrac{8}{100} \right)^{10}\approx 215{,}89$ (triệu đồng).\\
	Số  tiền lãi sau $10$ năm xấp xỉ là $215{,}89-199=115{,}89$ (triệu đồng).
	} 
\end{ex}

% nguồn: Ok_TuanLuan_Toan11_Lan2
\begin{ex}%[1D6H3-5]%[ma:1dc2783d]
	Nếu $D_0$ là  chênh lệch nhiệt độ ban đầu giữa một vật $M$ và các vật xung quanh, và nếu các vật xung quanh có nhiệt độ $T_S$, thì nhiệt độ của vật $M$ tại thời điểm $t$ được mô hình hóa bởi hàm số: $T(t)=T_S+D_0 \cdot\mathrm{e}^{-kt}\,(1)$ (trong đó $k$ là hằng số dương phụ thuộc vào vật $M$).
	Một con gà tây nướng  được lấy từ lò nướng khi nhiệt độ của nó đã đạt đến $195^\circ F$ và được đặt trên một bàn trong một căn phòng có nhiệt độ là $65^\circ F$.
	Nếu nhiệt  độ của gà tây là $150^\circ F$ sau nửa giờ, nhiệt độ của nó sau $60$ phút là bao nhiêu?
	 \loigiai{
		Ta có $T_S=65^\circ F$ và độ chênh lệch nhiệt độ là $D_0=195^\circ F-65^\circ F=130^\circ F$.\\ 
		Sau nửa giờ $\left(t=0{,}5\right)$ thì nhiệt độ của gà là $T=150^\circ F$.\\ 
		Áp dụng công thức $(1)$  ta được $150=65+130\cdot\mathrm{e}^{-0{,}5k}\Leftrightarrow\mathrm{e}^{-k}=\left(\frac{17}{26}\right)^2$.\\
		 Vậy $T(t)=65+130\cdot\left(\frac{17}{26} \right)^{2t}$.\\
		Suy ra nhiệt  độ của gà sau $60$ phút ($t=1$ giờ) là $65+130\cdot\left(\frac{17}{26}\right)^{2\cdot1} \approx121^\circ F$.
	 }
\end{ex}

% nguồn: Ok_TuanLuan_Toan11_Lan2
\begin{ex}%[1D6H3-5]%[ma:b0dc7b60]
	Lúc đầu trong ao có một số con ếch. Người ta ghi nhận số lượng ếch trong $5$  năm đầu như hình vẽ. Giả sử số lượng ếch tăng theo hàm số $n(t)=C\cdot a^t$.
	\begin{enumerate}
		\item Tính số lượng ếch lúc ban đầu. 
		\item Tìm hàm số biểu diễn số lượng ếch sau $t$ năm kể từ khi chúng xuất hiện  trong ao.
		\item Dự đoán số  lượng ếch sau $20$ năm.
	\end{enumerate}
	\begin{center}
		\begin{tikzpicture}[line cap=round,line join=round,font=\footnotesize,scale=0.8]
				\begin{axis}[
						width=12cm,height=8cm, 
						 xmin=0,xmax=6.5,
						 ymin=0,ymax=650,
						axis lines=middle, 
						 xlabel style={at={(axis description cs:0.95,0.02)},anchor=north east},
						 ylabel style={at={(axis description cs:0.02,0.95)},anchor=south west},
						xtick={0,1,2,3,4,5,6}, 
						ytick={0,100,200,300,400,500,600}, 
						 xticklabel style={/pgf/number format/fixed},
						yticklabel style={ 
						/pgf/number format/fixed, 
						/pgf/number format/use  comma=false
						}] 
					\addplot[blue] {100*(1.4)^x}; 
					\addplot[only marks,mark=*,mark size=2pt,black] coordinates {(2,196)}; 
					\addplot[black,dashed]  coordinates {(2,0)
						(2,196)}; 
					 \addplot[black,dashed] coordinates {(0,196)
						 (2,196)};
					\node[above left] at  (axis cs:2,196) {\small (2; 196)};
				\end{axis}
				 \draw (10.5,0) node[below]{\bf Năm};
				\draw (0,6.5) node[above]{\bf  Con};
			\end{tikzpicture}
	\end{center}
	 \loigiai{
		\begin{enumerate}
			 \item Số lượng ếch ban đầu là $100$ con tại $t=0$.
			\item Đồ thị hàm  số đi qua điểm $(0;100)$ và $(2;196)$ nên ta có
			$$\heva{&C\cdot a^0=100\\& C\cdot a^2=196} \Leftrightarrow  \heva{&C=100\\ &a^2=1{,}96} \Leftrightarrow\heva{&C=100\\& a=1{,}4}\Rightarrow n(t)=100\cdot 1{,}4^t.$$
			Vậy hàm số biểu diễn số lượng ếch sau $t$ năm kể từ khi chúng xuất hiện  trong ao là
			 $$H(t)=100+100\cdot 1{,}4^t.$$
			\item Lượng ếch sau $20$ năm:  $H(20)=100+100\cdot 1{,}4^{20} \approx 83\,768$ con.
		\end{enumerate}
	 }
\end{ex}

% nguồn: Ok_TuanLuan_Toan11_Lan2
\begin{ex}%[1D6H3-5]%[ma:4009ef94]
	Một người gửi $100$ triệu đồng vào một ngân hàng với lãi suất $6\%$/năm. Biết rằng nếu không rút tiền ra khỏi ngân hàng thì cứ sau mỗi năm số tiền lãi sẽ  được nhập vào gốc để tính lãi cho năm tiếp theo. Hỏi sau ít nhất bao nhiêu năm người đó nhận được số tiền nhiều hơn $300$ triệu đồng bao gồm cả gốc lẫn lãi?
	 \loigiai{
		Theo công thức tính lãi suất kép, ta có vốn tích luỹ sau $n$ năm là  $P_n=P(1+r)^n$ với $P$ là vốn ban đầu, $r$ là lãi suất.\\
		Theo đề bài ta có phương  trình
		$$300=100\left(1+\dfrac{6}{100}\right)^n\Leftrightarrow n=\log_{1{,}06}3\approx 19.$$ 
		Vậy sau ít nhất $19$ năm thì người đó nhận được số tiền nhiều hơn $300$ triệu đồng bao gồm cả gốc lẫn lãi. 
	} 
\end{ex}

% nguồn: Ok_TuanLuan_Toan11_Lan2
\begin{ex}%[1D6H3-5]%[ma:26ee0e7e]
	Một người gửi $50$ triệu đồng vào ngân hàng theo thể thức  lãi kép với lãi suất $5{,}5\%$/năm, kì hạn $1$ năm. Hỏi sau $4$ năm, người đó rút cả vốn lẫn lãi được bao nhiêu tiền?
	\loigiai{ 
		Số tiền cả vốn lẫn lãi sau $n$ năm  được tính theo công thức $A_n=A(1+r)^n$.\\
		Thay số  với $A=50$, $r=5{,}5\%$, $n=4$
		ta được số tiền là $$A_4=50\left(1+\dfrac{5{,}5}{100}\right)^4\approx 61{,}94\  \text{triệu đồng}.$$
	 }
\end{ex}

% nguồn: Ok_TuanLuan_Toan11_Lan2
\begin{ex}%[1D6H3-5]%[ma:16522933]
	Ông A gửi tiền  tiết kiệm với lãi suất $8{,}1\%$/ năm và lãi suất hằng năm được nhập vào vốn. Hỏi sau bao nhiêu năm ông A được số tiền gấp đôi số tiền ban đầu?
	 \loigiai{
		Gọi số tiền ban đầu ông A gửi tiết kiệm  là $B$.\\
		Theo công thức lãi kép  ta có số tiền sau $n$ năm là $B(1+0{,}081)^n$.\\
		Để số tiền tăng  gấp đôi thì $n$ phải thỏa mãn phương trình
		$$B(1+0{,}081)^n=2B \Leftrightarrow (1{,}081)^n=2\Leftrightarrow n=\log_{1{,}081}  2\approx 8{,}899.$$
		Như vậy sau  $9$ năm ông A sẽ thu được số tiền gấp đôi số tiền ban đầu.
	} 
\end{ex}

% nguồn: Ok_TuanLuan_Toan11_Lan2
\begin{ex}%[1D6H3-5]%[ma:8b1a7a8f]
	Lãi suất gửi tiền tiết kiệm  của các ngân hàng trong thời gian qua liên tục thay đổi. Bác Mạnh gửi vào một ngân hàng số tiền $5$ triệu đồng với lãi suất $0{,}7\%$/tháng. Sau sáu tháng gửi tiền, lãi suất tăng lên $0{,}9\%$/tháng. Đến tháng thứ $10$ sau khi gửi tiền, lãi suất giảm xuống $0{,}6\%$/tháng và giữ ổn định. Biết rằng nếu bác Mạnh không rút tiền ra khỏi ngân hàng thì cứ sau mỗi tháng, số tiền lãi sẽ được nhập vào vốn ban đầu. Sau một năm gửi tiền, bác Mạnh rút được số tiền là bao nhiêu?
	\loigiai{ 
		Số tiền bác Mạnh thu được $5(1+0{,}007)^6(1+0{,}009)^3(1+0{,}006)^3=5{,}452733453$ triệu  đồng.
	} 
\end{ex}

% nguồn: Ok_TuanLuan_Toan11_Lan2
\begin{ex}%[1D6H3-5]%[ma:46a2aebc]
	Một người gửi $200$  triệu đồng vào ngân hàng với lãi suất $0{,}2\%$/tháng. Biết rằng nếu không rút tiền khỏi ngân hàng thì cứ sau mỗi tháng số tiền lãi sẽ được nhập vào vốn ban đầu để tính lãi cho tháng tiếp theo. Hỏi sau đúng $10$ tháng người đó được lĩnh bao nhiêu tiền?
	 \loigiai{
		Theo công thức lãi kép ta có số tiền cả lãi và vốn  sau $10$ tháng là
		$$T=200(1+0{,}2\%)^{10} \approx 204{,}036\  \text{triệu đồng}.$$
	 }
\end{ex}

% nguồn: Ok_TuanLuan_Toan11_Lan2
\begin{ex}%[1D6H3-5]%[ma:ec1f1d35]
	Cường độ ánh sáng $I$ dưới mặt biển giảm dần theo độ sâu theo công thức $I=I_0\cdot a^d$ trong đó $I_0$ là cường độ ánh sáng tại mặt nước biển, $a$ là hằng số ($a>0$) và $d$ là độ sâu tính bằng mét tính từ mặt  nước biển.
	\begin{enumerate}
		\item  Có thể khẳng định rằng $0<a<1$ không? Vì sao?
		\item Biết  rằng cường độ ánh sáng tại độ sâu $1$ m bằng $0{,}95I_0$. Tìm $a$.
		\item Tại độ sâu $20$ m thì cường độ  ánh sáng bằng bao nhiêu phần trăm so với $I_0$.
	\end{enumerate}
	\loigiai{ 
		\begin{enumerate}
			\item Do cường độ ánh sáng $I$ dưới mặt biển giảm dần theo độ sâu nên hàm số $I=I_0 \cdot a^d$ nghịch biến nên  $0<a<1$.
			\item Ta có $0{,}95I_0=I_0\cdot a^1 \Rightarrow a=0{,}95$. 
			\item Ta có $I=I_0\cdot a^d=I_0\cdot 0{,}95^{20}\approx 0{,}36I_0$ nên tại độ sâu $20$\,m thì cường độ ánh sáng bằng $36\%$  so với $I_0$.
		\end{enumerate}
	} 
\end{ex}

% nguồn: Ok_TuanLuan_Toan11_Lan2
\begin{ex}%[1D6H3-5]%[ma:dfb78186]
	Một người gửi số tiền $300$ triệu đồng vào một ngân hàng  với lãi suất kép $6\%$ một năm. Biết rằng nếu không rút tiền ra khỏi ngân hàng thì cứ sau mỗi năm, số tiền lãi suất sẽ được nhập vào vốn ban đầu. Hỏi sau $3$ năm không rút tiền gốc và lãi, số tiền trong ngân hàng của người đó xấp xỉ bao nhiêu triệu đồng?
	\loigiai{ 
		Áp dụng công thức tính  lãi suất theo hình thức lãi kép $P=A(1+r)^n$.\\
		Trong đó:  $P$ là số tiền gồm vốn lẫn lãi tại thời điểm $n$ tính từ thời điểm gửi; $A$ là số tiền gửi vào ban đầu và $r$ là lãi suất.\\
		Theo bài toán ta có $A=300\,000\,000$, $n=3$, $r=6\%$.\\ 
		Vậy số tiền là $P=300\,000\,000\left(1+\dfrac{6}{100}\right)^3\approx 357$ triệu đồng. 
	 }
\end{ex}

% nguồn: Ok_TuanLuan_Toan11_Lan2
\begin{ex}%[1D6H3-5]%[ma:24f20e01]
	Một người gửi $100$ triệu đồng vào ngân hàng với lãi suất $0{,}4\%$/tháng. Biết rằng nếu không rút tiền ta khỏi ngân hàng thì cứ sau mỗi tháng, số tiền lãi sẽ  được lập vào vốn ban đầu để tính lãi cho tháng tiếp theo. Hỏi sau $6$ tháng, người đó được bao nhiêu tiền nếu trong khoảng thời gian này người đó không rút tiền ra và lãi xuất không thay đổi?
	 \loigiai{
		Áp dụng công thức lãi kép ta có sau đúng $6$ tháng, người đó lĩnh được số tiền 
		$$100\,000\,000\left(1+\dfrac{0{,}4}{100}\right)^6\approx 102\,424\,128\  \text{đồng}.$$
	 }
\end{ex}

% nguồn: Ok_TuanLuan_Toan11_Lan2
\begin{ex}%[1D6H3-5]%[ma:a2367c99]
	Ông A gửi $200$  triệu vào một ngân hàng theo hình thức lãi kép, với lãi suất là $6{,}5\%$ một năm và lãi suất không đổi trong suốt thời gian gửi. Sau $6$ năm, số tiền lãi của ông bằng bao nhiêu?
	 \loigiai{
		Số tiền cả gốc  lẫn lãi là $200(1+6{,}5\%)^6\approx 292$ triệu.\\
		 Vậy số tiền lãi là $292-200=92$ triệu.
	 }
\end{ex}

% nguồn: Ok_TuanLuan_Toan11_Lan2
\begin{ex}%[1D6H3-5]%[ma:26b9ab88]
	Một người gửi $100$ triệu đồng vào một ngân hàng với lãi suất $0{,}7\%$/tháng. Biết rằng nếu không rút tiền ra khỏi ngân hàng thì cứ sau mỗi tháng, số tiền lãi sẽ được nhập vào vốn ban đầu để tính lãi cho tháng tiếp  theo. Hỏi sau đúng $5$ tháng, người đó được lĩnh bao nhiêu tiền?
	\loigiai{ 
		Sau $5$ tháng, người đó  được lĩnh số tiền là
		$$S=100\,000\,000(1+0{,}7\%)^5\approx 103\,549\,000\ \text{đồng}.$$ 
	} 
\end{ex}

% nguồn: Ok_TuanLuan_Toan11_Lan2
\begin{ex}%[1D6H3-5]%[ma:da564a2a]
	Trong một nghiên cứu, một nhóm học sinh được xem cùng một danh sách các loài động vật và được kiểm tra lại xem họ còn nhớ bao nhiêu phần  trăm danh sách đó sau mỗi tháng. Giả sử sau $t$ tháng, khả năng nhớ trung bình của nhóm học sinh đó được tính theo công thức $M(t)=75-20\ln(t+1), 0\le t\le 12$ (đơn vị $\%$). Hãy tính khả năng nhớ trung bình của nhóm học sinh đó sau $6$ tháng.
	\loigiai{ 
		Áp dụng công thức $M(t)=75-20\ln (t+1)$, ta có \[M(6)=75-20\ln(6+1)=75-20\ln 7\approx 60{,}39.\] 
		Vậy khả năng nhớ trung bình của nhóm học sinh sau $6$ tháng là khoảng $60{,}39\%$. 
	} 
\end{ex}

% nguồn: Ok_TuanLuan_Toan11_Lan1
\begin{bt}%[1D6V3-5]%[ma:93143b24]
	Ông Đại mới  xin được việc làm nên gửi tiết kiệm vào ngân hàng với hình thức cứ mỗi đầu tháng đóng vào 5 triệu đồng với lãi suất $0{,}33\%/$ tháng. Tính số tiền mà ông Đại thu được từ ngân hàng sau 5 năm.
	 \loigiai{
		Với $a$ là số tiền ông Đại đóng vào  hằng tháng;\\
		$r\%$ lãi suất ông Đại gửi  tiết kiệm hằng tháng.\\
		Gọi $P_n$ là số tiền mà ông Đại thu được sau  $n$ tháng $(n\ge 1)$.\\
		Suy ra 
		\begin{itemize}
			\item  $P_1=a\cdot(1+r\%)$.
			\item $P_2=\left(P_1+a\right)(1+r\%)=a\cdot(1+r\%)^2+a\cdot(1+r\%)$. 
			\item $P_3=\left(P_2+a\right)(1+r\%)=a\cdot(1+r\%)^3+a\cdot(1+r\%)^2+a\cdot(1+r\%)$. 
			\item  \dotfill
			 \item $P_n=\left(P_{n-1}+a\right)(1+r\%)=a\cdot(1+r\%)^n+a\cdot(1+r\%)^{n-1}+\cdots+a\cdot(1+r\%)$.
		\end{itemize}
		Xét cấp số nhân có số hạng đầu là $u_1=a\cdot(1+r\%)$ và công bội $q=1+r\%$ thì \begin{align*}
		P_n=u_1+u_2+\cdots+u_n=u_1\cdot \dfrac{1-q^n}{1-q}. 
	\end{align*}
	Vậy số tiền ông Đại nhận được từ ngân hàng sau $5$ năm  là
	\begin{align*}
		 P_{60}=u_1\cdot \dfrac{1-q^{60}}{1-q}=5\cdot\left(1{,}0033\right)\cdot\dfrac{1-\left(1{,}0033\right)^{60}}{0{,}0033} \approx 332\text{ (triệu đồng)}.
	\end{align*}
 }
\end{bt}

% nguồn: Ok_TuanLuan_Toan11_Lan1
\begin{bt}%[1D6V3-5]%[ma:578109d6]
	Ông Bình vậy vốn ngân hàng với số tiền $100\,000\,000$ đồng. Ông dự định sau đúng $5$ năm thì trả hết nợ theo hình thức sau đúng một tháng kể từ ngày vậy, ông bắt đầu hoàn nợ, hai lần hoàn nợ liên tiếp cách nhau đúng một tháng, số tiền hoàn nợ ở mỗi lần là như nhau. Hỏi theo cách đó, số tiền $a$ mà ông sẽ phải trả cho ngân hàng trong mỗi lần hoàn nợ là bao nhiêu? Biết lãi suất hàng tháng là $1{,}2\%$ và không thay đổi trong thời  gian ông hoàn nợ.
	\loigiai{ 
		Gọi $m$, $r$, $T_n$, $a$ lần  lượt là số tiền vậy ngân hàng, lãi suất hàng tháng, tổng số tiền vậy còn lại sau $n$ tháng, số tiền trả đều đặn mỗi tháng.
		\begin{itemize}
			\item Sau khi hết tháng thứ  nhất $(n=1)$ thì còn lại $T_1=m(r+1)-a$.
			\item  Sau khi hết tháng thứ hai $(n=2)$ thì còn lại
			\begin{eqnarray*}
				 T_2&=&\left[m(r+1)-a\right](r+1)-a\\
				 &=&m(r+1)^2-a(r+1)-a\\
				 &=&m(r+1)^2-a(r+2)\\
				 &=&m(r+1)^2-\dfrac{a}{r} \left[(r+1)^2-1\right].
			\end{eqnarray*}
			\item  Sau khi hết tháng thứ ba $(n=3)$ thì còn
			\begin{eqnarray*}
				T_3&=&\left[m(r+1)^2-\dfrac{a}{r} \left[(r+1)^2-1\right]\right](r+1)-a\\ 
				&=&m(r+1)^3-\dfrac{a}{r} \left[(r+1)^3-1\right]. 
			\end{eqnarray*}
			 \item \dotfill
			\item Sau khi hết tháng thứ $n$ thì  còn lại
			\begin{align*}
				T_n=m(r+1)^n-\dfrac{a}{r} \left[(r+1)^n-1\right]. 
			\end{align*}
		\end{itemize}
		Áp dụng công thức  trên, ta có
		\begin{eqnarray*}
			T_n=0\Leftrightarrow  a=\dfrac{m(r+1)^n r}{(r+1)^n-1}=\dfrac{12\cdot 10^5 \Bigl(\dfrac{1{,}2}{100}+1\Bigr)^{60}}{\Bigl(\dfrac{1{,}2}{100}+1\Bigr)^{60}-1}.
		\end{eqnarray*}
	 }
\end{bt}

% nguồn: Ok_TuanLuan_Toan11_Lan1
\begin{bt}%[1D6V3-5]%[ma:a56a086e]
	Một người gửi $100$ triệu  đồng vào ngân hàng với lãi suất $0{,}4\% /$ tháng. Biết rằng nếu không rút tiền ta khỏi ngân hàng thì cứ sau mỗi tháng, số tiền lãi sẽ được lập vào vốn ban đầu để tính lãi cho tháng tiếp theo. Hỏi sau $6$ tháng, người đó được lĩnh số tiền bao nhiêu, nếu trong khoảng thời gian này người đó không rút tiền ra và lãi xuất không thay đổi?
	 \loigiai{
		Áp dụng công thức lãi kép ta có sau đúng $6$  tháng, người đó lĩnh được số tiền
		\begin{align*}
			 A_n=A_0 (1+r)^n=100\,000\,000\left(1+\dfrac{0{,}4}{100} \right)^6=102\,424\,128\text{ (đồng)}.
		\end{align*}
	} 
\end{bt}

% nguồn: Ok_TuanLuan_Toan11_Lan1
\begin{bt}%[1D6V3-5]%[ma:2d5dd062]
	Thực hiện một mẻ nuôi cấy vi khuẩn với $1000$ vi khuẩn ban  đầu, nhà sinh học phát hiện ra số lượng vi khuẩn tăng thêm $25 \%$ sau hai ngày.
	\begin{enumerate}
		\item Công thức $ P(t)= P_0 \cdot a^t$ cho phép tính số lượng vi khuẩn của mẻ nuôi cấy sau $t$ ngày kể từ thời điểm ban đầu. Xác định các tham số $P_0$ và $a$, $(a>0)$. Làm tròn đến hàng phần trăm. Lấy kết quả đã làm tròn ở ý a) để làm ý b) và ý c). 
		\item Sau $5$ ngày thì số lượng vi  khuẩn bằng bao nhiêu? Làm tròn kết quả đến hàng trăm.
		\item Sau bao nhiêu ngày thì số lượng vi khuẩn bằng gấp đôi số lượng ban đầu? Làm tròn kết quả đến hàng phần  chục.
	\end{enumerate}
	\loigiai{ 
		\begin{enumerate}
			 \item Ban đầu có $1000$ vi khuẩn nên $P_0=1000$.\\
			Sau $2$ ngày, số lượng vi khuẩn là $P(2)=125 \%\cdot P_0=125 \% \cdot  1000=1250$.\\
			Ta có $P(2)=P_0 \cdot a^2 \Leftrightarrow 1250=1000 \cdot a^2 \Leftrightarrow a^2=1{,}25 \Leftrightarrow a \approx 1{,}12$. 
			\item Số lượng vi khuẩn sau $5$ ngày là $P(5)=P_0 \cdot a^5=1000 \cdot 1{,}12^2 \approx 1800$ (vi  khuẩn).
			\item Với  $P(t)=2P_0$, ta có\\
			$P(t)=P_0 \cdot a^t \Leftrightarrow 2 P_0=P_0 \cdot 1{,}12^t \Leftrightarrow 1{,}12^t=2 \Leftrightarrow  t=\log _{1,12} 2 \approx 6{,}1$ (ngày).\\
			Vậy sau  $6{,}1$ ngày thì số lượng vi khuẩn vượt gấp đôi số lượng ban đầu.
		\end{enumerate}
	} 
\end{bt}

% nguồn: Ok_TuanLuan_Toan11_Lan1
\begin{bt}%[1D6V3-5]%[ma:8b4b9ce3]
	Lãi suất cho vay tại PVcomBank trong tháng $5/2022$ rất ưu đãi, ở mức $5\%/$năm, được áp dụng trong $6$ tháng đầu, từ tháng thứ $7$ trở đi ấn định mức lãi $12\%/$năm. Tại ngân hàng này, thời hạn cho vậy mua nhà tối đa là $20$ năm, mức vậy tối đa $85\%$ giá trị  tài sản đảm bảo. Một người có khả năng trả cố định hằng tháng là $15$ triệu. Giả sử người đó có thể mượn người thân $15\%$ giá trị căn nhà, nếu được sử dụng gói vậy ở trên với thời hạn tối đa và mức vậy tối đa thì có thể mua được căn nhà có giá trị tối đa khoảng?
	\loigiai{ 
		Gọi $A$ là số tiền tối đa người này có thể vậy, $A_i$ là  số tiền nợ sau tháng thứ $i$.
		\begin{itemize}
			\item $r_1=\dfrac{5\%}{12}$ là lãi suất$/1$  tháng, trong $6$ tháng đầu.
			\item $r_2=\dfrac{12\%}{12}=1\%$ là lãi suất$/1$ tháng, từ  tháng thứ $7$ trở đi.
			\item Sau $1$ tháng, số tiền gốc và lãi là $A(1+r)$, người đó trả $15$ triệu nên còn nợ \begin{align*}
			A_1=A(1+r)-15. 
		\end{align*}
		 \item Sau tháng thứ $2$
		\begin{align*}
			A_2=A_1  (1+r_1)-15=\left(A(1+r_1)-15\right)(1+r_1)-15=A(1+r_1)^2-\dfrac{15}{r_1} \left[((1+r_1))^2-1\right].
		\end{align*}
		\item Sau tháng thứ  $3A_3=A(1+r_1)^3-\dfrac{15}{r_1} \left[(1+r_1)^3-1\right]$.
		 \item \dotfill.
		\item Sau tháng  thứ $6A_6=A(1+r_1)^6-\dfrac{15}{r_1} \left[(1+r_1)^6-1\right]$.
		\item Sau tháng thứ  $7A_7=A_6 (1+r_2)-15$.
		 \item Sau tháng thứ $8A_8=A_6 (1+r_2)^2-\dfrac{15}{r_2} \left[(1+r_2)^2-1\right]$.
		\item \dotfill 
		 \item Sau tháng thứ $240A_{240}=A_6 (1+r_2)^{234}-\dfrac{15}{r_2} \left[(1+r_2)^{234}-1\right]$.
	\end{itemize}
	Vì phải  trả hết nợ sau $20$ năm nên
	\begin{eqnarray*}
		A_{240}=0&\Leftrightarrow& A_6=\dfrac{15\left[(1+r_2)^{234}-1\right]}{\left(1+r_2 \right)^{234} r_2} \approx 1\,353{,}819328\\ 
		 &\Rightarrow& A=\dfrac{A_6+\dfrac{15}{r_1} \left[(1+r_1)^6-1\right]}{\left(1+r_1 \right)^6} \approx 1\,409{,}163992.
	\end{eqnarray*}
	Vậy người này có thể mua được căn nhà có giá trị tối đa là \begin{align*}
	\dfrac{A}{85\%} \approx 1\,657{,}83999\text{ (triệu đồng)} \approx 1{,}65784\text{ (tỷ đồng)}. 
	\end{align*}
 }
\end{bt}

% nguồn: Ok_TuanLuan_Toan11_Lan1
\begin{bt}%[1D6V3-5]%[ma:14be810c]
	Một học sinh A khi đủ 18 tuổi được cha mẹ cho $200\, 000\, 000$ VNĐ. Số tiền này được  bảo quản trong ngân hàng MSB với kì hạn thanh toán $1$ năm và học sinh A chỉ nhận được số tiền này khi học xong $4$ năm đại học. Biết rằng khi đủ $22$ tuổi, số tiền mà học sinh A được nhận sẽ là $243\,101\, 250$ VNĐ. Vậy lãi suất kì hạn một năm của ngân hàng MSB là bao nhiêu?
	\loigiai{ 
		Gọi lãi suất kỳ hạn một năm của ngân hàng MSB là  $r$. \\
		Áp dụng công thức lãi suất kép $P=a\left(1+r\right)^n$, ta có 
		\begin{eqnarray*}
			243\,101\,250=200\,000\,000(1+r)^4 &\Leftrightarrow&  (1+r)^4=\dfrac{243\,101\,250}{200\,000\,000}\\
			&\Leftrightarrow& 1+r=\sqrt[{4}]{\dfrac{243\,101\,250}{200\,000\,000}}\\ &\Leftrightarrow& r=\sqrt[{4}]{\dfrac{243\,101\,250}{200\,000\,000}}-1\\ 
			&\Leftrightarrow& r\approx  0{,}05.
		\end{eqnarray*}
	} 
\end{bt}

% nguồn: Ok_TuanLuan_Toan11_Lan2
\begin{ex}%[1D6V3-5]%[ma:5f210a06]
	Ông Đại mới xin được việc làm nên gửi tiết kiệm vào ngân hàng với hình thức cứ mỗi  đầu tháng đóng vào 5 triệu đồng với lãi suất $0{,}33\%/$ tháng. Tính số tiền mà ông Đại thu được từ ngân hàng sau 5 năm.
	\loigiai{ 
		Với $a$ là số  tiền ông Đại đóng vào hằng tháng;\\
		 $r\%$ lãi suất ông Đại gửi tiết kiệm hằng tháng.\\
		Gọi  $P_n$ là số tiền mà ông Đại thu được sau $n$ tháng $(n\ge 1)$.\\
		 Suy ra
		\begin{itemize}
			\item  $P_1=a\cdot(1+r\%)$.
			\item  $P_2=\left(P_1+a\right)(1+r\%)=a\cdot(1+r\%)^2+a\cdot(1+r\%)$.
			 \item $P_3=\left(P_2+a\right)(1+r\%)=a\cdot(1+r\%)^3+a\cdot(1+r\%)^2+a\cdot(1+r\%)$.
			\item  \dotfill
			 \item $P_n=\left(P_{n-1}+a\right)(1+r\%)=a\cdot(1+r\%)^n+a\cdot(1+r\%)^{n-1}+\cdots+a\cdot(1+r\%)$.
		\end{itemize}
		Xét cấp số nhân có số hạng đầu là $u_1=a\cdot(1+r\%)$ và công bội $q=1+r\%$ thì \begin{align*}
		P_n=u_1+u_2+\cdots+u_n=u_1\cdot  \dfrac{1-q^n}{1-q}.
	\end{align*}
	Vậy số tiền ông Đại nhận được từ ngân hàng sau $5$ năm  là
	\begin{align*}
		 P_{60}=u_1\cdot \dfrac{1-q^{60}}{1-q}=5\cdot\left(1{,}0033\right)\cdot\dfrac{1-\left(1{,}0033\right)^{60}}{0{,}0033} \approx 332\text{ (triệu đồng)}.
	\end{align*}
 }
\end{ex}

% nguồn: Ok_TuanLuan_Toan11_Lan2
\begin{ex}%[1D6V3-5]%[ma:1ab15813]
	Ông Bình vậy vốn ngân hàng với số tiền $100\,000\,000$ đồng. Ông dự định sau đúng $5$ năm thì trả hết  nợ theo hình thức sau đúng một tháng kể từ ngày vậy, ông bắt đầu hoàn nợ, hai lần hoàn nợ liên tiếp cách nhau đúng một tháng, số tiền hoàn nợ ở mỗi lần là như nhau. Hỏi theo cách đó, số tiền $a$ mà ông sẽ phải trả cho ngân hàng trong mỗi lần hoàn nợ là bao nhiêu? Biết lãi suất hàng tháng là $1{,}2\%$ và không thay đổi trong thời gian ông hoàn nợ.
	 \loigiai{
		Gọi $m$, $r$, $T_n$, $a$ lần lượt là số tiền vậy ngân hàng, lãi  suất hàng tháng, tổng số tiền vậy còn lại sau $n$ tháng, số tiền trả đều đặn mỗi tháng.
		\begin{itemize}
			\item  Sau khi hết tháng thứ nhất $(n=1)$ thì còn lại $T_1=m(r+1)-a$.
			\item  Sau khi hết tháng thứ hai $(n=2)$ thì còn lại
			\begin{eqnarray*}
				T_2&=&\left[m(r+1)-a\right](r+1)-a\\ 
				&=&m(r+1)^2-a(r+1)-a\\ 
				&=&m(r+1)^2-a(r+2)\\ 
				 &=&m(r+1)^2-\dfrac{a}{r} \left[(r+1)^2-1\right].
			\end{eqnarray*}
			\item Sau khi  hết tháng thứ ba $(n=3)$ thì còn
			\begin{eqnarray*}
				T_3&=&\left[m(r+1)^2-\dfrac{a}{r} \left[(r+1)^2-1\right]\right](r+1)-a\\ 
				&=&m(r+1)^3-\dfrac{a}{r}  \left[(r+1)^3-1\right].
			\end{eqnarray*}
			 \item \dotfill
			\item Sau khi hết  tháng thứ $n$ thì còn lại
			\begin{align*}
				 T_n=m(r+1)^n-\dfrac{a}{r} \left[(r+1)^n-1\right].
			\end{align*}
		\end{itemize}
		Áp dụng công thức  trên, ta có
		\begin{eqnarray*}
			T_n=0\Leftrightarrow a=\dfrac{m(r+1)^n  r}{(r+1)^n-1}=\dfrac{12\cdot 10^5 \Bigl(\dfrac{1{,}2}{100}+1\Bigr)^{60}}{\Bigl(\dfrac{1{,}2}{100}+1\Bigr)^{60}-1}.
		\end{eqnarray*}
	} 
\end{ex}

% nguồn: Ok_TuanLuan_Toan11_Lan2
\begin{ex}%[1D6V3-5]%[ma:fadddfd4]
	Một người gửi $100$ triệu đồng vào ngân hàng với lãi suất $0{,}4\% /$ tháng. Biết rằng nếu không rút tiền ta khỏi ngân hàng thì cứ sau mỗi tháng, số tiền  lãi sẽ được lập vào vốn ban đầu để tính lãi cho tháng tiếp theo. Hỏi sau $6$ tháng, người đó được lĩnh số tiền bao nhiêu, nếu trong khoảng thời gian này người đó không rút tiền ra và lãi xuất không thay đổi?
	\loigiai{ 
		Áp dụng công thức lãi kép ta có sau đúng $6$ tháng, người  đó lĩnh được số tiền
		\begin{align*}
			 A_n=A_0 (1+r)^n=100\,000\,000\left(1+\dfrac{0{,}4}{100} \right)^6=102\,424\,128\text{ (đồng)}.
		\end{align*}
	} 
\end{ex}

% nguồn: Ok_TuanLuan_Toan11_Lan2
\begin{ex}%[1D6V3-5]%[ma:23e4de17]
	Thực hiện một mẻ nuôi cấy vi khuẩn với $1000$ vi khuẩn ban đầu, nhà  sinh học phát hiện ra số lượng vi khuẩn tăng thêm $25 \%$ sau hai ngày.
	\begin{enumerate}
		\item Công thức $ P(t)= P_0 \cdot a^t$ cho phép tính số lượng vi khuẩn của mẻ nuôi cấy sau $t$ ngày kể từ thời điểm ban đầu. Xác định các tham số $P_0$ và $a$, $(a>0)$. Làm tròn đến hàng phần trăm. Lấy kết quả đã làm tròn ở ý a) để làm ý b) và ý c). 
		\item Sau $5$ ngày thì số lượng vi  khuẩn bằng bao nhiêu? Làm tròn kết quả đến hàng trăm.
		\item Sau bao nhiêu ngày thì số lượng  vi khuẩn bằng gấp đôi số lượng ban đầu? Làm tròn kết quả đến hàng phần chục.
	\end{enumerate}
	 \loigiai{
		\begin{enumerate}
			\item Ban  đầu có $1000$ vi khuẩn nên $P_0=1000$.\\
			Sau $2$ ngày, số lượng vi khuẩn là $P(2)=125 \%\cdot P_0=125  \% \cdot 1000=1250$.\\
			Ta có $P(2)=P_0 \cdot a^2 \Leftrightarrow 1250=1000  \cdot a^2 \Leftrightarrow a^2=1{,}25 \Leftrightarrow a \approx 1{,}12$.
			\item Số lượng vi khuẩn sau $5$  ngày là $P(5)=P_0 \cdot a^5=1000 \cdot 1{,}12^2 \approx 1800$ (vi khuẩn).
			\item Với $P(t)=2P_0$,  ta có\\
			$P(t)=P_0 \cdot a^t \Leftrightarrow 2 P_0=P_0 \cdot 1{,}12^t \Leftrightarrow 1{,}12^t=2 \Leftrightarrow t=\log  _{1,12} 2 \approx 6{,}1$ (ngày).\\
			Vậy sau $6{,}1$ ngày thì số lượng vi khuẩn vượt gấp đôi số lượng ban  đầu.
		\end{enumerate}
	} 
\end{ex}

% nguồn: Ok_TuanLuan_Toan11_Lan2
\begin{ex}%[1D6V3-5]%[ma:7b187321]
	Lãi suất cho vay tại PVcomBank trong tháng $5/2022$ rất ưu đãi, ở mức $5\%/$năm, được áp dụng trong $6$ tháng đầu, từ tháng thứ $7$ trở đi ấn định mức lãi $12\%/$năm. Tại ngân hàng này, thời hạn cho vậy mua nhà tối đa là $20$ năm, mức vậy tối đa $85\%$ giá trị tài sản đảm bảo. Một người có khả năng trả cố định hằng tháng là $15$ triệu. Giả sử người đó có thể mượn người thân $15\%$ giá trị căn nhà, nếu được sử dụng gói vậy ở trên với thời hạn tối đa và mức vậy  tối đa thì có thể mua được căn nhà có giá trị tối đa khoảng?
	\loigiai{ 
		Gọi $A$ là số tiền tối đa người này có thể vậy,  $A_i$ là số tiền nợ sau tháng thứ $i$.
		\begin{itemize}
			\item $r_1=\dfrac{5\%}{12}$ là lãi suất$/1$  tháng, trong $6$ tháng đầu.
			\item $r_2=\dfrac{12\%}{12}=1\%$ là lãi  suất$/1$ tháng, từ tháng thứ $7$ trở đi.
			\item Sau $1$ tháng, số tiền gốc và lãi là $A(1+r)$, người đó trả $15$ triệu nên còn nợ \begin{align*}
			 A_1=A(1+r)-15.
		\end{align*}
		\item Sau tháng  thứ $2$
		\begin{align*}
			A_2=A_1 (1+r_1)-15=\left(A(1+r_1)-15\right)(1+r_1)-15=A(1+r_1)^2-\dfrac{15}{r_1}  \left[((1+r_1))^2-1\right].
		\end{align*}
		\item Sau tháng thứ $3A_3=A(1+r_1)^3-\dfrac{15}{r_1} \left[(1+r_1)^3-1\right]$. 
		\item \dotfill. 
		\item Sau  tháng thứ $6A_6=A(1+r_1)^6-\dfrac{15}{r_1} \left[(1+r_1)^6-1\right]$.
		\item Sau  tháng thứ $7A_7=A_6 (1+r_2)-15$.
		\item  Sau tháng thứ $8A_8=A_6 (1+r_2)^2-\dfrac{15}{r_2} \left[(1+r_2)^2-1\right]$.
		\item  \dotfill
		\item  Sau tháng thứ $240A_{240}=A_6 (1+r_2)^{234}-\dfrac{15}{r_2} \left[(1+r_2)^{234}-1\right]$.
	\end{itemize}
	Vì  phải trả hết nợ sau $20$ năm nên
	\begin{eqnarray*}
		A_{240}=0&\Leftrightarrow& A_6=\dfrac{15\left[(1+r_2)^{234}-1\right]}{\left(1+r_2 \right)^{234} r_2}  \approx 1\,353{,}819328\\
		&\Rightarrow& A=\dfrac{A_6+\dfrac{15}{r_1} \left[(1+r_1)^6-1\right]}{\left(1+r_1  \right)^6} \approx 1\,409{,}163992.
	\end{eqnarray*}
	Vậy người này có thể mua được căn nhà có giá trị tối đa là \begin{align*}
	\dfrac{A}{85\%}  \approx 1\,657{,}83999\text{ (triệu đồng)} \approx 1{,}65784\text{ (tỷ đồng)}.
	\end{align*}
 }
\end{ex}

% nguồn: Ok_TuanLuan_Toan11_Lan2
\begin{ex}%[1D6V3-5]%[ma:abee462f]
	Một học sinh A khi đủ 18 tuổi được cha mẹ cho $200\, 000\, 000$ VNĐ. Số tiền này được bảo quản trong ngân hàng MSB với kì hạn thanh toán $1$ năm và học sinh A chỉ  nhận được số tiền này khi học xong $4$ năm đại học. Biết rằng khi đủ $22$ tuổi, số tiền mà học sinh A được nhận sẽ là $243\,101\, 250$ VNĐ. Vậy lãi suất kì hạn một năm của ngân hàng MSB là bao nhiêu?
	\loigiai{ 
		Gọi lãi suất kỳ hạn một năm của ngân hàng  MSB là $r$. \\
		Áp  dụng công thức lãi suất kép $P=a\left(1+r\right)^n$, ta có
		\begin{eqnarray*}
			243\,101\,250=200\,000\,000(1+r)^4  &\Leftrightarrow& (1+r)^4=\dfrac{243\,101\,250}{200\,000\,000}\\
			&\Leftrightarrow& 1+r=\sqrt[{4}]{\dfrac{243\,101\,250}{200\,000\,000}}\\ &\Leftrightarrow& r=\sqrt[{4}]{\dfrac{243\,101\,250}{200\,000\,000}}-1\\ 
			&\Leftrightarrow& r\approx 0{,}05. 
		\end{eqnarray*}
	} 
\end{ex}

% nguồn: Ok_TuanLuan_Toan11_Lan2
\begin{ex}%[1D6V3-5]%[ma:d02743d5]
	Một người gửi $50$ triệu đồng vào ngân hàng theo thể thức lãi kép với lãi suất $5{,}5\%$/năm, kì hạn $1$ năm. Hỏi sau  $4$ năm, người đó rút cả vốn lẫn lãi được số tiền bao nhiêu?
	\loigiai{ 
		 Gọi số tiền ban đầu $A$. Lãi suất tính theo năm là $r$.\\
		Số  tiền cả vốn lẫn lãi sau $n$ năm được tính theo công thức $A_n=A\left(1+r\right)^n$.\\
		Thay  số với $A=50$; $r=5{,}5\%$, $n=4$ ta được số tiền là $A_4=50\cdot\left(1+\dfrac{5{,}5}{100} \right)^4=61{,}94$.
	 }
\end{ex}

% nguồn: Ok_TuanLuan_Toan11_Lan2
\begin{ex}%[1D6V3-5]%[ma:2c2555e9]
	Một người gửi $200$ vào ngân hàng với lãi suất $0{,}2\%$/tháng. Biết rằng nếu không rút tiền khỏi ngân hàng thì cứ sau mỗi tháng số tiền lãi sẽ được nhập vào vốn ban đầu để tính lãi cho tháng tiếp theo. Hỏi sau đúng $10$ tháng người đó được  lĩnh số tiền bao nhiêu?
	\loigiai{ 
		Theo công thức lãi kép  ta có số tiền cả lãi và vốn sau $10$ tháng là
		\[T=X\left(1+r\right)^n=200\left(1+0{,}2\% \right)^{10} \approx 204{,}036\text{  triệu đồng}.\]
	} 
\end{ex}

% nguồn: Ok_TuanLuan_Toan11_Lan2
\begin{ex}%[1D6V3-5]%[ma:0d9e8609]
	Ông Bình vậy vốn ngân hàng với số tiền $100\,000\,000$ đồng. Ông dự định sau đúng $5$ năm thì trả hết nợ theo hình thức: sau đúng một tháng kể từ ngày vay, ông bắt đầu hoàn nợ, hai lần hoàn nợ liên tiếp cách nhau đúng một tháng, số tiền hoàn nợ ở mỗi  lần là như nhau. Hỏi theo cách đó, số tiền $a$ mà ông sẽ phải trả cho ngân hàng trong mỗi lần hoàn nợ là bao nhiêu? Biết lãi suất hàng tháng là $1{,}2\%$ và không thay đổi trong thời gian ông hoàn nợ.
	\loigiai{ 
		Gọi $m$, $r$, $T_n$, $a$ lần lượt là số tiền vay ngân hàng, lãi suất hàng tháng, tổng số tiền vậy còn lại sau $n$  tháng, số tiền trả đều đặn mỗi tháng.\\
		Sau khi  hết tháng thứ nhất ($n=1$) thì còn lại $T_1=m(r+1)-a$.\\
		Sau khi hết tháng thứ hai ($n=2$) thì còn lại \begin{eqnarray*}
		 T_2&=&[m(r+1)-a](r+1)-a=m(r+1)^2-a(r+1)-a=m(r+1)^2-a(r+2)\\
		&=&m(r+1)^2-a\dfrac{[(r+1)^2-1]}{r}. 
	\end{eqnarray*}
	Sau khi hết tháng thứ ba  ($n=3$) thì còn
	\begin{eqnarray*}
		 T_3&=&[m(r+1)^2-a\dfrac{[(r+1)^2-1]}{r}](r+1)-a\\
		 &=&m(r+1)^3-a\dfrac{[(r+1)^3-1]}{r}.
	\end{eqnarray*}
	\vdots\\ 
	Sau khi hết tháng thứ $n$ thì còn lại  $T_n=m(r+1)^n-a\dfrac{[(r+1)^n-1]}{r}$.\\
	Áp dụng công thức  trên, ta có $$T_n=0\Leftrightarrow a=\dfrac{m(r+1)^n r}{(r+1)^n-1}=\dfrac{12\cdot 10^7\left(1+\dfrac{1{,}2}{100}\right)^{60} \dfrac{1{,}2}{100}}{\left(1+\dfrac{1{,}2}{100}\right)^{60}-1}\approx 2\,817\,138\ \text{đồng}.$$
	Số tiền mà ông Bình phải trả cho ngân hàng trong mỗi lần hoàn nợ là $a\approx 2\,817\,138$ đồng. 
} 
\end{ex}

% nguồn: Ok_TuanLuan_Toan11_Lan2
\begin{ex}%[1D6V3-5]%[ma:5b2b1ca4]
	Hai anh em An và Bình cùng vay tiền ở ngân hàng với lãi suất $0{,}7\%$ một tháng với tổng số tiền vay là $200$ triệu đồng. Sau đúng $1$ tháng kể từ khi vay, mỗi người bắt đầu trả nợ cho ngân hàng khoản vay của mình. Mỗi tháng hai người trả số tiền bằng nhau cho ngân hàng để trừ vào tiền gốc và lãi. Để trả  hết gốc và lãi cho ngân hàng thì An cần $10$ tháng, Bình cần $15$ tháng. Hỏi số tiền mà mỗi người trả cho ngân hàng mỗi tháng là bao nhiêu?
	\loigiai{ 
		Gọi số tiền vay ban đầu là $u_0$,  tiền trả hàng tháng là $x$, lãi suất hàng tháng là $0{,}7\%$.\\
		Số tiền  còn lại sau $1$ tháng là $u_1=u_0\cdot 1{,}007- x$.
		Số tiền  còn lại sau $2$ tháng là
		$$u_2=u_11{,}007-  x=u_01{,}007^2-1{,}007x - x=u_01{,}007^2- x(1+1{,}007).$$
		Số tiền còn lại  sau $n$ tháng là
		$$u_n=u_0\cdot 1{,}007^n-x\left(1+1{,}007+1{,}007^2+\cdots+1{,}007^{n-1}\right)=u_0\cdot 1{,}007^n-x\cdot\dfrac{1{,}007^n-1}{0,007}.$$ 
		Sau $n$ tháng thì hết nợ nên  $u_n=0\Leftrightarrow u_0=\dfrac{x\left(1{,}007^n-1\right)}{0,007\cdot 1{,}007^n}$.\\
		Để trả  hết nợ thì An cần $10$ tháng và Bình cần $15$ tháng, ta được:
		$$\dfrac{x\left(1{,}007^{10}-1\right)}{0,007\cdot 1{,}007^{10}}+\dfrac{x\left(1{,}007^{15} - 1\right)}{0,007\cdot  1{,}007^{15}}=2\cdot 10^8\Rightarrow x\approx 8397068{,}067.$$
	 }
\end{ex}

% nguồn: Ok_TuanLuan_Toan11_Lan2
\begin{ex}%[1D6V3-5]%[ma:dcdc3391]
	Ông Đại mới xin được việc làm nên gửi tiết kiệm vào ngân hàng với hình thức cứ mỗi đầu tháng đóng vào $5$ triệu đồng với lãi suất $0{,}33\%$/tháng. Tính số tiền  mà ông Đại thu được từ ngân hàng sau $5$ năm.
	\loigiai{ 
		Với $a$ là số tiền ông Đại đóng vào hằng tháng, $r$\,(\%) là lãi suất ông  Đại gửi tiết kiệm hằng tháng.\\
		Gọi $P_n$ là số tiền mà ông Đại thu được sau $n$ tháng  $(n\ge 1)$.
		Suy  ra
		\allowdisplaybreaks 
		\begin{eqnarray*}
			&&P_1=a\cdot\left(1+r\right).\\ 
			 &&P_2=\left(P_1+a\right)\left(1+r\right)=a\cdot \left(1+r\right)^2+a\cdot \left(1+r \right).\\
			&&P_3=\left(P_2+a\right)\left(1+r  \right)=a\cdot\left(1+r \right)^3+a\cdot\left(1+r \right)^2+a\cdot\left(1+r \right).\\
			 &&\dots\\
			&&P_n=\left(P_{n-1}+a\right)\left(1+r  \right)=a\cdot\left(1+r \right)^n+a\cdot\left(1+r \right)^{n-1}+\cdots+a\cdot\left(1+r \right).
		\end{eqnarray*}
		Xét cấp số nhân có số hạng đầu là $u_1=a\cdot\left(1+r \right)$ và  công bội $q=1+r$ thì
		\[P_n=u_1+u_2+\cdots+u_n=u_1\cdot \dfrac{1-q^n}{1-q}.\] 
		Vậy  số tiền ông Đại nhận được từ ngân hàng sau $5$ năm là
		\[P_{60}=u_1\cdot \dfrac{1-q^{60}}{1-q}=5\cdot\left(1{,}0033\right)\cdot\dfrac{1-\left(1{,}0033\right)^{60}}{0{,}0033} \approx  332\text{ triệu đồng}.\]
	} 
\end{ex}

% nguồn: Ok_TuanLuan_Toan11_Lan2
\begin{ex}%[1D6V3-5]%[ma:3d155c97]
	Chu kì bán rã của chất phóng xạ Plutolium $^{239}Pu$ là  $24\,360$ năm (tức là một lượng chất $^{239}Pu$ sau $24360$ năm phân hủy còn một nửa). Sự phân hủy này được tính theo công thức $S=A\mathrm{e}^{-rt}$, trong đó $A$ là lượng chất phóng xạ ban đầu, $r$ là tỉ lệ phân hủy hàng năm, $t$ là thời gian phân hủy, $S$ là lượng còn lại sau thời gian phân hủy $t$. Hỏi $20$ gam $^{239}Pu$ sau ít nhất bao nhiêu năm thì phân hủy còn $4$ gam?
	 \loigiai{
		Vì  $^{239}Pu$ có chu kì bán rã là $24\,360$ năm nên với $20$ gam $^{239}Pu$ ta có
		$$10=20\cdot\mathrm{e}^{-r\cdot 24\,360}\Leftrightarrow-r\cdot  24360=\ln\dfrac{1}{2}\Leftrightarrow r=\dfrac{\ln 2}{24\,360}.$$
		Theo bài ra ta có  phương trình $4=20\cdot\mathrm{e}^{-rt}\Leftrightarrow-rt=\ln\dfrac{1}{5}\Leftrightarrow rt=\ln 5\Leftrightarrow t=\dfrac{\ln 5}{r}$.\\
		Suy ra  $t\approx 56\,562{,}2$.\\
	Vậy  sau ít nhất $56\,563$ năm thì $20$ gam $^{239}Pu$ sẽ phân hủy còn $4$ gam.}
\end{ex}

% nguồn: Ok_TuanLuan_Toan11_Lan2
\begin{ex}%[1D6V3-5]%[ma:53513dc1]
	Anh Nam vậy tiền ngân hàng $1$ tỷ đồng theo phương thức trả góp với lãi suất $0{,}5\%$/tháng. Nếu cuối mỗi tháng bắt đầu từ tháng thứ nhất anh Nam trả $30$  triệu đồng. Hỏi sau bao nhiêu tháng anh Nam trả hết nợ?
	\loigiai{ 
		Gọi $a$ là số tiền vay, $r$ là lãi suất, $m$ là số tiền hàng tháng trả. 
		Số tiền  nợ sau tháng thứ nhất là $N_1=a(1+r)-m$.\\
		Số  tiền nợ sau tháng thứ hai là $N_2=[a(1+r)-m]+[a(1-r)-m]r-m=a(1+r)^2-m[(1+r)+1]$.\\
		Số tiền nợ sau $n$ tháng là \begin{eqnarray*}
		 N_n&=&a(1+r)^n-m[(1+r)^{n-1}+\cdots+(1+r)^2+\cdots+1]\\
		&=&a(1+r)^n-m\cdot\dfrac{(1+r)^n-1}{r}. 
	\end{eqnarray*}
	Sau $n$ tháng anh Nam trả hết nợ nên \begin{eqnarray*}
	N_n=a(1+r)^n-m\cdot\dfrac{(1+r)^n-1}{r}=0&\Leftrightarrow &1000(1+0{,}005)^n -  30\dfrac{(1+0{,}005)^n -1}{0,005}=0\\
	&\Leftrightarrow& n \approx 36{,}55. 
	\end{eqnarray*}
Vậy $37$ tháng  thì anh Nam trả hết nợ.}
\end{ex}

% nguồn: Ok_TuanLuan_Toan11_Lan2
\begin{ex}%[1D6V3-5]%[ma:8e063d8b]
	Giả sử một chất phóng xạ bị phân rã theo cách sao cho khối lượng $m(t)$ của chất còn lại (tính bằng kilôgam) sau $t$ ngày được cho bởi hàm  số $m(t)=13\mathrm{e}^{-0{,}015t}$.
	\begin{enumerate}
		\item Tìm  khối lượng của chất đó tại thời điểm $t=0$.
		\item Sau $45$ ngày khối lượng chất đó còn lại là bao  nhiêu?
	\end{enumerate}
	 \loigiai{
		\begin{enumerate}
			\item  Khi $t=0$, $m(0)=13\mathrm{e}^{-0{,}015 \times 0}=13\mathrm{e}^0=13$.\\
			 Vậy khối lượng của chất phóng xạ ban đầu là $13$ kg.
			\item Để tìm khối lượng chất phóng xạ còn lại sau $45$ ngày, ta sử dụng công thức $m(t)=13\mathrm{e}^{-0{,}015t}$ và thay $t=45$ vào ta được $m(45)=13\mathrm{e}^{-0{,}015\cdot 45}  \approx 6{,}19$ {kg}.
		\end{enumerate}
	 }
\end{ex}

% nguồn: Ok_TuanLuan_Toan11_Lan2
\begin{ex}%[1D6V3-5]%[ma:1950dd7a]
	Ông A vậy ngân hàng $100$ triệu đồng với lãi suất $1\%$ một tháng. Cứ sau mỗi tháng kể từ ngày vay ông trả góp số tiền $5$ triệu đồng. Hỏi sau ít nhất bao nhiêu tháng thì ông A trả hết nợ, biết tháng cuối cùng ông có thể trả số tiền  ít hơn $5$ triệu đồng?
	 \loigiai{
		Sau  $n$ tháng, ông A còn vay số tiền là
		 $$100(1+r)^n-5[(1+r)^{n-1}+(1+r)^{n-2}+\cdots+1]=100(1+r)^n-5\dfrac{(1+r)^n-1}{r}$$
		với $r$ là lãi  suất $1$ tháng.\\
		Để tháng thứ $n$ ông trả hết nợ thì  $$100(1{,}01)^n - 5\dfrac{(1{,}01)^n-1}{0,01}=0\Leftrightarrow (1{,}01)^n=\dfrac{5}{4}\Leftrightarrow n=\log_{1{,}01}\dfrac{5}{4} \approx 23.$$
		Vậy sau ít nhât $23$ tháng  thì ông A trả hết nợ.
	} 
\end{ex}

% nguồn: Ok_TuanLuan_Toan11_Lan2
\begin{ex}%[1D6V3-5]%[ma:9046906d]
	Anh A vậy ngân hàng $600\,000\,000$ đồng để mua xe ô tô với lãi suất $7{,}8\%$ một năm. Anh A bắt đầu trả nợ cho ngân hàng theo cách: sau đúng $1$ năm kể từ ngày vậy anh bắt đầu trả nợ và hai lần trả nợ liên tiếp cách nhau đúng $1$ năm. Số tiền trả nợ là như nhau ở mỗi lần và sau đúng $8$ năm thì anh A trả hết nợ. Biết  rằng lãi suất ngân hàng không thay đổi trong suốt thời gian anh A trả nợ. Số tiền anh A trả nợ ngân hàng trong mỗi lần là bao nhiêu?
	 \loigiai{
		Gọi $M$ là số  tiền anh A trả hàng năm.\\
		Sau năm thứ $1$, số tiền còn lại $V_1=600(1+r)-M$.\\ 
		Sau năm thứ $2$, số  tiền còn lại $V_2=V_1(1+r)-M=600(1+r)^2- M (1+r)-M$.\\
		 \vdots\\
		Sau năm thứ $n$, số  tiền còn lại
		\begin{eqnarray*}
			 V_n&=&600(1+r)^n-M(1+r)^{n-1}-\cdots-M(1+r)-M\\
			&=&600(1+r)^n-M\cdot\dfrac{(1+r)^n-1}{r}. 
		\end{eqnarray*}
		Sau $8$ năm anh A trả hết nợ  nên ta có
		\begin{eqnarray*}
			600(1+r)^8-M\cdot\dfrac{(1+r)^8-1}{r}=0&\Leftrightarrow&  M=\dfrac{600(1+r)^8.r}{(1+r)^8-1}\\
			&\Leftrightarrow &M=\dfrac{600(1+7{,}8\%)^8\cdot  7{,}8\%}{(1+7{,}8\%)^8-1} \approx 103{,}618\ \text{triệu đồng}.
		\end{eqnarray*}
		Số tiền mà  anh A trả nợ ngân hàng mỗi lần xấp xỉ $103{,}618$ triệu đồng.
	} 
\end{ex}

% nguồn: Ok_TuanLuan_Toan11_Lan2
\begin{ex}%[1D6V3-5]%[ma:e5708c55]
	Một người gửi số tiền $300$ triệu đồng vào một ngân hàng với lãi suất kép $6\%$ một năm. Biết rằng nếu không rút tiền ra khỏi ngân hàng thì cứ sau mỗi năm,  số tiền lãi suất sẽ được nhập vào vốn ban đầu. Hỏi sau $3$ năm không rút tiền gốc và lãi, số tiền trong ngân hàng của người đó là bao nhiêu?
	\loigiai{ 
		Áp dụng công thức tính lãi suất theo  hình thức lãi kép $P=A\left(1+r\right)^n$.\\
		Trong đó: $P$ là số tiền gồm vốn lẫn lãi tại thời điểm $n$ tính từ thời điểm gửi; $A$ là  số tiền gửi vào ban đầu và $r\left(\% \right)$ là lãi suất.
		Với  $\heva{&A=300\,000\,000\\&n=3\\&r=6\%.}$\\
		Khi  đó $P=300\,000\,000\left(1+6\% \right)^3=357\,304\,800\approx 357\,305\,000$.
	 }
\end{ex}

% nguồn: Ok_TuanLuan_Toan11_Lan2
\begin{ex}%[1D6V3-5]%[ma:c636bed4]
	Một người gửi $100$ triệu đồng vào ngân hàng với lãi suất $0{,}4$\%/tháng. Biết rằng nếu không rút tiền ta khỏi ngân hàng thì cứ sau mỗi tháng, số tiền lãi sẽ được lập vào vốn ban đầu để tính lãi cho tháng tiếp theo. Hỏi sau $6$ tháng, người đó được lĩnh số tiền bao nhiêu, nếu trong khoảng thời gian này người đó  không rút tiền ra và lãi xuất không thay đổi?
	 \loigiai{
		Áp dụng công thức lãi kép ta có sau đúng  $6$ tháng, người đó lĩnh được số tiền là
		 \[A_n=A_0 (1+r)^n=10^8\cdot \left(1+\dfrac{0{,}4}{100} \right)^6=102\,424\,128.\]
	} 
\end{ex}

% nguồn: Ok_TuanLuan_Toan11_Lan2
\begin{ex}%[1D6V3-5]%[ma:e133267d]
	Lãi suất cho vay tại PVcomBank trong tháng 5/2022 rất ưu đãi, ở mức $5\%$/năm, được  áp dụng trong $6$ tháng đầu, từ tháng thứ $7$ trở đi ấn định mức lãi $12\%$/năm. Tại ngân hàng này, thời hạn cho vay mua nhà tối đa là $20$ năm, mức vay tối đa $85\%$ giá trị tài sản đảm bảo. Một người có khả năng trả cố định hằng tháng là $15$ triệu. Giả sử người đó có thể mượn người thân $15\%$ giá trị căn nhà, nếu được sử dụng gói vay ở trên với thời hạn tối đa và mức vay tối đa thì có thể mua được căn nhà có giá trị tối đa bao nhiêu tiền?
	\loigiai{ 
		Gọi $A$ là số tiền tối đa người này có thể vậy, $A_i$ là số tiền  nợ sau tháng thứ $i$.\\
		$r_1=\dfrac{5\%}{12}$ là lãi suất/$1$ tháng,  trong $6$ tháng đầu.\\
		$r_2=\dfrac{12\%}{12}=1\%$ là lãi suất/1 tháng, từ tháng thứ  $7$ trở đi.\\
		Sau 1 tháng, số tiền gốc và lãi là $A\left(1+r\right)$, người  đó trả $15$ triệu nên còn nợ $A_1=A\left(1+r\right)-15$.\\
		Sau  tháng thứ 2:
		 \allowdisplaybreaks
		\begin{eqnarray*}
			A_2&=&A_1  \left(1+r_1 \right)-15\\
			&=&\left(A\left(1+r_1 \right)-15\right)\left(1+r_1 \right)-15\\ 
			&=&A\left(1+r_1  \right)^2-\dfrac{15}{r_1} \left[(1+r_1)^2-1\right].
		\end{eqnarray*}
		Sau tháng thứ 3: $A_3=A\left(1+r_1  \right)^3-\dfrac{15}{r_1} \left[(1+r_1)^3-1\right]$.\\
		 {\dots}\\
		 Sau tháng thứ 6: $A_6=A\left(1+r_1 \right)^6-\dfrac{15}{r_1} \left[(1+r_1)^6-1\right]$.\\
		Sau tháng thứ  7: $A_7=A_6 \left(1+r_2 \right)-15$.\\
		Sau tháng thứ 8:  $A_8=A_6 \left(1+r_2 \right)^2-\dfrac{15}{r_2} \left[(1+r_2)^2-1\right]$.\\
		 {\dots}\\
		Sau tháng thứ $240$: $A_{240}=A_6  \left(1+r_2 \right)^{234}-\dfrac{15}{r_2} \left[(1+r_2)^{234}-1\right]$.\\
		Vì phải trả  hết nợ sau $20$ năm nên
		\allowdisplaybreaks 
		\begin{eqnarray*}
			 &&A_{240}=0\\
			&\Leftrightarrow& A_6=\dfrac{15\left[(1+r_2)^{234}-1\right]}{\left(1+r_2 \right)^{234}  r_2} \approx 1\,353,819.\\
			&\Rightarrow& A=\dfrac{A_6+\dfrac{15}{r_1} \left[(1+r_1)^6-1\right]}{\left(1+r_1 \right)^6}  \approx 1\,409{,}164.
		\end{eqnarray*}
		Vậy người này có thể  mua được căn nhà có giá trị tối đa là
		\[\dfrac{A}{85\%} \approx 1\,657{,}84 \text{ triệu đồng} \approx 1{,}65784\text{  tỷ đồng}.\]
	} 
\end{ex}

% nguồn: Ok_TuanLuan_Toan11_Lan2
\begin{ex}%[1D6V3-5]%[ma:cfe847c1]
	Một học sinh A khi đủ $18$ tuổi được cha mẹ cho $200\,000\,000$ VNĐ. Số tiền này được gửi  vào ngân hàng MSB với kì hạn thanh toán $1$ năm và học sinh A chỉ nhận được số tiền này khi học xong $4$ năm đại học. Biết rằng khi đủ $22$ tuổi, số tiền mà học sinh A được nhận sẽ là $243\,101\,250$ VNĐ. Hỏi lãi suất kì hạn một năm của ngân hàng MSB là bao nhiêu?
	\loigiai{ 
		Gọi lãi suất kỳ hạn một năm của ngân hàng MSB là  $r$.\\
		Áp dụng  công thức lãi suất kép
		$P=a\left(1+r\right)^n$.\\  Khi đó
		 \allowdisplaybreaks
		\begin{eqnarray*}
			&&243\,101\,250=200\,000\,000\left(1+r\right)^4\\ 
			 &\Leftrightarrow& \left(1+r\right)^4=\dfrac{243\,101\,250}{200\,000\,000}\\ &\Leftrightarrow& 1+r=\sqrt[4]{\dfrac{243\,101\,250}{200\,000\,000}}\\ &\Leftrightarrow& r=\sqrt[4]{\dfrac{243\,101\,250}{200\,000\,000}}-1\\
			&\Leftrightarrow& r=0{,}05. 
		\end{eqnarray*}
	 }
\end{ex}

% nguồn: Ok_TuanLuan_Toan11_Lan2
\begin{ex}%[1D6V3-5]%[ma:b16d0094]
	Một nhóm bạn thực hiện dự án khởi nghiệp làm tinh dầu tự nhiên từ cây xả. Trong bản kế hoạch nhóm đề ra vậy ngân hàng $300$ triệu đồng theo phương thức trả góp với lãi suất $0{,}5\%$/tháng. Nếu cuối mỗi tháng bắt đầu từ tháng thứ chín nhóm bắt đầu trả trả $10$ triệu  đồng. Hỏi sau bao nhiêu tháng kể từ ngày vậy nhóm trả hết nợ?
	 \loigiai{
		Gọi $a$ là tổng số tiền vậy số tiền  vậy ngân hàng sau $8$ tháng, $r$ là lãi, $m$ là số tiền hàng tháng trả.\\
		 Ta có $a=300(1+0{,}5\%)$.\\
		Số tiền nợ sau tháng thứ chín là $N_1=a(1+r)-m$.\\ 
		Số tiền nợ sau  tháng thứ mười là
		$N_2=[a(1+r)-m]+[a(1-r)-m]r-m=a(1+r)^2-m[(1+r)+1]$.\\ 
		Số tiền  nợ sau tháng thứ $n+9$ là $N_n=a(1+r)^n-m\cdot\dfrac{(1+r)^n-1}{r}$.\\
		Sau $n+9$ tháng  trả hết nợ nên
		\begin{eqnarray*}
			 &&N_n=a(1+r)^n-m\cdot\dfrac{(1+r)^n-1}{r}=0\\
			&\Leftrightarrow  & \left[300\cdot (1+0{,}5\%)^8\right](1+0{,}5\%)^n-10\cdot \dfrac{(1+0{,}5\%)^n-1}{0{,}5\%}=0\\
			&\Rightarrow  &n \approx 23.
		\end{eqnarray*}
		Vậy sau $31$  tháng thì nhóm bạn trả hết nợ.
	} 
\end{ex}

% nguồn: Ok_TuanLuan_Toan11_Lan2
\begin{ex}%[1D6V3-5]%[ma:653edfb2]
	Số người  trong cộng đồng sinh viên đã nghe một tin đồn nào đó là $N=P\left(1-\mathrm{e}^{-0,15d} \right)$ trong đó $P$ là tổng số sinh viên của cộng đồng và $d$ là số ngày trôi qua kể từ khi tin đồn bắt đầu. Trong một cộng đồng $1\,000$ sinh viên, cần bao nhiêu ngày để $450$ sinh viên nghe được tin đồn?
	 \loigiai{
		Ta có 
		\allowdisplaybreaks 
		\begin{eqnarray*}
			 &&N=P\left(1-\mathrm{e}^{-0,15d} \right)\\
			&\Leftrightarrow& 450=1\,000\cdot\left(1-\mathrm{e}^{-0,15d} \right)\\ &\Leftrightarrow& \mathrm{e}^{-0,15d}=\ln  \dfrac{11}{20}\\
			&\Leftrightarrow& d\approx 3{,}98. 
		\end{eqnarray*}
		Vậy cần  $4$ ngày để $450$ sinh viên nghe được tin đồn.
	} 
\end{ex}

% nguồn: Ok_TuanLuan_Toan11_Lan1
\begin{bt}%[1D6C3-5]%[ma:81c78938]
	Mức cường độ âm $L$ (đơn vị: $dB$) được tính bởi công thức $L=10\log\left(\dfrac{I}{10^{-12}}\right)$, trong đó $I$ (đơn vị: W/m$^2$) là cường độ âm. Mức cường độ âm ở một  khu dân cư được quy định là dưới $60\ dB$. Hỏi cường độ âm của khu vực đó phải dưới bao nhiêu W/m$^2$?
	\loigiai{ 
		 Ta có $L<60\Leftrightarrow10\log\left(\dfrac{I}{10^{-12}}\right)<60\Leftrightarrow\log\left(\dfrac{I}{10^{-12}}\right)<6$\\
		$\Leftrightarrow\dfrac{I}{10^{-12}}<10^6\Leftrightarrow I<10^{-6}$ (do  $10>1$).\\
		Vậy cường độ âm ở khu vực đó phải dưới $10^{-6}$  (W/m$^2$).
	} 
\end{bt}

% nguồn: Ok_TuanLuan_Toan11_Lan1
\begin{bt}%[1D6C3-5]%[ma:1623609b]
	Dân số thế giới được tính theo công thức $S=A\mathrm{e}^{nr}$ trong đó $A$ là dân số của năm lấy làm mốc tính, $S$ là dân số sau $n$ năm, $r$ là tỉ lệ tăng dân số hằng năm. Cho biết năm 2005 Việt Nam có khoảng $80902400$ người và tỉ lệ tăng dân số là $1,47\%$ một năm. Như vậy, nếu tỉ lệ tăng  dân số hàng năm không đổi thì tối thiểu đến năm bao nhiêu dân của Việt Nam có khoảng $93713000$ người?
	\loigiai{ 
		 Ta có $S=A\cdot\mathrm{e}^{nr}\Leftrightarrow\mathrm{e}^{nr}=\dfrac{S}{A}\Leftrightarrow nr=\ln\dfrac{S}{A}\Leftrightarrow n=\dfrac{1}{r}\ln\dfrac{S}{A}$\\
		với $S=93713700$ người, $A=80902400$ người,  $r=\dfrac{1,47}{100}=0,0147$.\\
		Suy  ra $n=\dfrac{1}{0,0147}\ln\dfrac{93713000}{80902400}\approx10$.\\
		 Vậy tối thiểu đến năm 2015 thì dân số của Việt Nam có khoảng $93713000$ người.
	} 
\end{bt}

% nguồn: Ok_TuanLuan_Toan11_Lan1
\begin{bt}%[1D6C3-5]%[ma:b07438f5]
	Người ta xây dựng một cây cầu vượt giao thông hình parabol nối hai điểm có khoảng cách  là $400$m. Độ dốc của mặt cầu không vượt quá $10^{\circ}$ (độ dốc tại một điểm được xác định bởi góc giữa phương tiếp xúc với mặt cầu và phương ngang). Tính chiều cao giới hạn từ đỉnh cầu đến mặt đường (làm tròn kết quả đến chữ số thập phân thứ nhất).
	\begin{center}
		\begin{tikzpicture}[>=latex,line join=round,line cap=round,scale=1]
				\coordinate (M)  at (-1.5,1);
				 \coordinate (P) at (-3,1);
				\coordinate (N)  at (1,4);
				\draw[samples=200,smooth,red,line width=1] plot[domain=-4:4] (\x,{sqrt  (25-(\x)^2)-3})--(-4,0);
				\draw[samples=200,smooth,blue,line width=1]  plot[domain=-4:-1] (\x,{(3*(\x)+13)/4});
				\draw[dashed] (-4,0)  node[below]{$A$}
					(4,0)  node[below]{$B$}
					(0,0) node[below]{$400$m}--(0,1)  node[right]{?}--(0,2)
					(P)  node[above]{$P$}--(M);
				 \draw pic["",draw=blue,-,angle eccentricity=1.2,angle radius=0.4cm] {angle=M--P--N};
			\end{tikzpicture}
	\end{center}
	 \loigiai{
		Chọn hệ  trục toạ độ như hình vẽ, sao cho đỉnh cầu là gốc tọa độ và mặt cắt của cây cầu có hình dạng parabol $y=-ax^2$ (với $a$ là hằng số dương).
		\begin{center}
			\begin{tikzpicture}[>=latex,line join=round,line cap=round,scale=1]
					\draw[->]  (-5,2)--(5,2) node[above]{$x$};
					\draw[->] (0,-1)--(0,3)  node[right]{$y$};
					\coordinate  (M) at (-1.5,1);
					\coordinate (P) at  (-3,1);
					\coordinate (N) at (1,4); 
					 \draw[dashed] (-4,2) node[above]{$-200$}|-(0,0) node[below right]{$I$}-|(4,2) node[above]{$200$};
					\draw[samples=200,smooth,red,line width=1] plot[domain=-4:4] (\x,{sqrt  (25-(\x)^2)-3})--(-4,0);
					\draw[samples=200,smooth,blue,line  width=1] plot[domain=-4:-1] (\x,{(3*(\x)+13)/4});
					\draw[dashed] (-4,0) node[below]{$A$} 
						(4,0)  node[below]{$B$}
						(0,0)--(0,1)  node[right]{?}--(0,2)
						(P) node[above]{$P$}--(M); 
					\draw  pic["",draw=blue,-,angle eccentricity=1.2,angle radius=0.4cm] {angle=M--P--N};
				\end{tikzpicture}
		\end{center}
		Hệ số góc của tiếp tuyến của parabol  bằng $k=y^{\prime}\left(x_0\right)=-2ax_0$, $-200\leq x_0\leq200$.\\
		Hệ số góc xác định độ dốc của mặt cầu (độ dốc dương)  là $|k|=2a|x|\leq400a$.\\
		Vì độ dốc của mặt cầu không  vượt quá $10^{\circ}$ nên ta có:
		$400a\leq\tan10^{\circ}\Leftrightarrow  a\leq\dfrac{4,408174518}{10000}$.\\
		Chiều cao giới hạn từ đỉnh cầu đến mặt đường là đoạn $OI$, cũng chính là độ  lớn của tung độ điểm $B$ khi $a$ đạt giá trị lớn nhất.\\
		Do đó, $OI=\left|-a\cdot200^2\right|=17,6$  (m).\\
		Vậy chiều cao giới hạn từ đỉnh cầu đến mặt đường là $17,6$ m. 
	} 
\end{bt}

% nguồn: Ok_TuanLuan_Toan11_Lan2
\begin{ex}%[1D6C3-5]%[ma:0a6a3d6e]
	Mức cường độ âm $L$ (đơn vị: $dB$) được tính bởi công thức $L=10\log\left(\dfrac{I}{10^{-12}}\right)$, trong đó $I$ (đơn vị: W/m$^2$) là cường độ âm. Mức cường độ âm ở một khu dân cư được quy định  là dưới $60\ dB$. Hỏi cường độ âm của khu vực đó phải dưới bao nhiêu W/m$^2$?
	 \loigiai{
		Ta  có $L<60\Leftrightarrow10\log\left(\dfrac{I}{10^{-12}}\right)<60\Leftrightarrow\log\left(\dfrac{I}{10^{-12}}\right)<6$\\
		$\Leftrightarrow\dfrac{I}{10^{-12}}<10^6\Leftrightarrow  I<10^{-6}$ (do $10>1$).\\
		Vậy cường độ âm ở khu vực  đó phải dưới $10^{-6}$ (W/m$^2$).
	 }
\end{ex}

% nguồn: Ok_TuanLuan_Toan11_Lan2
\begin{ex}%[1D6C3-5]%[ma:9c83ce5a]
	Dân số thế giới được tính  theo công thức $S=A\mathrm{e}^{nr}$ trong đó $A$ là dân số của năm lấy làm mốc tính, $S$ là dân số sau $n$ năm, $r$ là tỉ lệ tăng dân số hằng năm. Cho biết năm 2005 Việt Nam có khoảng $80902400$ người và tỉ lệ tăng dân số là $1,47\%$ một năm. Như vậy, nếu tỉ lệ tăng dân số hàng năm không đổi thì tối thiểu đến năm bao nhiêu dân của Việt Nam có khoảng $93713000$ người?
	\loigiai{ 
		Ta có $S=A\cdot\mathrm{e}^{nr}\Leftrightarrow\mathrm{e}^{nr}=\dfrac{S}{A}\Leftrightarrow  nr=\ln\dfrac{S}{A}\Leftrightarrow n=\dfrac{1}{r}\ln\dfrac{S}{A}$\\
		với $S=93713700$ người, $A=80902400$  người, $r=\dfrac{1,47}{100}=0,0147$.\\
		 Suy ra $n=\dfrac{1}{0,0147}\ln\dfrac{93713000}{80902400}\approx10$.\\
		Vậy tối thiểu đến năm 2015 thì dân  số của Việt Nam có khoảng $93713000$ người.
	 }
\end{ex}

% nguồn: Ok_TuanLuan_Toan11_Lan2
\begin{ex}%[1D6C3-5]%[ma:5e0be94c]
	Người ta xây dựng một cây cầu vượt giao thông hình parabol nối hai điểm có khoảng cách là $400$m. Độ dốc của mặt cầu không vượt quá $10^{\circ}$ (độ dốc tại một điểm được xác định bởi góc giữa phương tiếp xúc với mặt cầu và phương ngang). Tính chiều cao giới hạn từ đỉnh cầu  đến mặt đường (làm tròn kết quả đến chữ số thập phân thứ nhất).
	\begin{center}
		\begin{tikzpicture}[>=latex,line join=round,line cap=round,scale=1]
				\coordinate (M)  at (-1.5,1);
				 \coordinate (P) at (-3,1);
				\coordinate (N) at (1,4); 
				\draw[samples=200,smooth,red,line  width=1] plot[domain=-4:4] (\x,{sqrt (25-(\x)^2)-3})--(-4,0);
				\draw[samples=200,smooth,blue,line  width=1] plot[domain=-4:-1] (\x,{(3*(\x)+13)/4});
				 \draw[dashed] (-4,0) node[below]{$A$}
					 (4,0) node[below]{$B$}
					(0,0)  node[below]{$400$m}--(0,1) node[right]{?}--(0,2)
					(P) node[above]{$P$}--(M); 
				 \draw pic["",draw=blue,-,angle eccentricity=1.2,angle radius=0.4cm] {angle=M--P--N};
			\end{tikzpicture}
	\end{center}
	 \loigiai{
		Chọn hệ trục toạ độ như hình vẽ, sao cho đỉnh cầu là gốc tọa  độ và mặt cắt của cây cầu có hình dạng parabol $y=-ax^2$ (với $a$ là hằng số dương).
		\begin{center}
			\begin{tikzpicture}[>=latex,line join=round,line cap=round,scale=1]
					\draw[->] (-5,2)--(5,2)  node[above]{$x$};
					\draw[->] (0,-1)--(0,3) node[right]{$y$}; 
					\coordinate  (M) at (-1.5,1);
					\coordinate  (P) at (-3,1);
					\coordinate (N) at  (1,4);
					\draw[dashed] (-4,2) node[above]{$-200$}|-(0,0)  node[below right]{$I$}-|(4,2) node[above]{$200$};
					\draw[samples=200,smooth,red,line  width=1] plot[domain=-4:4] (\x,{sqrt (25-(\x)^2)-3})--(-4,0);
					\draw[samples=200,smooth,blue,line width=1] plot[domain=-4:-1]  (\x,{(3*(\x)+13)/4});
					\draw[dashed] (-4,0)  node[below]{$A$}
						 (4,0) node[below]{$B$}
						(0,0)--(0,1) node[right]{?}--(0,2) 
						(P)  node[above]{$P$}--(M);
					\draw pic["",draw=blue,-,angle  eccentricity=1.2,angle radius=0.4cm] {angle=M--P--N};
				\end{tikzpicture}
		\end{center}
		Hệ số góc của tiếp tuyến của parabol bằng $k=y^{\prime}\left(x_0\right)=-2ax_0$, $-200\leq  x_0\leq200$.\\
		Hệ số góc xác định độ dốc của mặt  cầu (độ dốc dương) là $|k|=2a|x|\leq400a$.\\
		Vì độ  dốc của mặt cầu không vượt quá $10^{\circ}$ nên ta có:
		 $400a\leq\tan10^{\circ}\Leftrightarrow a\leq\dfrac{4,408174518}{10000}$.\\
		Chiều cao giới  hạn từ đỉnh cầu đến mặt đường là đoạn $OI$, cũng chính là độ lớn của tung độ điểm $B$ khi $a$ đạt giá trị lớn nhất.\\
		 Do đó, $OI=\left|-a\cdot200^2\right|=17,6$ (m).\\
		 Vậy chiều cao giới hạn từ đỉnh cầu đến mặt đường là $17,6$ m.
	 }
\end{ex}

% nguồn: Ok_TuanLuan_Toan11_Lan2
\begin{ex}%[1D6C3-5]%[ma:e3c81261]
	Ông Bình vay vốn ngân hàng với số tiền $100\,000\,000$ đồng. Ông dự định sau đúng $5$ năm thì trả hết nợ theo hình thức: sau đúng một tháng kể từ ngày vay, ông bắt đầu hoàn nợ, hai lần hoàn nợ liên tiếp cách nhau đúng một tháng, số tiền hoàn nợ ở  mỗi lần là như nhau. Hỏi theo cách đó, số tiền $a$ mà ông sẽ phải trả cho ngân hàng trong mỗi lần hoàn nợ là bao nhiêu? Biết lãi suất hàng tháng là $1{,}2\%$ và không thay đổi trong thời gian ông hoàn nợ.
	 \loigiai{
		Gọi $m$, $r$, $T_n$, $a$ lần lượt là số tiền vay ngân hàng, lãi suất hàng tháng, tổng số tiền vay  còn lại sau $n$ tháng, số tiền trả đều đặn mỗi tháng.\\
		Sau khi hết  tháng thứ nhất $\left(n=1\right)$ thì còn lại là $T_1=m\left(r+1\right)-a$.\\
		Sau khi  hết tháng thứ hai $\left(n=2\right)$ thì còn lại là
		 \allowdisplaybreaks
		\begin{eqnarray*}
			T_2&=&\left[m\left(r+1\right)-a\right]\left(r+1\right)-a\\ 
			&=&m\left(r+1\right)^2-a\left(r+1\right)-a\\ 
			&=&m\left(r+1\right)^2-a\left(r+2\right)\\ 
			&=&m\left(r+1\right)^2-\dfrac{a}{r} \left[\left(r+1\right)^2-1\right]. 
		\end{eqnarray*}
		Sau khi hết tháng  thứ ba $(n=3)$ thì còn lại là
		 \allowdisplaybreaks
		\begin{eqnarray*}
			T_3&=&\left[m\left(r+1\right)^2-\dfrac{a}{r} \left[\left(r+1\right)^2-1\right]\right]\left(r+1\right)-a\\ 
			&=&m\left(r+1\right)^3-\dfrac{a}{r}  \left[\left(r+1\right)^3-1\right].
		\end{eqnarray*}
		 \dots\\
		Sau khi hết tháng thứ $n$ thì  còn lại là $T_n=m\left(r+1\right)^n-\dfrac{a}{r} \left[\left(r+1\right)^n-1\right]$.\\
		Áp dụng công thức trên, ta có \[T_n=0\Leftrightarrow a=\dfrac{m\left(r+1\right)^n r}{\left(r+1\right)^n-1}=\dfrac{12\cdot  10^5\left(\dfrac{1{,}2}{100}+1\right)^{60}}{\left(\dfrac{1{,}2}{100}+1\right)^{60}-1}=2\,347\,614{,}532.\]
	 }
\end{ex}
